\makeatletter
\def\input@path{{./}{../}{../../}{../../../}{../../../../}{preamble/}{../preamble/}{../../preamble/}{../../../preamble/}{../../../../preamble/}}
\makeatother

\input{preamble_G10_T1_practice_activities_c1_c3.tex}
\input{graphs/highlights.tex}
% \GTenHighlightsfalse % Uncomment to disable instructional highlights.
\input{graphs/specific_highlights.tex}


\title{10th Grade Midterm Analysis of Decisions: C4--C5\\Solutions}
\author{}
\date{}

\begin{document}
	\maketitle
	
	\setcounter{tocdepth}{1}
	\setcounter{secnumdepth}{2}
	
	\phantomsection
	\label{toc}
	\tableofcontents
	
	\section*{Evaluated criteria}
	\begin{enumerate}[label=\textbf{C\arabic*}, start=4]
		\item Applies the Maximin and Maximax criteria to make optimistic or conservative decisions without probabilities, selecting the alternative with the highest possible payoff or the best worst-case, interpreting risks or limitations.
		\item Constructs opportunity loss tables and applies the Minimax regret criterion to select the alternative that minimises the maximum opport
	\end{enumerate}
	
	\begin{center}
		\textit{All payoffs are in thousands of dollars.}\\
		\textit{Regret \(=\) best payoff in the state \(-\) payoff.}
	\end{center}
	
	% ================================================================
	% ================================================================
	\section{Introduction}
	\label{sec:introduction}
\phantomsection\label{tone:legend}
\PracticeHighlightLegend{G10-T1-C4-C5}

	
	This section presents the general procedure for analysing a \GTenToneHL{payoff}{payoff} table under
	uncertainty. The \GTenToneHL{alternative}{alternatives} are denoted by \(\GTenToneHL{alternative}{A_i}\), the \GTenToneHL{state}{states} \GTenToneHL{state}{of} \GTenToneHL{state}{nature} by
	\(\GTenToneHL{state}{S_j}\), and the \GTenToneHL{payoff}{payoff} obtained by choosing \GTenToneHL{alternative}{alternative} \(\GTenToneHL{alternative}{A_i}\) when state
	\(\GTenToneHL{state}{S_j}\) occurs by
	\[
	\GTenToneHL{payoff}{P}(\GTenToneHL{alternative}{A_i},\GTenToneHL{state}{S_j}).
	\]
	
	The analysis uses three decision criteria:
	
	\begin{enumerate}
		\item \textbf{\GTenToneHL{best}{Maximax}:} choose the \GTenToneHL{alternative}{alternative} with the largest possible
		\GTenToneHL{payoff}{payoff}.
		\item \textbf{\GTenToneHL{worst}{Maximin}:} choose the \GTenToneHL{alternative}{alternative} with the largest worst-case
		\GTenToneHL{payoff}{payoff}.
		\item \textbf{\GTenToneHL{maxregret}{Minimax} \GTenToneHL{maxregret}{Regret}:} choose the \GTenToneHL{alternative}{alternative} whose maximum
		\GTenToneHL{regret}{opportunity} \GTenToneHL{regret}{loss} is the smallest.
	\end{enumerate}
	
	Because the \GTenToneHL{payoff}{payoffs} in this section are represented symbolically, the final
	choice depends on the numerical values assigned to the expressions
	\(\GTenToneHL{payoff}{P}(\GTenToneHL{alternative}{A_i},\GTenToneHL{state}{S_j})\).
	
	\SectionProblemsTableOfContents{
		\SectionProblemLink{sec:general-problem-2x2}
		{General \(2\times 2\) Problem} \\
		\SectionProblemLink{sec:general-problem-4x3}
		{General \(4\times 3\) Problem}
	}
	
	
	% ----------------------------------------------------------------
	\subsection{General \(2\times 2\) Problem}
	\label{sec:general-problem-2x2}
	
	A decision maker must choose between two \GTenToneHL{alternative}{alternatives}, \(\GTenToneHL{alternative}{A_1}\) and \(\GTenToneHL{alternative}{A_2}\).
	Two \GTenToneHL{state}{states} \GTenToneHL{state}{of} \GTenToneHL{state}{nature}, \(\GTenToneHL{state}{S_1}\) and \(\GTenToneHL{state}{S_2}\), may occur. The \GTenToneHL{payoff}{payoff} associated
	with each \GTenToneHL{alternative}{alternative}--state combination is represented by \(\GTenToneHL{payoff}{P}(\GTenToneHL{alternative}{A_i},\GTenToneHL{state}{S_j})\),
	where \(i\) identifies the \GTenToneHL{alternative}{alternative} and \(j\) identifies the state of
	nature.
	
	\begin{center}
		\textit{General \GTenToneHL{payoff}{payoff} table} \\
		\begin{tabular}{l c c}
			\toprule
			State of nature & Alternative \(\GTenToneHL{alternative}{A_1}\) & Alternative \(\GTenToneHL{alternative}{A_2}\) \\
			\midrule
			State \(\GTenToneHL{state}{S_1}\) & \(\GTenToneHL{payoff}{P}(\GTenToneHL{alternative}{A_1},\GTenToneHL{state}{S_1})\) & \(\GTenToneHL{payoff}{P}(\GTenToneHL{alternative}{A_2},\GTenToneHL{state}{S_1})\) \\
			State \(\GTenToneHL{state}{S_2}\) & \(\GTenToneHL{payoff}{P}(\GTenToneHL{alternative}{A_1},\GTenToneHL{state}{S_2})\) & \(\GTenToneHL{payoff}{P}(\GTenToneHL{alternative}{A_2},\GTenToneHL{state}{S_2})\) \\
			\bottomrule
		\end{tabular}
	\end{center}
	
	Use \GTenToneHL{best}{Maximax}, \GTenToneHL{worst}{Maximin}, and \GTenToneHL{maxregret}{Minimax} \GTenToneHL{maxregret}{Regret} to determine the decision rules for
	this general \GTenToneHL{payoff}{payoff} table.
	
	\textbf{Questions/Tasks.}
	\begin{enumerate}
		\CriterionTask
		{subsec:general-problem-2x2-c4}{4}
		{Applies the Maximin and Maximax criteria}

		\CriterionTask
		{subsec:general-problem-2x2-c5}{5}
		{Constructs opportunity loss tables and applies the Minimax regret criterion}
	\end{enumerate}
	
	
	\subsubsection*{C4 - Applies the Maximin and Maximax criteria}
	\phantomsection
	\label{subsec:general-problem-2x2-c4}
	
	\textbf{Step 1: Find the largest \GTenToneHL{payoff}{payoff} for each \GTenToneHL{alternative}{alternative}.}
	
	For \(\GTenToneHL{alternative}{A_1}\), compare its \GTenToneHL{payoff}{payoffs} under \(\GTenToneHL{state}{S_1}\) and \(\GTenToneHL{state}{S_2}\):
	\[
	\GTenToneHL{best}{M_1}=\max\{\GTenToneHL{payoff}{P}(\GTenToneHL{alternative}{A_1},\GTenToneHL{state}{S_1}),\GTenToneHL{payoff}{P}(\GTenToneHL{alternative}{A_1},\GTenToneHL{state}{S_2})\}.
	\]
	
	For \(\GTenToneHL{alternative}{A_2}\), compare its \GTenToneHL{payoff}{payoffs} under \(\GTenToneHL{state}{S_1}\) and \(\GTenToneHL{state}{S_2}\):
	\[
	\GTenToneHL{best}{M_2}=\max\{\GTenToneHL{payoff}{P}(\GTenToneHL{alternative}{A_2},\GTenToneHL{state}{S_1}),\GTenToneHL{payoff}{P}(\GTenToneHL{alternative}{A_2},\GTenToneHL{state}{S_2})\}.
	\]
	
	Therefore,
	\[
	\begin{aligned}
		\GTenToneHL{best}{M_1}&=\max\{\GTenToneHL{payoff}{P}(\GTenToneHL{alternative}{A_1},\GTenToneHL{state}{S_1}),\GTenToneHL{payoff}{P}(\GTenToneHL{alternative}{A_1},\GTenToneHL{state}{S_2})\},\\
		\GTenToneHL{best}{M_2}&=\max\{\GTenToneHL{payoff}{P}(\GTenToneHL{alternative}{A_2},\GTenToneHL{state}{S_1}),\GTenToneHL{payoff}{P}(\GTenToneHL{alternative}{A_2},\GTenToneHL{state}{S_2})\}.
	\end{aligned}
	\]
	
	\textbf{Step 2: Apply the \GTenToneHL{best}{Maximax} criterion.}
	
	Compare the largest \GTenToneHL{payoff}{payoffs} \(\GTenToneHL{best}{M_1}\) and \(\GTenToneHL{best}{M_2}\):
	\[
	M^*=\max\{\GTenToneHL{best}{M_1},\GTenToneHL{best}{M_2}\}.
	\]
	
	Equivalently,
	\[
	M^*
	=
	\max\left\{
	\max\{\GTenToneHL{payoff}{P}(\GTenToneHL{alternative}{A_1},\GTenToneHL{state}{S_1}),\GTenToneHL{payoff}{P}(\GTenToneHL{alternative}{A_1},\GTenToneHL{state}{S_2})\},
	\max\{\GTenToneHL{payoff}{P}(\GTenToneHL{alternative}{A_2},\GTenToneHL{state}{S_1}),\GTenToneHL{payoff}{P}(\GTenToneHL{alternative}{A_2},\GTenToneHL{state}{S_2})\}
	\right\}.
	\]
	
	The \GTenToneHL{best}{Maximax} decision is
	\[
	A_{\text{Maximax}}
	=
	\operatorname*{arg\,max}_{i\in\{1,2\}}
	\left\{
	\max_{j\in\{1,2\}}\GTenToneHL{payoff}{P}(\GTenToneHL{alternative}{A_i},\GTenToneHL{state}{S_j})
	\right\}.
	\]
	
	Thus, choose \(\GTenToneHL{alternative}{A_1}\) if \(\GTenToneHL{best}{M_1}>\GTenToneHL{best}{M_2}\), choose \(\GTenToneHL{alternative}{A_2}\) if \(\GTenToneHL{best}{M_2}>\GTenToneHL{best}{M_1}\), and
	report a tie if \(\GTenToneHL{best}{M_1}=\GTenToneHL{best}{M_2}\).
	
	\textbf{Step 3: Find the smallest \GTenToneHL{payoff}{payoff} for each \GTenToneHL{alternative}{alternative}.}
	
	For \(\GTenToneHL{alternative}{A_1}\),
	\[
	\GTenToneHL{worst}{m_1}=\min\{\GTenToneHL{payoff}{P}(\GTenToneHL{alternative}{A_1},\GTenToneHL{state}{S_1}),\GTenToneHL{payoff}{P}(\GTenToneHL{alternative}{A_1},\GTenToneHL{state}{S_2})\}.
	\]
	
	For \(\GTenToneHL{alternative}{A_2}\),
	\[
	\GTenToneHL{worst}{m_2}=\min\{\GTenToneHL{payoff}{P}(\GTenToneHL{alternative}{A_2},\GTenToneHL{state}{S_1}),\GTenToneHL{payoff}{P}(\GTenToneHL{alternative}{A_2},\GTenToneHL{state}{S_2})\}.
	\]
	
	Therefore,
	\[
	\begin{aligned}
		\GTenToneHL{worst}{m_1}&=\min\{\GTenToneHL{payoff}{P}(\GTenToneHL{alternative}{A_1},\GTenToneHL{state}{S_1}),\GTenToneHL{payoff}{P}(\GTenToneHL{alternative}{A_1},\GTenToneHL{state}{S_2})\},\\
		\GTenToneHL{worst}{m_2}&=\min\{\GTenToneHL{payoff}{P}(\GTenToneHL{alternative}{A_2},\GTenToneHL{state}{S_1}),\GTenToneHL{payoff}{P}(\GTenToneHL{alternative}{A_2},\GTenToneHL{state}{S_2})\}.
	\end{aligned}
	\]
	
	\textbf{Step 4: Apply the \GTenToneHL{worst}{Maximin} criterion.}
	
	Compare the smallest \GTenToneHL{payoff}{payoffs} \(\GTenToneHL{worst}{m_1}\) and \(\GTenToneHL{worst}{m_2}\):
	\[
	m^*=\max\{\GTenToneHL{worst}{m_1},\GTenToneHL{worst}{m_2}\}.
	\]
	
	Equivalently,
	\[
	m^*
	=
	\max\left\{
	\min\{\GTenToneHL{payoff}{P}(\GTenToneHL{alternative}{A_1},\GTenToneHL{state}{S_1}),\GTenToneHL{payoff}{P}(\GTenToneHL{alternative}{A_1},\GTenToneHL{state}{S_2})\},
	\min\{\GTenToneHL{payoff}{P}(\GTenToneHL{alternative}{A_2},\GTenToneHL{state}{S_1}),\GTenToneHL{payoff}{P}(\GTenToneHL{alternative}{A_2},\GTenToneHL{state}{S_2})\}
	\right\}.
	\]
	
	The \GTenToneHL{worst}{Maximin} decision is
	\[
	A_{\text{Maximin}}
	=
	\operatorname*{arg\,max}_{i\in\{1,2\}}
	\left\{
	\min_{j\in\{1,2\}}\GTenToneHL{payoff}{P}(\GTenToneHL{alternative}{A_i},\GTenToneHL{state}{S_j})
	\right\}.
	\]
	
	Thus, choose \(\GTenToneHL{alternative}{A_1}\) if \(\GTenToneHL{worst}{m_1}>\GTenToneHL{worst}{m_2}\), choose \(\GTenToneHL{alternative}{A_2}\) if \(\GTenToneHL{worst}{m_2}>\GTenToneHL{worst}{m_1}\), and
	report a tie if \(\GTenToneHL{worst}{m_1}=\GTenToneHL{worst}{m_2}\).
	
	
	\subsubsection*{C5 - Constructs opportunity loss tables and applies the Minimax regret criterion}
	\phantomsection
	\label{subsec:general-problem-2x2-c5}
	
	\textbf{Step 1: Find the best \GTenToneHL{payoff}{payoff} in each \GTenToneHL{state}{state} \GTenToneHL{state}{of} \GTenToneHL{state}{nature}.}
	
	Under \(\GTenToneHL{state}{S_1}\), compare the \GTenToneHL{payoff}{payoffs} of \(\GTenToneHL{alternative}{A_1}\) and \(\GTenToneHL{alternative}{A_2}\):
	\[
	\GTenToneHL{benchmark}{b_1}=\max\{\GTenToneHL{payoff}{P}(\GTenToneHL{alternative}{A_1},\GTenToneHL{state}{S_1}),\GTenToneHL{payoff}{P}(\GTenToneHL{alternative}{A_2},\GTenToneHL{state}{S_1})\}.
	\]
	
	Under \(\GTenToneHL{state}{S_2}\), compare the \GTenToneHL{payoff}{payoffs} of \(\GTenToneHL{alternative}{A_1}\) and \(\GTenToneHL{alternative}{A_2}\):
	\[
	\GTenToneHL{benchmark}{b_2}=\max\{\GTenToneHL{payoff}{P}(\GTenToneHL{alternative}{A_1},\GTenToneHL{state}{S_2}),\GTenToneHL{payoff}{P}(\GTenToneHL{alternative}{A_2},\GTenToneHL{state}{S_2})\}.
	\]
	
	Therefore,
	\[
	\begin{aligned}
		\GTenToneHL{benchmark}{b_1}&=\max\{\GTenToneHL{payoff}{P}(\GTenToneHL{alternative}{A_1},\GTenToneHL{state}{S_1}),\GTenToneHL{payoff}{P}(\GTenToneHL{alternative}{A_2},\GTenToneHL{state}{S_1})\},\\
		\GTenToneHL{benchmark}{b_2}&=\max\{\GTenToneHL{payoff}{P}(\GTenToneHL{alternative}{A_1},\GTenToneHL{state}{S_2}),\GTenToneHL{payoff}{P}(\GTenToneHL{alternative}{A_2},\GTenToneHL{state}{S_2})\}.
	\end{aligned}
	\]
	
	\textbf{Step 2: Calculate the \GTenToneHL{regret}{regret} for every \GTenToneHL{alternative}{alternative}--state
		combination.}
	
	The \GTenToneHL{regret}{regret} associated with \GTenToneHL{alternative}{alternative} \(\GTenToneHL{alternative}{A_i}\) under state \(\GTenToneHL{state}{S_j}\) is
	\[
	\GTenToneHL{regret}{r_{ij}}=\GTenToneHL{benchmark}{b_j}-\GTenToneHL{payoff}{P}(\GTenToneHL{alternative}{A_i},\GTenToneHL{state}{S_j}).
	\]
	
	Consequently,
	\[
	\begin{aligned}
		\GTenToneHL{regret}{r_{11}}&=\GTenToneHL{benchmark}{b_1}-\GTenToneHL{payoff}{P}(\GTenToneHL{alternative}{A_1},\GTenToneHL{state}{S_1}),&
		\GTenToneHL{regret}{r_{12}}&=\GTenToneHL{benchmark}{b_2}-\GTenToneHL{payoff}{P}(\GTenToneHL{alternative}{A_1},\GTenToneHL{state}{S_2}),\\
		\GTenToneHL{regret}{r_{21}}&=\GTenToneHL{benchmark}{b_1}-\GTenToneHL{payoff}{P}(\GTenToneHL{alternative}{A_2},\GTenToneHL{state}{S_1}),&
		\GTenToneHL{regret}{r_{22}}&=\GTenToneHL{benchmark}{b_2}-\GTenToneHL{payoff}{P}(\GTenToneHL{alternative}{A_2},\GTenToneHL{state}{S_2}).
	\end{aligned}
	\]
	
	The general \GTenToneHL{regret}{regret} table is therefore
	
	\begin{center}
		\begin{tabular}{l c c}
			\toprule
			Alternative & State \(\GTenToneHL{state}{S_1}\) & State \(\GTenToneHL{state}{S_2}\) \\
			\midrule
			\(\GTenToneHL{alternative}{A_1}\)
			& \(\GTenToneHL{benchmark}{b_1}-\GTenToneHL{payoff}{P}(\GTenToneHL{alternative}{A_1},\GTenToneHL{state}{S_1})\)
			& \(\GTenToneHL{benchmark}{b_2}-\GTenToneHL{payoff}{P}(\GTenToneHL{alternative}{A_1},\GTenToneHL{state}{S_2})\) \\
			\(\GTenToneHL{alternative}{A_2}\)
			& \(\GTenToneHL{benchmark}{b_1}-\GTenToneHL{payoff}{P}(\GTenToneHL{alternative}{A_2},\GTenToneHL{state}{S_1})\)
			& \(\GTenToneHL{benchmark}{b_2}-\GTenToneHL{payoff}{P}(\GTenToneHL{alternative}{A_2},\GTenToneHL{state}{S_2})\) \\
			\bottomrule
		\end{tabular}
	\end{center}
	
	\textbf{Step 3: Find the \GTenToneHL{maxregret}{maximum} \GTenToneHL{maxregret}{regret} for each \GTenToneHL{alternative}{alternative}.}
	
	For \(\GTenToneHL{alternative}{A_1}\),
	\[
	\GTenToneHL{maxregret}{R_1}=
	\max\{
	\GTenToneHL{benchmark}{b_1}-\GTenToneHL{payoff}{P}(\GTenToneHL{alternative}{A_1},\GTenToneHL{state}{S_1}),
	\GTenToneHL{benchmark}{b_2}-\GTenToneHL{payoff}{P}(\GTenToneHL{alternative}{A_1},\GTenToneHL{state}{S_2})
	\}.
	\]
	
	For \(\GTenToneHL{alternative}{A_2}\),
	\[
	\GTenToneHL{maxregret}{R_2}=
	\max\{
	\GTenToneHL{benchmark}{b_1}-\GTenToneHL{payoff}{P}(\GTenToneHL{alternative}{A_2},\GTenToneHL{state}{S_1}),
	\GTenToneHL{benchmark}{b_2}-\GTenToneHL{payoff}{P}(\GTenToneHL{alternative}{A_2},\GTenToneHL{state}{S_2})
	\}.
	\]
	
	Thus,
	\[
	\begin{aligned}
		\GTenToneHL{maxregret}{R_1}&=\max\{\GTenToneHL{regret}{r_{11}},\GTenToneHL{regret}{r_{12}}\},\\
		\GTenToneHL{maxregret}{R_2}&=\max\{\GTenToneHL{regret}{r_{21}},\GTenToneHL{regret}{r_{22}}\}.
	\end{aligned}
	\]
	
	\textbf{Step 4: Apply the \GTenToneHL{maxregret}{Minimax} \GTenToneHL{maxregret}{Regret} criterion.}
	
	Compare the maximum regrets \(\GTenToneHL{maxregret}{R_1}\) and \(\GTenToneHL{maxregret}{R_2}\):
	\[
	R^*=\min\{\GTenToneHL{maxregret}{R_1},\GTenToneHL{maxregret}{R_2}\}.
	\]
	
	The \GTenToneHL{maxregret}{Minimax} \GTenToneHL{maxregret}{Regret} decision is
	\[
	A_{\text{Minimax Regret}}
	=
	\operatorname*{arg\,min}_{i\in\{1,2\}}
	\left\{
	\max_{j\in\{1,2\}}
	\left[\GTenToneHL{benchmark}{b_j}-\GTenToneHL{payoff}{P}(\GTenToneHL{alternative}{A_i},\GTenToneHL{state}{S_j})\right]
	\right\}.
	\]
	
	Thus, choose \(\GTenToneHL{alternative}{A_1}\) if \(\GTenToneHL{maxregret}{R_1}<\GTenToneHL{maxregret}{R_2}\), choose \(\GTenToneHL{alternative}{A_2}\) if \(\GTenToneHL{maxregret}{R_2}<\GTenToneHL{maxregret}{R_1}\), and
	report a tie if \(\GTenToneHL{maxregret}{R_1}=\GTenToneHL{maxregret}{R_2}\).
	
	
	\subsubsection*{General decisions summary}
	
	\begin{center}
		\begin{tabular}{l c}
			\toprule
			Criterion & General decision rule \\
			\midrule
			\GTenToneHL{best}{Maximax}
			& \(\displaystyle
			\operatorname*{arg\,max}_{i\in\{1,2\}}
			\max_{j\in\{1,2\}}\GTenToneHL{payoff}{P}(\GTenToneHL{alternative}{A_i},\GTenToneHL{state}{S_j})\) \\[0.25cm]
			Maximin
			& \(\displaystyle
			\operatorname*{arg\,max}_{i\in\{1,2\}}
			\min_{j\in\{1,2\}}\GTenToneHL{payoff}{P}(\GTenToneHL{alternative}{A_i},\GTenToneHL{state}{S_j})\) \\[0.25cm]
			Minimax Regret
			& \(\displaystyle
			\operatorname*{arg\,min}_{i\in\{1,2\}}
			\max_{j\in\{1,2\}}
			\left[\GTenToneHL{benchmark}{b_j}-\GTenToneHL{payoff}{P}(\GTenToneHL{alternative}{A_i},\GTenToneHL{state}{S_j})\right]\) \\
			\bottomrule
		\end{tabular}
	\end{center}
	
	\secInicio[sec:general-problem-2x2][sec:introduction]
	
	
	% ----------------------------------------------------------------
	\subsection{General \(4\times 3\) Problem}
	\label{sec:general-problem-4x3}
	
	A decision maker must choose among three alternatives, \(\GTenToneHL{alternative}{A_1}\), \(\GTenToneHL{alternative}{A_2}\), and
	\(\GTenToneHL{alternative}{A_3}\). Four states of nature, \(\GTenToneHL{state}{S_1}\), \(\GTenToneHL{state}{S_2}\), \(\GTenToneHL{state}{S_3}\), and \(\GTenToneHL{state}{S_4}\), may
	occur. The payoff associated with each alternative--state combination is
	represented by \(\GTenToneHL{payoff}{P}(\GTenToneHL{alternative}{A_i},\GTenToneHL{state}{S_j})\).
	
	\begin{center}
		\textit{General payoff table} \\
		\begin{tabular}{l c c c}
			\toprule
			State of nature
			& Alternative \(\GTenToneHL{alternative}{A_1}\)
			& Alternative \(\GTenToneHL{alternative}{A_2}\)
			& Alternative \(\GTenToneHL{alternative}{A_3}\) \\
			\midrule
			State \(\GTenToneHL{state}{S_1}\)
			& \(\GTenToneHL{payoff}{P}(\GTenToneHL{alternative}{A_1},\GTenToneHL{state}{S_1})\)
			& \(\GTenToneHL{payoff}{P}(\GTenToneHL{alternative}{A_2},\GTenToneHL{state}{S_1})\)
			& \(\GTenToneHL{payoff}{P}(\GTenToneHL{alternative}{A_3},\GTenToneHL{state}{S_1})\) \\
			State \(\GTenToneHL{state}{S_2}\)
			& \(\GTenToneHL{payoff}{P}(\GTenToneHL{alternative}{A_1},\GTenToneHL{state}{S_2})\)
			& \(\GTenToneHL{payoff}{P}(\GTenToneHL{alternative}{A_2},\GTenToneHL{state}{S_2})\)
			& \(\GTenToneHL{payoff}{P}(\GTenToneHL{alternative}{A_3},\GTenToneHL{state}{S_2})\) \\
			State \(\GTenToneHL{state}{S_3}\)
			& \(\GTenToneHL{payoff}{P}(\GTenToneHL{alternative}{A_1},\GTenToneHL{state}{S_3})\)
			& \(\GTenToneHL{payoff}{P}(\GTenToneHL{alternative}{A_2},\GTenToneHL{state}{S_3})\)
			& \(\GTenToneHL{payoff}{P}(\GTenToneHL{alternative}{A_3},\GTenToneHL{state}{S_3})\) \\
			State \(\GTenToneHL{state}{S_4}\)
			& \(\GTenToneHL{payoff}{P}(\GTenToneHL{alternative}{A_1},\GTenToneHL{state}{S_4})\)
			& \(\GTenToneHL{payoff}{P}(\GTenToneHL{alternative}{A_2},\GTenToneHL{state}{S_4})\)
			& \(\GTenToneHL{payoff}{P}(\GTenToneHL{alternative}{A_3},\GTenToneHL{state}{S_4})\) \\
			\bottomrule
		\end{tabular}
	\end{center}
	
	Use Maximax, Maximin, and Minimax Regret to determine the decision rules for
	this general payoff table.
	
	\textbf{Questions/Tasks.}
	\begin{enumerate}
		\CriterionTask
		{subsec:general-problem-4x3-c4}{4}
		{Applies the Maximin and Maximax criteria}

		\CriterionTask
		{subsec:general-problem-4x3-c5}{5}
		{Constructs opportunity loss tables and applies the Minimax regret criterion}
	\end{enumerate}
	
	
	\subsubsection*{C4 - Applies the Maximin and Maximax criteria}
	\phantomsection
	\label{subsec:general-problem-4x3-c4}
	
	\textbf{Step 1: Find the largest payoff for each alternative.}
	
	For the three alternatives,
	\[
	\begin{aligned}
		\GTenToneHL{best}{M_1}&=\max\{
		\GTenToneHL{payoff}{P}(\GTenToneHL{alternative}{A_1},\GTenToneHL{state}{S_1}),\GTenToneHL{payoff}{P}(\GTenToneHL{alternative}{A_1},\GTenToneHL{state}{S_2}),\GTenToneHL{payoff}{P}(\GTenToneHL{alternative}{A_1},\GTenToneHL{state}{S_3}),\GTenToneHL{payoff}{P}(\GTenToneHL{alternative}{A_1},\GTenToneHL{state}{S_4})
		\},\\
		\GTenToneHL{best}{M_2}&=\max\{
		\GTenToneHL{payoff}{P}(\GTenToneHL{alternative}{A_2},\GTenToneHL{state}{S_1}),\GTenToneHL{payoff}{P}(\GTenToneHL{alternative}{A_2},\GTenToneHL{state}{S_2}),\GTenToneHL{payoff}{P}(\GTenToneHL{alternative}{A_2},\GTenToneHL{state}{S_3}),\GTenToneHL{payoff}{P}(\GTenToneHL{alternative}{A_2},\GTenToneHL{state}{S_4})
		\},\\
		\GTenToneHL{best}{M_3}&=\max\{
		\GTenToneHL{payoff}{P}(\GTenToneHL{alternative}{A_3},\GTenToneHL{state}{S_1}),\GTenToneHL{payoff}{P}(\GTenToneHL{alternative}{A_3},\GTenToneHL{state}{S_2}),\GTenToneHL{payoff}{P}(\GTenToneHL{alternative}{A_3},\GTenToneHL{state}{S_3}),\GTenToneHL{payoff}{P}(\GTenToneHL{alternative}{A_3},\GTenToneHL{state}{S_4})
		\}.
	\end{aligned}
	\]
	
	\textbf{Step 2: Apply the \GTenToneHL{best}{Maximax} criterion.}
	
	Compare the three largest \GTenToneHL{payoff}{payoffs}:
	\[
	M^*=\max\{\GTenToneHL{best}{M_1},\GTenToneHL{best}{M_2},\GTenToneHL{best}{M_3}\}.
	\]
	
	Equivalently,
	\[
	\begin{aligned}
		M^*=\max\{&
		\max\{
		\GTenToneHL{payoff}{P}(\GTenToneHL{alternative}{A_1},\GTenToneHL{state}{S_1}),\GTenToneHL{payoff}{P}(\GTenToneHL{alternative}{A_1},\GTenToneHL{state}{S_2}),\GTenToneHL{payoff}{P}(\GTenToneHL{alternative}{A_1},\GTenToneHL{state}{S_3}),\GTenToneHL{payoff}{P}(\GTenToneHL{alternative}{A_1},\GTenToneHL{state}{S_4})
		\},\\
		&
		\max\{
		\GTenToneHL{payoff}{P}(\GTenToneHL{alternative}{A_2},\GTenToneHL{state}{S_1}),\GTenToneHL{payoff}{P}(\GTenToneHL{alternative}{A_2},\GTenToneHL{state}{S_2}),\GTenToneHL{payoff}{P}(\GTenToneHL{alternative}{A_2},\GTenToneHL{state}{S_3}),\GTenToneHL{payoff}{P}(\GTenToneHL{alternative}{A_2},\GTenToneHL{state}{S_4})
		\},\\
		&
		\max\{
		\GTenToneHL{payoff}{P}(\GTenToneHL{alternative}{A_3},\GTenToneHL{state}{S_1}),\GTenToneHL{payoff}{P}(\GTenToneHL{alternative}{A_3},\GTenToneHL{state}{S_2}),\GTenToneHL{payoff}{P}(\GTenToneHL{alternative}{A_3},\GTenToneHL{state}{S_3}),\GTenToneHL{payoff}{P}(\GTenToneHL{alternative}{A_3},\GTenToneHL{state}{S_4})
		\}
		\}.
	\end{aligned}
	\]
	
	The \GTenToneHL{best}{Maximax} decision is
	\[
	A_{\text{Maximax}}
	=
	\operatorname*{arg\,max}_{i\in\{1,2,3\}}
	\left\{
	\max_{j\in\{1,2,3,4\}}\GTenToneHL{payoff}{P}(\GTenToneHL{alternative}{A_i},\GTenToneHL{state}{S_j})
	\right\}.
	\]
	
	Choose the \GTenToneHL{alternative}{alternative} whose value \(\GTenToneHL{best}{M_i}\) is the largest. If two or more
	\GTenToneHL{alternative}{alternatives} have the same largest value, the \GTenToneHL{best}{Maximax} criterion produces a
	tie.
	
	\textbf{Step 3: Find the smallest \GTenToneHL{payoff}{payoff} for each \GTenToneHL{alternative}{alternative}.}
	
	For the three \GTenToneHL{alternative}{alternatives},
	\[
	\begin{aligned}
		\GTenToneHL{worst}{m_1}&=\min\{
		\GTenToneHL{payoff}{P}(\GTenToneHL{alternative}{A_1},\GTenToneHL{state}{S_1}),\GTenToneHL{payoff}{P}(\GTenToneHL{alternative}{A_1},\GTenToneHL{state}{S_2}),\GTenToneHL{payoff}{P}(\GTenToneHL{alternative}{A_1},\GTenToneHL{state}{S_3}),\GTenToneHL{payoff}{P}(\GTenToneHL{alternative}{A_1},\GTenToneHL{state}{S_4})
		\},\\
		\GTenToneHL{worst}{m_2}&=\min\{
		\GTenToneHL{payoff}{P}(\GTenToneHL{alternative}{A_2},\GTenToneHL{state}{S_1}),\GTenToneHL{payoff}{P}(\GTenToneHL{alternative}{A_2},\GTenToneHL{state}{S_2}),\GTenToneHL{payoff}{P}(\GTenToneHL{alternative}{A_2},\GTenToneHL{state}{S_3}),\GTenToneHL{payoff}{P}(\GTenToneHL{alternative}{A_2},\GTenToneHL{state}{S_4})
		\},\\
		\GTenToneHL{worst}{m_3}&=\min\{
		\GTenToneHL{payoff}{P}(\GTenToneHL{alternative}{A_3},\GTenToneHL{state}{S_1}),\GTenToneHL{payoff}{P}(\GTenToneHL{alternative}{A_3},\GTenToneHL{state}{S_2}),\GTenToneHL{payoff}{P}(\GTenToneHL{alternative}{A_3},\GTenToneHL{state}{S_3}),\GTenToneHL{payoff}{P}(\GTenToneHL{alternative}{A_3},\GTenToneHL{state}{S_4})
		\}.
	\end{aligned}
	\]
	
	\textbf{Step 4: Apply the \GTenToneHL{worst}{Maximin} criterion.}
	
	Compare the three smallest \GTenToneHL{payoff}{payoffs}:
	\[
	m^*=\max\{\GTenToneHL{worst}{m_1},\GTenToneHL{worst}{m_2},\GTenToneHL{worst}{m_3}\}.
	\]
	
	Equivalently,
	\[
	\begin{aligned}
		m^*=\max\{&
		\min\{
		\GTenToneHL{payoff}{P}(\GTenToneHL{alternative}{A_1},\GTenToneHL{state}{S_1}),\GTenToneHL{payoff}{P}(\GTenToneHL{alternative}{A_1},\GTenToneHL{state}{S_2}),\GTenToneHL{payoff}{P}(\GTenToneHL{alternative}{A_1},\GTenToneHL{state}{S_3}),\GTenToneHL{payoff}{P}(\GTenToneHL{alternative}{A_1},\GTenToneHL{state}{S_4})
		\},\\
		&
		\min\{
		\GTenToneHL{payoff}{P}(\GTenToneHL{alternative}{A_2},\GTenToneHL{state}{S_1}),\GTenToneHL{payoff}{P}(\GTenToneHL{alternative}{A_2},\GTenToneHL{state}{S_2}),\GTenToneHL{payoff}{P}(\GTenToneHL{alternative}{A_2},\GTenToneHL{state}{S_3}),\GTenToneHL{payoff}{P}(\GTenToneHL{alternative}{A_2},\GTenToneHL{state}{S_4})
		\},\\
		&
		\min\{
		\GTenToneHL{payoff}{P}(\GTenToneHL{alternative}{A_3},\GTenToneHL{state}{S_1}),\GTenToneHL{payoff}{P}(\GTenToneHL{alternative}{A_3},\GTenToneHL{state}{S_2}),\GTenToneHL{payoff}{P}(\GTenToneHL{alternative}{A_3},\GTenToneHL{state}{S_3}),\GTenToneHL{payoff}{P}(\GTenToneHL{alternative}{A_3},\GTenToneHL{state}{S_4})
		\}
		\}.
	\end{aligned}
	\]
	
	The \GTenToneHL{worst}{Maximin} decision is
	\[
	A_{\text{Maximin}}
	=
	\operatorname*{arg\,max}_{i\in\{1,2,3\}}
	\left\{
	\min_{j\in\{1,2,3,4\}}\GTenToneHL{payoff}{P}(\GTenToneHL{alternative}{A_i},\GTenToneHL{state}{S_j})
	\right\}.
	\]
	
	Choose the \GTenToneHL{alternative}{alternative} whose value \(\GTenToneHL{worst}{m_i}\) is the largest. If two or more
	\GTenToneHL{alternative}{alternatives} have the same largest value, the \GTenToneHL{worst}{Maximin} criterion produces a
	tie.
	
	
	\subsubsection*{C5 - Constructs opportunity loss tables and applies the Minimax regret criterion}
	\phantomsection
	\label{subsec:general-problem-4x3-c5}
	
	\textbf{Step 1: Find the best \GTenToneHL{payoff}{payoff} in each \GTenToneHL{state}{state} \GTenToneHL{state}{of} \GTenToneHL{state}{nature}.}
	
	For the four states,
	\[
	\begin{aligned}
		\GTenToneHL{benchmark}{b_1}&=\max\{\GTenToneHL{payoff}{P}(\GTenToneHL{alternative}{A_1},\GTenToneHL{state}{S_1}),\GTenToneHL{payoff}{P}(\GTenToneHL{alternative}{A_2},\GTenToneHL{state}{S_1}),\GTenToneHL{payoff}{P}(\GTenToneHL{alternative}{A_3},\GTenToneHL{state}{S_1})\},\\
		\GTenToneHL{benchmark}{b_2}&=\max\{\GTenToneHL{payoff}{P}(\GTenToneHL{alternative}{A_1},\GTenToneHL{state}{S_2}),\GTenToneHL{payoff}{P}(\GTenToneHL{alternative}{A_2},\GTenToneHL{state}{S_2}),\GTenToneHL{payoff}{P}(\GTenToneHL{alternative}{A_3},\GTenToneHL{state}{S_2})\},\\
		\GTenToneHL{benchmark}{b_3}&=\max\{\GTenToneHL{payoff}{P}(\GTenToneHL{alternative}{A_1},\GTenToneHL{state}{S_3}),\GTenToneHL{payoff}{P}(\GTenToneHL{alternative}{A_2},\GTenToneHL{state}{S_3}),\GTenToneHL{payoff}{P}(\GTenToneHL{alternative}{A_3},\GTenToneHL{state}{S_3})\},\\
		\GTenToneHL{benchmark}{b_4}&=\max\{\GTenToneHL{payoff}{P}(\GTenToneHL{alternative}{A_1},\GTenToneHL{state}{S_4}),\GTenToneHL{payoff}{P}(\GTenToneHL{alternative}{A_2},\GTenToneHL{state}{S_4}),\GTenToneHL{payoff}{P}(\GTenToneHL{alternative}{A_3},\GTenToneHL{state}{S_4})\}.
	\end{aligned}
	\]
	
	\textbf{Step 2: Calculate the \GTenToneHL{regret}{regret} for every \GTenToneHL{alternative}{alternative}--state
		combination.}
	
	For every \GTenToneHL{alternative}{alternative} \(\GTenToneHL{alternative}{A_i}\) and state \(\GTenToneHL{state}{S_j}\), the \GTenToneHL{regret}{regret} is
	\[
	\GTenToneHL{regret}{r_{ij}}=\GTenToneHL{benchmark}{b_j}-\GTenToneHL{payoff}{P}(\GTenToneHL{alternative}{A_i},\GTenToneHL{state}{S_j}).
	\]
	
	The general \GTenToneHL{regret}{regret} table is therefore
	
	\begin{center}
		\begin{tabular}{l c c c c}
			\toprule
			Alternative
			& State \(\GTenToneHL{state}{S_1}\)
			& State \(\GTenToneHL{state}{S_2}\)
			& State \(\GTenToneHL{state}{S_3}\)
			& State \(\GTenToneHL{state}{S_4}\) \\
			\midrule
			\(\GTenToneHL{alternative}{A_1}\)
			& \(\GTenToneHL{benchmark}{b_1}-\GTenToneHL{payoff}{P}(\GTenToneHL{alternative}{A_1},\GTenToneHL{state}{S_1})\)
			& \(\GTenToneHL{benchmark}{b_2}-\GTenToneHL{payoff}{P}(\GTenToneHL{alternative}{A_1},\GTenToneHL{state}{S_2})\)
			& \(\GTenToneHL{benchmark}{b_3}-\GTenToneHL{payoff}{P}(\GTenToneHL{alternative}{A_1},\GTenToneHL{state}{S_3})\)
			& \(\GTenToneHL{benchmark}{b_4}-\GTenToneHL{payoff}{P}(\GTenToneHL{alternative}{A_1},\GTenToneHL{state}{S_4})\) \\
			\(\GTenToneHL{alternative}{A_2}\)
			& \(\GTenToneHL{benchmark}{b_1}-\GTenToneHL{payoff}{P}(\GTenToneHL{alternative}{A_2},\GTenToneHL{state}{S_1})\)
			& \(\GTenToneHL{benchmark}{b_2}-\GTenToneHL{payoff}{P}(\GTenToneHL{alternative}{A_2},\GTenToneHL{state}{S_2})\)
			& \(\GTenToneHL{benchmark}{b_3}-\GTenToneHL{payoff}{P}(\GTenToneHL{alternative}{A_2},\GTenToneHL{state}{S_3})\)
			& \(\GTenToneHL{benchmark}{b_4}-\GTenToneHL{payoff}{P}(\GTenToneHL{alternative}{A_2},\GTenToneHL{state}{S_4})\) \\
			\(\GTenToneHL{alternative}{A_3}\)
			& \(\GTenToneHL{benchmark}{b_1}-\GTenToneHL{payoff}{P}(\GTenToneHL{alternative}{A_3},\GTenToneHL{state}{S_1})\)
			& \(\GTenToneHL{benchmark}{b_2}-\GTenToneHL{payoff}{P}(\GTenToneHL{alternative}{A_3},\GTenToneHL{state}{S_2})\)
			& \(\GTenToneHL{benchmark}{b_3}-\GTenToneHL{payoff}{P}(\GTenToneHL{alternative}{A_3},\GTenToneHL{state}{S_3})\)
			& \(\GTenToneHL{benchmark}{b_4}-\GTenToneHL{payoff}{P}(\GTenToneHL{alternative}{A_3},\GTenToneHL{state}{S_4})\) \\
			\bottomrule
		\end{tabular}
	\end{center}
	
	\textbf{Step 3: Find the \GTenToneHL{maxregret}{maximum} \GTenToneHL{maxregret}{regret} for each \GTenToneHL{alternative}{alternative}.}
	
	For \(\GTenToneHL{alternative}{A_1}\),
	\[
	\GTenToneHL{maxregret}{R_1}
	=
	\max\{
	\GTenToneHL{benchmark}{b_1}-\GTenToneHL{payoff}{P}(\GTenToneHL{alternative}{A_1},\GTenToneHL{state}{S_1}),
	\GTenToneHL{benchmark}{b_2}-\GTenToneHL{payoff}{P}(\GTenToneHL{alternative}{A_1},\GTenToneHL{state}{S_2}),
	\GTenToneHL{benchmark}{b_3}-\GTenToneHL{payoff}{P}(\GTenToneHL{alternative}{A_1},\GTenToneHL{state}{S_3}),
	\GTenToneHL{benchmark}{b_4}-\GTenToneHL{payoff}{P}(\GTenToneHL{alternative}{A_1},\GTenToneHL{state}{S_4})
	\}.
	\]
	
	For \(\GTenToneHL{alternative}{A_2}\),
	\[
	\GTenToneHL{maxregret}{R_2}
	=
	\max\{
	\GTenToneHL{benchmark}{b_1}-\GTenToneHL{payoff}{P}(\GTenToneHL{alternative}{A_2},\GTenToneHL{state}{S_1}),
	\GTenToneHL{benchmark}{b_2}-\GTenToneHL{payoff}{P}(\GTenToneHL{alternative}{A_2},\GTenToneHL{state}{S_2}),
	\GTenToneHL{benchmark}{b_3}-\GTenToneHL{payoff}{P}(\GTenToneHL{alternative}{A_2},\GTenToneHL{state}{S_3}),
	\GTenToneHL{benchmark}{b_4}-\GTenToneHL{payoff}{P}(\GTenToneHL{alternative}{A_2},\GTenToneHL{state}{S_4})
	\}.
	\]
	
	For \(\GTenToneHL{alternative}{A_3}\),
	\[
	\GTenToneHL{maxregret}{R_3}
	=
	\max\{
	\GTenToneHL{benchmark}{b_1}-\GTenToneHL{payoff}{P}(\GTenToneHL{alternative}{A_3},\GTenToneHL{state}{S_1}),
	\GTenToneHL{benchmark}{b_2}-\GTenToneHL{payoff}{P}(\GTenToneHL{alternative}{A_3},\GTenToneHL{state}{S_2}),
	\GTenToneHL{benchmark}{b_3}-\GTenToneHL{payoff}{P}(\GTenToneHL{alternative}{A_3},\GTenToneHL{state}{S_3}),
	\GTenToneHL{benchmark}{b_4}-\GTenToneHL{payoff}{P}(\GTenToneHL{alternative}{A_3},\GTenToneHL{state}{S_4})
	\}.
	\]
	
	Equivalently,
	\[
	\begin{aligned}
		\GTenToneHL{maxregret}{R_1}&=\max\{\GTenToneHL{regret}{r_{11}},\GTenToneHL{regret}{r_{12}},\GTenToneHL{regret}{r_{13}},\GTenToneHL{regret}{r_{14}}\},\\
		\GTenToneHL{maxregret}{R_2}&=\max\{\GTenToneHL{regret}{r_{21}},\GTenToneHL{regret}{r_{22}},\GTenToneHL{regret}{r_{23}},\GTenToneHL{regret}{r_{24}}\},\\
		\GTenToneHL{maxregret}{R_3}&=\max\{\GTenToneHL{regret}{r_{31}},\GTenToneHL{regret}{r_{32}},\GTenToneHL{regret}{r_{33}},\GTenToneHL{regret}{r_{34}}\}.
	\end{aligned}
	\]
	
	\textbf{Step 4: Apply the \GTenToneHL{maxregret}{Minimax} \GTenToneHL{maxregret}{Regret} criterion.}
	
	Compare the maximum regrets:
	\[
	R^*=\min\{\GTenToneHL{maxregret}{R_1},\GTenToneHL{maxregret}{R_2},\GTenToneHL{maxregret}{R_3}\}.
	\]
	
	The \GTenToneHL{maxregret}{Minimax} \GTenToneHL{maxregret}{Regret} decision is
	\[
	A_{\text{Minimax Regret}}
	=
	\operatorname*{arg\,min}_{i\in\{1,2,3\}}
	\left\{
	\max_{j\in\{1,2,3,4\}}
	\left[\GTenToneHL{benchmark}{b_j}-\GTenToneHL{payoff}{P}(\GTenToneHL{alternative}{A_i},\GTenToneHL{state}{S_j})\right]
	\right\}.
	\]
	
	Choose the \GTenToneHL{alternative}{alternative} whose value \(\GTenToneHL{maxregret}{R_i}\) is the smallest. If two or more
	\GTenToneHL{alternative}{alternatives} have the same smallest value, the \GTenToneHL{maxregret}{Minimax} \GTenToneHL{maxregret}{Regret} criterion
	produces a tie.
	
	
	\subsubsection*{General decisions summary}
	
	\begin{center}
		\begin{tabular}{l c}
			\toprule
			Criterion & General decision rule \\
			\midrule
			\GTenToneHL{best}{Maximax}
			& \(\displaystyle
			\operatorname*{arg\,max}_{i\in\{1,2,3\}}
			\max_{j\in\{1,2,3,4\}}\GTenToneHL{payoff}{P}(\GTenToneHL{alternative}{A_i},\GTenToneHL{state}{S_j})\) \\[0.25cm]
			Maximin
			& \(\displaystyle
			\operatorname*{arg\,max}_{i\in\{1,2,3\}}
			\min_{j\in\{1,2,3,4\}}\GTenToneHL{payoff}{P}(\GTenToneHL{alternative}{A_i},\GTenToneHL{state}{S_j})\) \\[0.25cm]
			Minimax Regret
			& \(\displaystyle
			\operatorname*{arg\,min}_{i\in\{1,2,3\}}
			\max_{j\in\{1,2,3,4\}}
			\left[\GTenToneHL{benchmark}{b_j}-\GTenToneHL{payoff}{P}(\GTenToneHL{alternative}{A_i},\GTenToneHL{state}{S_j})\right]\) \\
			\bottomrule
		\end{tabular}
	\end{center}
	
	\secInicio[sec:general-problem-4x3][sec:introduction]
	
	
	% ================================================================
	
\phantomsection\label{tone:reference}
\subsection*{Notation reference}
\GTenToneReference{C4-C5}
\section{Problems \(2\times 2\)}
	\label{sec:problems-2x2}
	
	\SectionProblemsTableOfContents{
		\SectionProblemLink{sec:food-truck-location}{Food Truck Location (ABB)} \\
		\SectionProblemLink{sec:promotional-campaign}{Promotional Campaign (ABA)} \\
		\SectionProblemLink{sec:equipment-lease}{Equipment Lease (AAA)}
	}
	
	% ----------------------------------------------------------------
	\subsection{Food Truck Location}
	\label{sec:food-truck-location}
	
	A food truck operator must choose between Location A and Location B. Demand can
	be strong or weak. The final payoff table, in thousands of dollars, is given
	below. Use Maximax, Maximin, and Minimax Regret to select a location.
	
	\begin{center}
		\textit{Payoff table} \\
		\begin{tabular}{l c c}
			\toprule
			State of nature & Location A & Location B \\
			\midrule
			Strong demand & \GTenToneHL{payoff}{100} & \GTenToneHL{payoff}{60} \\
			Weak demand   & \GTenToneHL{payoff}{0}   & \GTenToneHL{payoff}{50} \\
			\bottomrule
		\end{tabular}
	\end{center}
	
	\textbf{Questions/Tasks.}
	\begin{enumerate}
		\CriterionTask
		{subsec:food-truck-location-c4}{4}
		{Applies the Maximin and Maximax criteria}

		\CriterionTask
		{subsec:food-truck-location-c5}{5}
		{Constructs opportunity loss tables and applies the Minimax regret criterion}
	\end{enumerate}
	
	\subsubsection*{C4 - Applies the Maximin and Maximax criteria}
	\phantomsection
	\label{subsec:food-truck-location-c4}
	
	\[
	\begin{aligned}
		\max(\text{Location A}) &= \max\{100,0\}=100,\\
		\max(\text{Location B}) &= \max\{60,50\}=60.
	\end{aligned}
	\]
	\[
	\max\{\max(\text{Location A}),\max(\text{Location B})\}
	=\max\{100,60\}
	=100.
	\]
	The \GTenToneHL{best}{Maximax} choice is Location A because it has the highest possible \GTenToneHL{payoff}{payoff}.
	
	\[
	\begin{aligned}
		\min(\text{Location A}) &= \min\{100,0\}=0,\\
		\min(\text{Location B}) &= \min\{60,50\}=50.
	\end{aligned}
	\]
	\[
	\max\{\min(\text{Location A}),\min(\text{Location B})\}
	=\max\{0,50\}
	=50.
	\]
	The \GTenToneHL{worst}{Maximin} choice is Location B because it has the largest worst-case \GTenToneHL{payoff}{payoff}.
	
	\subsubsection*{C5 - Constructs opportunity loss tables and applies the Minimax regret criterion}
	\phantomsection
	\label{subsec:food-truck-location-c5}
	
	Best \GTenToneHL{payoff}{payoff} in each state:
	\[
	\begin{aligned}
		\text{Strong demand: }&\max\{100,60\}=100,\\
		\text{Weak demand: }&\max\{0,50\}=50.
	\end{aligned}
	\]
	
	Regret table (best \GTenToneHL{payoff}{payoff} \( - \) \GTenToneHL{payoff}{payoff}):
	\begin{center}
		\begin{tabular}{l c c c}
			\toprule
			Alternative & Strong demand & Weak demand & \GTenToneHL{maxregret}{Maximum} \GTenToneHL{maxregret}{regret} \\
			\midrule
			Location A & \(100-100=0\) & \(50-0=50\)  & \GTenToneHL{maxregret}{50} \\
			Location B & \(100-60=40\) & \(50-50=0\) & \GTenToneHL{maxregret}{40} \\
			\bottomrule
		\end{tabular}
	\end{center}
	
	The \GTenToneHL{maxregret}{Minimax} \GTenToneHL{maxregret}{Regret} choice is Location B because it has the smallest maximum
	\GTenToneHL{regret}{regret}, \(40\).
	
	\subsubsection*{Decisions summary}
	
	\begin{center}
		\begin{tabular}{l c}
			\toprule
			Criterion & Final decision \\
			\midrule
			\GTenToneHL{best}{Maximax}        & Location A \\
			\GTenToneHL{worst}{Maximin}        & Location B \\
			\GTenToneHL{maxregret}{Minimax} \GTenToneHL{maxregret}{Regret} & Location B \\
			\bottomrule
		\end{tabular}
	\end{center}
	
	\secInicio[sec:food-truck-location][sec:problems-2x2]
	
	
	% ----------------------------------------------------------------
	\subsection{Promotional Campaign}
	\label{sec:promotional-campaign}
	
	A business must choose between Campaign A and Campaign B. Customer response can
	be strong or weak. The final \GTenToneHL{payoff}{payoff} table, in thousands of dollars, is given
	below. Use \GTenToneHL{best}{Maximax}, \GTenToneHL{worst}{Maximin}, and \GTenToneHL{maxregret}{Minimax} \GTenToneHL{maxregret}{Regret} to select a campaign.
	
	\begin{center}
		\textit{Payoff table} \\
		\begin{tabular}{l c c}
			\toprule
			State of nature & Campaign A & Campaign B \\
			\midrule
			Strong response & \GTenToneHL{payoff}{100} & \GTenToneHL{payoff}{80} \\
			Weak response   & \GTenToneHL{payoff}{20}  & \GTenToneHL{payoff}{30} \\
			\bottomrule
		\end{tabular}
	\end{center}
	
	\textbf{Questions/Tasks.}
	\begin{enumerate}
		\CriterionTask
		{subsec:promotional-campaign-c4}{4}
		{Applies the Maximin and Maximax criteria}

		\CriterionTask
		{subsec:promotional-campaign-c5}{5}
		{Constructs opportunity loss tables and applies the Minimax regret criterion}
	\end{enumerate}
	
	\subsubsection*{C4 - Applies the Maximin and Maximax criteria}
	\phantomsection
	\label{subsec:promotional-campaign-c4}
	
	\[
	\begin{aligned}
		\max(\text{Campaign A}) &= \max\{100,20\}=100,\\
		\max(\text{Campaign B}) &= \max\{80,30\}=80.
	\end{aligned}
	\]
	\[
	\max\{\max(\text{Campaign A}),\max(\text{Campaign B})\}
	=\max\{100,80\}
	=100.
	\]
	The \GTenToneHL{best}{Maximax} choice is Campaign A because it has the highest possible \GTenToneHL{payoff}{payoff}.
	
	\[
	\begin{aligned}
		\min(\text{Campaign A}) &= \min\{100,20\}=20,\\
		\min(\text{Campaign B}) &= \min\{80,30\}=30.
	\end{aligned}
	\]
	\[
	\max\{\min(\text{Campaign A}),\min(\text{Campaign B})\}
	=\max\{20,30\}
	=30.
	\]
	The \GTenToneHL{worst}{Maximin} choice is Campaign B because it has the largest worst-case \GTenToneHL{payoff}{payoff}.
	
	\subsubsection*{C5 - Constructs opportunity loss tables and applies the Minimax regret criterion}
	\phantomsection
	\label{subsec:promotional-campaign-c5}
	
	Best \GTenToneHL{payoff}{payoff} in each state:
	\[
	\begin{aligned}
		\text{Strong response: }&\max\{100,80\}=100,\\
		\text{Weak response: }&\max\{20,30\}=30.
	\end{aligned}
	\]
	
	Regret table (best \GTenToneHL{payoff}{payoff} \( - \) \GTenToneHL{payoff}{payoff}):
	\begin{center}
		\begin{tabular}{l c c c}
			\toprule
			Alternative & Strong response & Weak response & \GTenToneHL{maxregret}{Maximum} \GTenToneHL{maxregret}{regret} \\
			\midrule
			Campaign A & \(100-100=0\) & \(30-20=10\) & \GTenToneHL{maxregret}{10} \\
			Campaign B & \(100-80=20\) & \(30-30=0\)  & \GTenToneHL{maxregret}{20} \\
			\bottomrule
		\end{tabular}
	\end{center}
	
	The \GTenToneHL{maxregret}{Minimax} \GTenToneHL{maxregret}{Regret} choice is Campaign A because it has the smallest maximum
	\GTenToneHL{regret}{regret}, \(10\).
	
	\subsubsection*{Decisions summary}
	
	\begin{center}
		\begin{tabular}{l c}
			\toprule
			Criterion & Final decision \\
			\midrule
			\GTenToneHL{best}{Maximax}        & Campaign A \\
			\GTenToneHL{worst}{Maximin}        & Campaign B \\
			\GTenToneHL{maxregret}{Minimax} \GTenToneHL{maxregret}{Regret} & Campaign A \\
			\bottomrule
		\end{tabular}
	\end{center}
	
	\secInicio[sec:promotional-campaign][sec:problems-2x2]
	
	
	% ----------------------------------------------------------------
	\subsection{Equipment Lease}
	\label{sec:equipment-lease}
	
	A company must choose between Lease A and Lease B. Market conditions can be
	favourable or unfavourable. The final \GTenToneHL{payoff}{payoff} table, in thousands of dollars, is
	given below. Use \GTenToneHL{best}{Maximax}, \GTenToneHL{worst}{Maximin}, and \GTenToneHL{maxregret}{Minimax} \GTenToneHL{maxregret}{Regret} to select a lease.
	
	\begin{center}
		\textit{Payoff table} \\
		\begin{tabular}{l c c}
			\toprule
			State of nature & Lease A & Lease B \\
			\midrule
			Favourable market   & \GTenToneHL{payoff}{70} & \GTenToneHL{payoff}{60} \\
			Unfavourable market & \GTenToneHL{payoff}{40} & \GTenToneHL{payoff}{30} \\
			\bottomrule
		\end{tabular}
	\end{center}
	
	\textbf{Questions/Tasks.}
	\begin{enumerate}
		\CriterionTask
		{subsec:equipment-lease-c4}{4}
		{Applies the Maximin and Maximax criteria}

		\CriterionTask
		{subsec:equipment-lease-c5}{5}
		{Constructs opportunity loss tables and applies the Minimax regret criterion}
	\end{enumerate}
	
	\subsubsection*{C4 - Applies the Maximin and Maximax criteria}
	\phantomsection
	\label{subsec:equipment-lease-c4}
	
	\[
	\begin{aligned}
		\max(\text{Lease A}) &= \max\{70,40\}=70,\\
		\max(\text{Lease B}) &= \max\{60,30\}=60.
	\end{aligned}
	\]
	\[
	\max\{\max(\text{Lease A}),\max(\text{Lease B})\}
	=\max\{70,60\}
	=70.
	\]
	The \GTenToneHL{best}{Maximax} choice is Lease A because it has the highest possible \GTenToneHL{payoff}{payoff}.
	
	\[
	\begin{aligned}
		\min(\text{Lease A}) &= \min\{70,40\}=40,\\
		\min(\text{Lease B}) &= \min\{60,30\}=30.
	\end{aligned}
	\]
	\[
	\max\{\min(\text{Lease A}),\min(\text{Lease B})\}
	=\max\{40,30\}
	=40.
	\]
	The \GTenToneHL{worst}{Maximin} choice is Lease A because it has the largest worst-case \GTenToneHL{payoff}{payoff}.
	
	\subsubsection*{C5 - Constructs opportunity loss tables and applies the Minimax regret criterion}
	\phantomsection
	\label{subsec:equipment-lease-c5}
	
	Best \GTenToneHL{payoff}{payoff} in each state:
	\[
	\begin{aligned}
		\text{Favourable market: }&\max\{70,60\}=70,\\
		\text{Unfavourable market: }&\max\{40,30\}=40.
	\end{aligned}
	\]
	
	Regret table (best \GTenToneHL{payoff}{payoff} \( - \) \GTenToneHL{payoff}{payoff}):
	\begin{center}
		\begin{tabular}{l c c c}
			\toprule
			Alternative & Favourable & Unfavourable & \GTenToneHL{maxregret}{Maximum} \GTenToneHL{maxregret}{regret} \\
			\midrule
			Lease A & \(70-70=0\)  & \(40-40=0\)  & \GTenToneHL{maxregret}{0} \\
			Lease B & \(70-60=10\) & \(40-30=10\) & \GTenToneHL{maxregret}{10} \\
			\bottomrule
		\end{tabular}
	\end{center}
	
	The \GTenToneHL{maxregret}{Minimax} \GTenToneHL{maxregret}{Regret} choice is Lease A because it has the smallest maximum
	\GTenToneHL{regret}{regret}, \(0\).
	
	\subsubsection*{Decisions summary}
	
	\begin{center}
		\begin{tabular}{l c}
			\toprule
			Criterion & Final decision \\
			\midrule
			\GTenToneHL{best}{Maximax}        & Lease A \\
			\GTenToneHL{worst}{Maximin}        & Lease A \\
			\GTenToneHL{maxregret}{Minimax} \GTenToneHL{maxregret}{Regret} & Lease A \\
			\bottomrule
		\end{tabular}
	\end{center}
	
	\secInicio[sec:equipment-lease][sec:problems-2x2]
	
	% ================================================================
	\section{Problems \(2\times 3\)}
	\label{sec:problems-2x3}
	
	\SectionProblemsTableOfContents{
		\SectionProblemLink{sec:production-plan}{Production Plan (ABB)} \\
		\SectionProblemLink{sec:marketing-strategy}{Marketing Strategy (ABA)} \\
		\SectionProblemLink{sec:event-venue}{Event Venue (ABC)} \\
		\SectionProblemLink{sec:distribution-plan}{Distribution Plan (AAA)}
	}
	
	% ----------------------------------------------------------------
	\subsection{Production Plan}
	\label{sec:production-plan}
	
	A manufacturer must choose Production Plan A, Plan B, or Plan C. Demand can be
	strong or weak. The final \GTenToneHL{payoff}{payoff} table, in thousands of dollars, is given below.
	Use \GTenToneHL{best}{Maximax}, \GTenToneHL{worst}{Maximin}, and \GTenToneHL{maxregret}{Minimax} \GTenToneHL{maxregret}{Regret} to select a plan.
	
	\begin{center}
		\textit{Payoff table} \\
		\begin{tabular}{l c c c}
			\toprule
			State of nature & Plan A & Plan B & Plan C \\
			\midrule
			Strong demand & \GTenToneHL{payoff}{100} & \GTenToneHL{payoff}{70} & \GTenToneHL{payoff}{80} \\
			Weak demand   & \GTenToneHL{payoff}{0}   & \GTenToneHL{payoff}{60} & \GTenToneHL{payoff}{20} \\
			\bottomrule
		\end{tabular}
	\end{center}
	
	\textbf{Questions/Tasks.}
	\begin{enumerate}
		\CriterionTask
		{subsec:production-plan-c4}{4}
		{Applies the Maximin and Maximax criteria}

		\CriterionTask
		{subsec:production-plan-c5}{5}
		{Constructs opportunity loss tables and applies the Minimax regret criterion}
	\end{enumerate}
	
	\subsubsection*{C4 - Applies the Maximin and Maximax criteria}
	\phantomsection
	\label{subsec:production-plan-c4}
	
	\[
	\begin{aligned}
		\max(\text{Plan A}) &= \max\{100,0\}=100,\\
		\max(\text{Plan B}) &= \max\{70,60\}=70,\\
		\max(\text{Plan C}) &= \max\{80,20\}=80.
	\end{aligned}
	\]
	\[
	\max\{\max(\text{Plan A}),\max(\text{Plan B}),\max(\text{Plan C})\}
	=\max\{100,70,80\}
	=100.
	\]
	The \GTenToneHL{best}{Maximax} choice is Plan A because it has the highest possible \GTenToneHL{payoff}{payoff}.
	
	\[
	\begin{aligned}
		\min(\text{Plan A}) &= \min\{100,0\}=0,\\
		\min(\text{Plan B}) &= \min\{70,60\}=60,\\
		\min(\text{Plan C}) &= \min\{80,20\}=20.
	\end{aligned}
	\]
	\[
	\max\{\min(\text{Plan A}),\min(\text{Plan B}),\min(\text{Plan C})\}
	=\max\{0,60,20\}
	=60.
	\]
	The \GTenToneHL{worst}{Maximin} choice is Plan B because it has the largest worst-case \GTenToneHL{payoff}{payoff}.
	
	\subsubsection*{C5 - Constructs opportunity loss tables and applies the Minimax regret criterion}
	\phantomsection
	\label{subsec:production-plan-c5}
	
	Best \GTenToneHL{payoff}{payoff} in each state:
	\[
	\begin{aligned}
		\text{Strong demand: }&\max\{100,70,80\}=100,\\
		\text{Weak demand: }&\max\{0,60,20\}=60.
	\end{aligned}
	\]
	
	Regret table (best \GTenToneHL{payoff}{payoff} \( - \) \GTenToneHL{payoff}{payoff}):
	\begin{center}
		\begin{tabular}{l c c c}
			\toprule
			Alternative & Strong demand & Weak demand & \GTenToneHL{maxregret}{Maximum} \GTenToneHL{maxregret}{regret} \\
			\midrule
			Plan A & \(100-100=0\) & \(60-0=60\)   & \GTenToneHL{maxregret}{60} \\
			Plan B & \(100-70=30\) & \(60-60=0\)  & \GTenToneHL{maxregret}{30} \\
			Plan C & \(100-80=20\) & \(60-20=40\) & \GTenToneHL{maxregret}{40} \\
			\bottomrule
		\end{tabular}
	\end{center}
	
	The \GTenToneHL{maxregret}{Minimax} \GTenToneHL{maxregret}{Regret} choice is Plan B because it has the smallest maximum
	\GTenToneHL{regret}{regret}, \(30\).
	
	\subsubsection*{Decisions summary}
	
	\begin{center}
		\begin{tabular}{l c}
			\toprule
			Criterion & Final decision \\
			\midrule
			\GTenToneHL{best}{Maximax}        & Plan A \\
			\GTenToneHL{worst}{Maximin}        & Plan B \\
			\GTenToneHL{maxregret}{Minimax} \GTenToneHL{maxregret}{Regret} & Plan B \\
			\bottomrule
		\end{tabular}
	\end{center}
	
	\secInicio[sec:production-plan][sec:problems-2x3]
	
	% ----------------------------------------------------------------
	\subsection{Marketing Strategy}
	\label{sec:marketing-strategy}
	
	A retailer must select Strategy A, Strategy B, or Strategy C for a new marketing
	campaign. Customer response can be strong or weak. The final \GTenToneHL{payoff}{payoff} table, in
	thousands of dollars, is given below. Use \GTenToneHL{best}{Maximax}, \GTenToneHL{worst}{Maximin}, and \GTenToneHL{maxregret}{Minimax} \GTenToneHL{maxregret}{Regret}
	to select a strategy.
	
	\begin{center}
		\textit{Payoff table} \\
		\begin{tabular}{l c c c}
			\toprule
			State of nature & Strategy A & Strategy B & Strategy C \\
			\midrule
			Strong response & \GTenToneHL{rate}{100} & \GTenToneHL{rate}{80} & \GTenToneHL{rate}{70} \\
			Weak response   & \GTenToneHL{rate}{40}  & \GTenToneHL{rate}{50} & \GTenToneHL{rate}{45} \\
			\bottomrule
		\end{tabular}
	\end{center}
	
	\textbf{Questions/Tasks.}
	\begin{enumerate}
		\CriterionTask
		{subsec:marketing-strategy-c4}{4}
		{Applies the Maximin and Maximax criteria}

		\CriterionTask
		{subsec:marketing-strategy-c5}{5}
		{Constructs opportunity loss tables and applies the Minimax regret criterion}
	\end{enumerate}
	
	\subsubsection*{C4 - Applies the Maximin and Maximax criteria}
	\phantomsection
	\label{subsec:marketing-strategy-c4}
	
	\[
	\begin{aligned}
		\max(\text{Strategy A}) &= \max\{100,40\}=100,\\
		\max(\text{Strategy B}) &= \max\{80,50\}=80,\\
		\max(\text{Strategy C}) &= \max\{70,45\}=70.
	\end{aligned}
	\]
	\[
	\max\{\max(\text{Strategy A}),\max(\text{Strategy B}),
	\max(\text{Strategy C})\}
	=\max\{100,80,70\}
	=100.
	\]
	The \GTenToneHL{best}{Maximax} choice is Strategy A because it has the highest possible \GTenToneHL{payoff}{payoff}.
	
	\[
	\begin{aligned}
		\min(\text{Strategy A}) &= \min\{100,40\}=40,\\
		\min(\text{Strategy B}) &= \min\{80,50\}=50,\\
		\min(\text{Strategy C}) &= \min\{70,45\}=45.
	\end{aligned}
	\]
	\[
	\max\{\min(\text{Strategy A}),\min(\text{Strategy B}),
	\min(\text{Strategy C})\}
	=\max\{40,50,45\}
	=50.
	\]
	The \GTenToneHL{worst}{Maximin} choice is Strategy B because it has the largest worst-case \GTenToneHL{payoff}{payoff}.
	
	\subsubsection*{C5 - Constructs opportunity loss tables and applies the Minimax regret criterion}
	\phantomsection
	\label{subsec:marketing-strategy-c5}
	
	Best \GTenToneHL{payoff}{payoff} in each state:
	\[
	\begin{aligned}
		\text{Strong response: }&\max\{100,80,70\}=100,\\
		\text{Weak response: }&\max\{40,50,45\}=50.
	\end{aligned}
	\]
	
	Regret table (best \GTenToneHL{payoff}{payoff} \( - \) \GTenToneHL{payoff}{payoff}):
	\begin{center}
		\begin{tabular}{l c c c}
			\toprule
			Alternative & Strong response & Weak response & \GTenToneHL{maxregret}{Maximum} \GTenToneHL{maxregret}{regret} \\
			\midrule
			Strategy A & \(100-100=0\) & \(50-40=10\) & \GTenToneHL{maxregret}{10} \\
			Strategy B & \(100-80=20\) & \(50-50=0\)  & \GTenToneHL{maxregret}{20} \\
			Strategy C & \(100-70=30\) & \(50-45=5\)  & \GTenToneHL{maxregret}{30} \\
			\bottomrule
		\end{tabular}
	\end{center}
	
	The \GTenToneHL{maxregret}{Minimax} \GTenToneHL{maxregret}{Regret} choice is Strategy A because it has the smallest maximum
	\GTenToneHL{regret}{regret}, \(10\).
	
	\subsubsection*{Decisions summary}
	
	\begin{center}
		\begin{tabular}{l c}
			\toprule
			Criterion & Final decision \\
			\midrule
			\GTenToneHL{best}{Maximax}        & Strategy A \\
			\GTenToneHL{worst}{Maximin}        & Strategy B \\
			\GTenToneHL{maxregret}{Minimax} \GTenToneHL{maxregret}{Regret} & Strategy A \\
			\bottomrule
		\end{tabular}
	\end{center}
	
	\secInicio[sec:marketing-strategy][sec:problems-2x3]
	
	
	% ----------------------------------------------------------------
	\subsection{Event Venue}
	\label{sec:event-venue}
	
	An event organiser must select Venue A, Venue B, or Venue C. Attendance can be
	high or low. The final \GTenToneHL{payoff}{payoff} table, in thousands of dollars, is given below.
	Use \GTenToneHL{best}{Maximax}, \GTenToneHL{worst}{Maximin}, and \GTenToneHL{maxregret}{Minimax} \GTenToneHL{maxregret}{Regret} to select a venue.
	
	\begin{center}
		\textit{Payoff table} \\
		\begin{tabular}{l c c c}
			\toprule
			State of nature & Venue A & Venue B & Venue C \\
			\midrule
			High attendance & \GTenToneHL{payoff}{100} & \GTenToneHL{payoff}{70} & \GTenToneHL{payoff}{80} \\
			Low attendance  & \GTenToneHL{payoff}{0}   & \GTenToneHL{payoff}{60} & \GTenToneHL{payoff}{50} \\
			\bottomrule
		\end{tabular}
	\end{center}
	
	\textbf{Questions/Tasks.}
	\begin{enumerate}
		\CriterionTask
		{subsec:event-venue-c4}{4}
		{Applies the Maximin and Maximax criteria}

		\CriterionTask
		{subsec:event-venue-c5}{5}
		{Constructs opportunity loss tables and applies the Minimax regret criterion}
	\end{enumerate}
	
	\subsubsection*{C4 - Applies the Maximin and Maximax criteria}
	\phantomsection
	\label{subsec:event-venue-c4}
	
	\[
	\begin{aligned}
		\max(\text{Venue A}) &= \max\{100,0\}=100,\\
		\max(\text{Venue B}) &= \max\{70,60\}=70,\\
		\max(\text{Venue C}) &= \max\{80,50\}=80.
	\end{aligned}
	\]
	\[
	\max\{\max(\text{Venue A}),\max(\text{Venue B}),\max(\text{Venue C})\}
	=\max\{100,70,80\}
	=100.
	\]
	The \GTenToneHL{best}{Maximax} choice is Venue A because it has the highest possible \GTenToneHL{payoff}{payoff}.
	
	\[
	\begin{aligned}
		\min(\text{Venue A}) &= \min\{100,0\}=0,\\
		\min(\text{Venue B}) &= \min\{70,60\}=60,\\
		\min(\text{Venue C}) &= \min\{80,50\}=50.
	\end{aligned}
	\]
	\[
	\max\{\min(\text{Venue A}),\min(\text{Venue B}),\min(\text{Venue C})\}
	=\max\{0,60,50\}
	=60.
	\]
	The \GTenToneHL{worst}{Maximin} choice is Venue B because it has the largest worst-case \GTenToneHL{payoff}{payoff}.
	
	\subsubsection*{C5 - Constructs opportunity loss tables and applies the Minimax regret criterion}
	\phantomsection
	\label{subsec:event-venue-c5}
	
	Best \GTenToneHL{payoff}{payoff} in each state:
	\[
	\begin{aligned}
		\text{High attendance: }&\max\{100,70,80\}=100,\\
		\text{Low attendance: }&\max\{0,60,50\}=60.
	\end{aligned}
	\]
	
	Regret table (best \GTenToneHL{payoff}{payoff} \( - \) \GTenToneHL{payoff}{payoff}):
	\begin{center}
		\begin{tabular}{l c c c}
			\toprule
			Alternative & High attendance & Low attendance & \GTenToneHL{maxregret}{Maximum} \GTenToneHL{maxregret}{regret} \\
			\midrule
			Venue A & \(100-100=0\) & \(60-0=60\)   & \GTenToneHL{maxregret}{60} \\
			Venue B & \(100-70=30\) & \(60-60=0\)  & \GTenToneHL{maxregret}{30} \\
			Venue C & \(100-80=20\) & \(60-50=10\) & \GTenToneHL{maxregret}{20} \\
			\bottomrule
		\end{tabular}
	\end{center}
	
	The \GTenToneHL{maxregret}{Minimax} \GTenToneHL{maxregret}{Regret} choice is Venue C because it has the smallest maximum
	\GTenToneHL{regret}{regret}, \(20\).
	
	\subsubsection*{Decisions summary}
	
	\begin{center}
		\begin{tabular}{l c}
			\toprule
			Criterion & Final decision \\
			\midrule
			\GTenToneHL{best}{Maximax}        & Venue A \\
			\GTenToneHL{worst}{Maximin}        & Venue B \\
			\GTenToneHL{maxregret}{Minimax} \GTenToneHL{maxregret}{Regret} & Venue C \\
			\bottomrule
		\end{tabular}
	\end{center}
	
	\secInicio[sec:event-venue][sec:problems-2x3]
	
	
	% ----------------------------------------------------------------
	\subsection{Distribution Plan}
	\label{sec:distribution-plan}
	
	A company must select Distribution Plan A, Plan B, or Plan C. Market conditions
	can be favourable or unfavourable. The final \GTenToneHL{payoff}{payoff} table, in thousands of
	dollars, is given below. Use \GTenToneHL{best}{Maximax}, \GTenToneHL{worst}{Maximin}, and \GTenToneHL{maxregret}{Minimax} \GTenToneHL{maxregret}{Regret} to select a
	distribution plan.
	
	\begin{center}
		\textit{Payoff table} \\
		\begin{tabular}{l c c c}
			\toprule
			State of nature & Plan A & Plan B & Plan C \\
			\midrule
			Favourable market   & \GTenToneHL{payoff}{90} & \GTenToneHL{payoff}{80} & \GTenToneHL{payoff}{70} \\
			Unfavourable market & \GTenToneHL{payoff}{60} & \GTenToneHL{payoff}{40} & \GTenToneHL{payoff}{50} \\
			\bottomrule
		\end{tabular}
	\end{center}
	
	\textbf{Questions/Tasks.}
	\begin{enumerate}
		\CriterionTask
		{subsec:distribution-plan-c4}{4}
		{Applies the Maximin and Maximax criteria}

		\CriterionTask
		{subsec:distribution-plan-c5}{5}
		{Constructs opportunity loss tables and applies the Minimax regret criterion}
	\end{enumerate}
	
	\subsubsection*{C4 - Applies the Maximin and Maximax criteria}
	\phantomsection
	\label{subsec:distribution-plan-c4}
	
	\[
	\begin{aligned}
		\max(\text{Plan A}) &= \max\{90,60\}=90,\\
		\max(\text{Plan B}) &= \max\{80,40\}=80,\\
		\max(\text{Plan C}) &= \max\{70,50\}=70.
	\end{aligned}
	\]
	\[
	\max\{\max(\text{Plan A}),\max(\text{Plan B}),\max(\text{Plan C})\}
	=\max\{90,80,70\}
	=90.
	\]
	The \GTenToneHL{best}{Maximax} choice is Plan A because it has the highest possible \GTenToneHL{payoff}{payoff}.
	
	\[
	\begin{aligned}
		\min(\text{Plan A}) &= \min\{90,60\}=60,\\
		\min(\text{Plan B}) &= \min\{80,40\}=40,\\
		\min(\text{Plan C}) &= \min\{70,50\}=50.
	\end{aligned}
	\]
	\[
	\max\{\min(\text{Plan A}),\min(\text{Plan B}),\min(\text{Plan C})\}
	=\max\{60,40,50\}
	=60.
	\]
	The \GTenToneHL{worst}{Maximin} choice is Plan A because it has the largest worst-case \GTenToneHL{payoff}{payoff}.
	
	\subsubsection*{C5 - Constructs opportunity loss tables and applies the Minimax regret criterion}
	\phantomsection
	\label{subsec:distribution-plan-c5}
	
	Best \GTenToneHL{payoff}{payoff} in each state:
	\[
	\begin{aligned}
		\text{Favourable market: }&\max\{90,80,70\}=90,\\
		\text{Unfavourable market: }&\max\{60,40,50\}=60.
	\end{aligned}
	\]
	
	Regret table (best \GTenToneHL{payoff}{payoff} \( - \) \GTenToneHL{payoff}{payoff}):
	\begin{center}
		\begin{tabular}{l c c c}
			\toprule
			Alternative & Favourable & Unfavourable & \GTenToneHL{maxregret}{Maximum} \GTenToneHL{maxregret}{regret} \\
			\midrule
			Plan A & \(90-90=0\)  & \(60-60=0\)  & \GTenToneHL{maxregret}{0} \\
			Plan B & \(90-80=10\) & \(60-40=20\) & \GTenToneHL{maxregret}{20} \\
			Plan C & \(90-70=20\) & \(60-50=10\) & \GTenToneHL{maxregret}{20} \\
			\bottomrule
		\end{tabular}
	\end{center}
	
	The \GTenToneHL{maxregret}{Minimax} \GTenToneHL{maxregret}{Regret} choice is Plan A because it has the smallest maximum
	\GTenToneHL{regret}{regret}, \(0\).
	
	\subsubsection*{Decisions summary}
	
	\begin{center}
		\begin{tabular}{l c}
			\toprule
			Criterion & Final decision \\
			\midrule
			\GTenToneHL{best}{Maximax}        & Plan A \\
			\GTenToneHL{worst}{Maximin}        & Plan A \\
			\GTenToneHL{maxregret}{Minimax} \GTenToneHL{maxregret}{Regret} & Plan A \\
			\bottomrule
		\end{tabular}
	\end{center}
	
	\secInicio[sec:distribution-plan][sec:problems-2x3]
	
	
	% ================================================================
	\section{Problems \(3\times 2\)}
	\label{sec:problems-3x2}
	
	\SectionProblemsTableOfContents{
		\SectionProblemLink{sec:routing-system}{Routing System (AAB)} \\
		\SectionProblemLink{sec:service-model-3x2}{Service Model (ABA)} \\
		\SectionProblemLink{sec:inventory-policy-3x2}{Inventory Policy (ABB)} \\
		\SectionProblemLink{sec:logistics-contract-3x2}{Logistics Contract (AAA)}
	}
	
	% ----------------------------------------------------------------
	\subsection{Routing System}
	\label{sec:routing-system}
	
	A shipping company must choose between System A and System B. Three states of
	nature represent different operating conditions. The final \GTenToneHL{payoff}{payoff} table, in
	thousands of dollars, is given below. Use \GTenToneHL{best}{Maximax}, \GTenToneHL{worst}{Maximin}, and \GTenToneHL{maxregret}{Minimax} \GTenToneHL{maxregret}{Regret}
	to select a routing system.
	
	\begin{center}
		\textit{Payoff table} \\
		\begin{tabular}{l c c}
			\toprule
			State of nature & System A & System B \\
			\midrule
			State \(\GTenToneHL{state}{S_1}\) & \GTenToneHL{payoff}{100} & \GTenToneHL{payoff}{70} \\
			State \(\GTenToneHL{state}{S_2}\) & \GTenToneHL{payoff}{50}  & \GTenToneHL{payoff}{49} \\
			State \(\GTenToneHL{state}{S_3}\) & \GTenToneHL{payoff}{50}  & \GTenToneHL{payoff}{90} \\
			\bottomrule
		\end{tabular}
	\end{center}
	
	\textbf{Questions/Tasks.}
	\begin{enumerate}
		\CriterionTask
		{subsec:routing-system-c4}{4}
		{Applies the Maximin and Maximax criteria}

		\CriterionTask
		{subsec:routing-system-c5}{5}
		{Constructs opportunity loss tables and applies the Minimax regret criterion}
	\end{enumerate}
	
	\subsubsection*{C4 - Applies the Maximin and Maximax criteria}
	\phantomsection
	\label{subsec:routing-system-c4}
	
	\[
	\begin{aligned}
		\max(\text{System A}) &= \max\{100,50,50\}=100,\\
		\max(\text{System B}) &= \max\{70,49,90\}=90.
	\end{aligned}
	\]
	\[
	\max\{\max(\text{System A}),\max(\text{System B})\}
	=\max\{100,90\}
	=100.
	\]
	The \GTenToneHL{best}{Maximax} choice is System A because it has the highest possible \GTenToneHL{payoff}{payoff}.
	
	\[
	\begin{aligned}
		\min(\text{System A}) &= \min\{100,50,50\}=50,\\
		\min(\text{System B}) &= \min\{70,49,90\}=49.
	\end{aligned}
	\]
	\[
	\max\{\min(\text{System A}),\min(\text{System B})\}
	=\max\{50,49\}
	=50.
	\]
	The \GTenToneHL{worst}{Maximin} choice is System A because it has the largest worst-case \GTenToneHL{payoff}{payoff}.
	
	\subsubsection*{C5 - Constructs opportunity loss tables and applies the Minimax regret criterion}
	\phantomsection
	\label{subsec:routing-system-c5}
	
	Best \GTenToneHL{payoff}{payoff} in each state:
	\[
	\begin{aligned}
		\text{State }\GTenToneHL{state}{S_1}:\;&\max\{100,70\}=100,\\
		\text{State }\GTenToneHL{state}{S_2}:\;&\max\{50,49\}=50,\\
		\text{State }\GTenToneHL{state}{S_3}:\;&\max\{50,90\}=90.
	\end{aligned}
	\]
	
	Regret table (best \GTenToneHL{payoff}{payoff} \( - \) \GTenToneHL{payoff}{payoff}):
	\begin{center}
		\begin{tabular}{l c c c c}
			\toprule
			Alternative & \(\GTenToneHL{state}{S_1}\) & \(\GTenToneHL{state}{S_2}\) & \(\GTenToneHL{state}{S_3}\) & \GTenToneHL{maxregret}{Maximum} \GTenToneHL{maxregret}{regret} \\
			\midrule
			System A
			& \(100-100=0\)
			& \(50-50=0\)
			& \(90-50=40\)
			& \GTenToneHL{maxregret}{40} \\
			System B
			& \(100-70=30\)
			& \(50-49=1\)
			& \(90-90=0\)
			& \GTenToneHL{maxregret}{30} \\
			\bottomrule
		\end{tabular}
	\end{center}
	
	The \GTenToneHL{maxregret}{Minimax} \GTenToneHL{maxregret}{Regret} choice is System B because it has the smallest maximum
	\GTenToneHL{regret}{regret}, \(30\).
	
	\subsubsection*{Decisions summary}
	
	\begin{center}
		\begin{tabular}{l c}
			\toprule
			Criterion & Final decision \\
			\midrule
			\GTenToneHL{best}{Maximax}        & System A \\
			\GTenToneHL{worst}{Maximin}        & System A \\
			\GTenToneHL{maxregret}{Minimax} \GTenToneHL{maxregret}{Regret} & System B \\
			\bottomrule
		\end{tabular}
	\end{center}
	
	\secInicio[sec:routing-system][sec:problems-3x2]
	
	
	% ----------------------------------------------------------------
	\subsection{Service Model}
	\label{sec:service-model-3x2}
	
	A service company must choose between Service Model A and Service Model B.
	Demand can be high, moderate, or low. The final \GTenToneHL{payoff}{payoff} table, in thousands of
	dollars, is given below. Use \GTenToneHL{best}{Maximax}, \GTenToneHL{worst}{Maximin}, and \GTenToneHL{maxregret}{Minimax} \GTenToneHL{maxregret}{Regret} to select a
	service model.
	
	\begin{center}
		\textit{Payoff table} \\
		\begin{tabular}{l c c}
			\toprule
			State of nature & Model A & Model B \\
			\midrule
			High demand     & \GTenToneHL{payoff}{100} & \GTenToneHL{payoff}{80} \\
			Moderate demand & \GTenToneHL{payoff}{45}  & \GTenToneHL{payoff}{50} \\
			Low demand      & \GTenToneHL{payoff}{40}  & \GTenToneHL{payoff}{50} \\
			\bottomrule
		\end{tabular}
	\end{center}
	
	\textbf{Questions/Tasks.}
	\begin{enumerate}
		\CriterionTask
		{subsec:service-model-3x2-c4}{4}
		{Applies the Maximin and Maximax criteria}

		\CriterionTask
		{subsec:service-model-3x2-c5}{5}
		{Constructs opportunity loss tables and applies the Minimax regret criterion}
	\end{enumerate}
	
	\subsubsection*{C4 - Applies the Maximin and Maximax criteria}
	\phantomsection
	\label{subsec:service-model-3x2-c4}
	
	\[
	\begin{aligned}
		\max(\text{Model A}) &= \max\{100,45,40\}=100,\\
		\max(\text{Model B}) &= \max\{80,50,50\}=80.
	\end{aligned}
	\]
	\[
	\max\{\max(\text{Model A}),\max(\text{Model B})\}
	=\max\{100,80\}
	=100.
	\]
	The \GTenToneHL{best}{Maximax} choice is Model A because it has the highest possible \GTenToneHL{payoff}{payoff}.
	
	\[
	\begin{aligned}
		\min(\text{Model A}) &= \min\{100,45,40\}=40,\\
		\min(\text{Model B}) &= \min\{80,50,50\}=50.
	\end{aligned}
	\]
	\[
	\max\{\min(\text{Model A}),\min(\text{Model B})\}
	=\max\{40,50\}
	=50.
	\]
	The \GTenToneHL{worst}{Maximin} choice is Model B because it has the largest worst-case \GTenToneHL{payoff}{payoff}.
	
	\subsubsection*{C5 - Constructs opportunity loss tables and applies the Minimax regret criterion}
	\phantomsection
	\label{subsec:service-model-3x2-c5}
	
	Best \GTenToneHL{payoff}{payoff} in each state:
	\[
	\begin{aligned}
		\text{High demand: }&\max\{100,80\}=100,\\
		\text{Moderate demand: }&\max\{45,50\}=50,\\
		\text{Low demand: }&\max\{40,50\}=50.
	\end{aligned}
	\]
	
	Regret table (best \GTenToneHL{payoff}{payoff} \( - \) \GTenToneHL{payoff}{payoff}):
	\begin{center}
		\begin{tabular}{l c c c c}
			\toprule
			Alternative & High & Moderate & Low & \GTenToneHL{maxregret}{Maximum} \GTenToneHL{maxregret}{regret} \\
			\midrule
			Model A
			& \(100-100=0\)
			& \(50-45=5\)
			& \(50-40=10\)
			& \GTenToneHL{maxregret}{10} \\
			Model B
			& \(100-80=20\)
			& \(50-50=0\)
			& \(50-50=0\)
			& \GTenToneHL{maxregret}{20} \\
			\bottomrule
		\end{tabular}
	\end{center}
	
	The \GTenToneHL{maxregret}{Minimax} \GTenToneHL{maxregret}{Regret} choice is Model A because it has the smallest maximum
	\GTenToneHL{regret}{regret}, \(10\).
	
	\subsubsection*{Decisions summary}
	
	\begin{center}
		\begin{tabular}{l c}
			\toprule
			Criterion & Final decision \\
			\midrule
			\GTenToneHL{best}{Maximax}        & Model A \\
			\GTenToneHL{worst}{Maximin}        & Model B \\
			\GTenToneHL{maxregret}{Minimax} \GTenToneHL{maxregret}{Regret} & Model A \\
			\bottomrule
		\end{tabular}
	\end{center}
	
	\secInicio[sec:service-model-3x2][sec:problems-3x2]
	
	
	% ----------------------------------------------------------------
	\subsection{Inventory Policy}
	\label{sec:inventory-policy-3x2}
	
	A retailer must choose between Inventory Policy A and Inventory Policy B.
	Demand can be high, moderate, or low. The final \GTenToneHL{payoff}{payoff} table, in thousands of
	dollars, is given below. Use \GTenToneHL{best}{Maximax}, \GTenToneHL{worst}{Maximin}, and \GTenToneHL{maxregret}{Minimax} \GTenToneHL{maxregret}{Regret} to select an
	inventory policy.
	
	\begin{center}
		\textit{Payoff table} \\
		\begin{tabular}{l c c}
			\toprule
			State of nature & Policy A & Policy B \\
			\midrule
			High demand     & \GTenToneHL{payoff}{100} & \GTenToneHL{payoff}{70} \\
			Moderate demand & \GTenToneHL{payoff}{20}  & \GTenToneHL{payoff}{60} \\
			Low demand      & \GTenToneHL{payoff}{0}   & \GTenToneHL{payoff}{50} \\
			\bottomrule
		\end{tabular}
	\end{center}
	
	\textbf{Questions/Tasks.}
	\begin{enumerate}
		\CriterionTask
		{subsec:inventory-policy-3x2-c4}{4}
		{Applies the Maximin and Maximax criteria}

		\CriterionTask
		{subsec:inventory-policy-3x2-c5}{5}
		{Constructs opportunity loss tables and applies the Minimax regret criterion}
	\end{enumerate}
	
	\subsubsection*{C4 - Applies the Maximin and Maximax criteria}
	\phantomsection
	\label{subsec:inventory-policy-3x2-c4}
	
	\[
	\begin{aligned}
		\max(\text{Policy A}) &= \max\{100,20,0\}=100,\\
		\max(\text{Policy B}) &= \max\{70,60,50\}=70.
	\end{aligned}
	\]
	\[
	\max\{\max(\text{Policy A}),\max(\text{Policy B})\}
	=\max\{100,70\}
	=100.
	\]
	The \GTenToneHL{best}{Maximax} choice is Policy A because it has the highest possible \GTenToneHL{payoff}{payoff}.
	
	\[
	\begin{aligned}
		\min(\text{Policy A}) &= \min\{100,20,0\}=0,\\
		\min(\text{Policy B}) &= \min\{70,60,50\}=50.
	\end{aligned}
	\]
	\[
	\max\{\min(\text{Policy A}),\min(\text{Policy B})\}
	=\max\{0,50\}
	=50.
	\]
	The \GTenToneHL{worst}{Maximin} choice is Policy B because it has the largest worst-case \GTenToneHL{payoff}{payoff}.
	
	\subsubsection*{C5 - Constructs opportunity loss tables and applies the Minimax regret criterion}
	\phantomsection
	\label{subsec:inventory-policy-3x2-c5}
	
	Best \GTenToneHL{payoff}{payoff} in each state:
	\[
	\begin{aligned}
		\text{High demand: }&\max\{100,70\}=100,\\
		\text{Moderate demand: }&\max\{20,60\}=60,\\
		\text{Low demand: }&\max\{0,50\}=50.
	\end{aligned}
	\]
	
	Regret table (best \GTenToneHL{payoff}{payoff} \( - \) \GTenToneHL{payoff}{payoff}):
	\begin{center}
		\begin{tabular}{l c c c c}
			\toprule
			Alternative & High & Moderate & Low & \GTenToneHL{maxregret}{Maximum} \GTenToneHL{maxregret}{regret} \\
			\midrule
			Policy A
			& \(100-100=0\)
			& \(60-20=40\)
			& \(50-0=50\)
			& \GTenToneHL{maxregret}{50} \\
			Policy B
			& \(100-70=30\)
			& \(60-60=0\)
			& \(50-50=0\)
			& \GTenToneHL{maxregret}{30} \\
			\bottomrule
		\end{tabular}
	\end{center}
	
	The \GTenToneHL{maxregret}{Minimax} \GTenToneHL{maxregret}{Regret} choice is Policy B because it has the smallest maximum
	\GTenToneHL{regret}{regret}, \(30\).
	
	\subsubsection*{Decisions summary}
	
	\begin{center}
		\begin{tabular}{l c}
			\toprule
			Criterion & Final decision \\
			\midrule
			\GTenToneHL{best}{Maximax}        & Policy A \\
			\GTenToneHL{worst}{Maximin}        & Policy B \\
			\GTenToneHL{maxregret}{Minimax} \GTenToneHL{maxregret}{Regret} & Policy B \\
			\bottomrule
		\end{tabular}
	\end{center}
	
	\secInicio[sec:inventory-policy-3x2][sec:problems-3x2]
	
	
	% ----------------------------------------------------------------
	\subsection{Logistics Contract}
	\label{sec:logistics-contract-3x2}
	
	A company must choose between Logistics Contract A and Logistics Contract B.
	Market conditions can be favourable, stable, or unfavourable. The final \GTenToneHL{payoff}{payoff}
	table, in thousands of dollars, is given below. Use \GTenToneHL{best}{Maximax}, \GTenToneHL{worst}{Maximin}, and
	\GTenToneHL{maxregret}{Minimax} \GTenToneHL{maxregret}{Regret} to select a logistics contract.
	
	\begin{center}
		\textit{Payoff table} \\
		\begin{tabular}{l c c}
			\toprule
			State of nature & Contract A & Contract B \\
			\midrule
			Favourable market   & \GTenToneHL{payoff}{90} & \GTenToneHL{payoff}{80} \\
			Stable market       & \GTenToneHL{payoff}{70} & \GTenToneHL{payoff}{50} \\
			Unfavourable market & \GTenToneHL{payoff}{60} & \GTenToneHL{payoff}{40} \\
			\bottomrule
		\end{tabular}
	\end{center}
	
	\textbf{Questions/Tasks.}
	\begin{enumerate}
		\CriterionTask
		{subsec:logistics-contract-3x2-c4}{4}
		{Applies the Maximin and Maximax criteria}

		\CriterionTask
		{subsec:logistics-contract-3x2-c5}{5}
		{Constructs opportunity loss tables and applies the Minimax regret criterion}
	\end{enumerate}
	
	\subsubsection*{C4 - Applies the Maximin and Maximax criteria}
	\phantomsection
	\label{subsec:logistics-contract-3x2-c4}
	
	\[
	\begin{aligned}
		\max(\text{Contract A}) &= \max\{90,70,60\}=90,\\
		\max(\text{Contract B}) &= \max\{80,50,40\}=80.
	\end{aligned}
	\]
	\[
	\max\{\max(\text{Contract A}),\max(\text{Contract B})\}
	=\max\{90,80\}
	=90.
	\]
	The \GTenToneHL{best}{Maximax} choice is Contract A because it has the highest possible \GTenToneHL{payoff}{payoff}.
	
	\[
	\begin{aligned}
		\min(\text{Contract A}) &= \min\{90,70,60\}=60,\\
		\min(\text{Contract B}) &= \min\{80,50,40\}=40.
	\end{aligned}
	\]
	\[
	\max\{\min(\text{Contract A}),\min(\text{Contract B})\}
	=\max\{60,40\}
	=60.
	\]
	The \GTenToneHL{worst}{Maximin} choice is Contract A because it has the largest worst-case \GTenToneHL{payoff}{payoff}.
	
	\subsubsection*{C5 - Constructs opportunity loss tables and applies the Minimax regret criterion}
	\phantomsection
	\label{subsec:logistics-contract-3x2-c5}
	
	Best \GTenToneHL{payoff}{payoff} in each state:
	\[
	\begin{aligned}
		\text{Favourable market: }&\max\{90,80\}=90,\\
		\text{Stable market: }&\max\{70,50\}=70,\\
		\text{Unfavourable market: }&\max\{60,40\}=60.
	\end{aligned}
	\]
	
	Regret table (best \GTenToneHL{payoff}{payoff} \( - \) \GTenToneHL{payoff}{payoff}):
	\begin{center}
		\begin{tabular}{l c c c c}
			\toprule
			Alternative & Favourable & Stable & Unfavourable & \GTenToneHL{maxregret}{Maximum} \GTenToneHL{maxregret}{regret} \\
			\midrule
			Contract A
			& \(90-90=0\)
			& \(70-70=0\)
			& \(60-60=0\)
			& \GTenToneHL{maxregret}{0} \\
			Contract B
			& \(90-80=10\)
			& \(70-50=20\)
			& \(60-40=20\)
			& \GTenToneHL{maxregret}{20} \\
			\bottomrule
		\end{tabular}
	\end{center}
	
	The \GTenToneHL{maxregret}{Minimax} \GTenToneHL{maxregret}{Regret} choice is Contract A because it has the smallest maximum
	\GTenToneHL{regret}{regret}, \(0\).
	
	\subsubsection*{Decisions summary}
	
	\begin{center}
		\begin{tabular}{l c}
			\toprule
			Criterion & Final decision \\
			\midrule
			\GTenToneHL{best}{Maximax}        & Contract A \\
			\GTenToneHL{worst}{Maximin}        & Contract A \\
			\GTenToneHL{maxregret}{Minimax} \GTenToneHL{maxregret}{Regret} & Contract A \\
			\bottomrule
		\end{tabular}
	\end{center}
	
	\secInicio[sec:logistics-contract-3x2][sec:problems-3x2]
	
	
	% ================================================================
	\section{Problems \(3\times 3\)}
	\label{sec:problems-3x3}
	
	\SectionProblemsTableOfContents{
		\SectionProblemLink{sec:technology-platform-3x3}{Technology Platform (AAB)} \\
		\SectionProblemLink{sec:service-model-3x3}{Service Model (ABA)} \\
		\SectionProblemLink{sec:inventory-policy-3x3}{Inventory Policy (ABB)} \\
		\SectionProblemLink{sec:event-venue-3x3}{Event Venue (ABC)} \\
		\SectionProblemLink{sec:distribution-plan-3x3}{Distribution Plan (AAA)}
	}
	
	
	% ----------------------------------------------------------------
	\subsection{Technology Platform}
	\label{sec:technology-platform-3x3}
	
	A company must choose Technology Platform A, Platform B, or Platform C.
	Operating conditions can be favourable, stable, or disrupted. The final
	\GTenToneHL{payoff}{payoff} table, in thousands of dollars, is given below. Use \GTenToneHL{best}{Maximax}, \GTenToneHL{worst}{Maximin},
	and \GTenToneHL{maxregret}{Minimax} \GTenToneHL{maxregret}{Regret} to select a technology platform.
	
	\begin{center}
		\textit{Payoff table} \\
		\begin{tabular}{l c c c}
			\toprule
			State of nature & Platform A & Platform B & Platform C \\
			\midrule
			Favourable conditions & \GTenToneHL{payoff}{100} & \GTenToneHL{payoff}{70} & \GTenToneHL{payoff}{60} \\
			Stable conditions     & \GTenToneHL{payoff}{50}  & \GTenToneHL{payoff}{49} & \GTenToneHL{payoff}{40} \\
			Disrupted conditions  & \GTenToneHL{payoff}{50}  & \GTenToneHL{payoff}{90} & \GTenToneHL{payoff}{45} \\
			\bottomrule
		\end{tabular}
	\end{center}
	
	\textbf{Questions/Tasks.}
	\begin{enumerate}
		\CriterionTask
		{subsec:technology-platform-3x3-c4}{4}
		{Applies the Maximin and Maximax criteria}

		\CriterionTask
		{subsec:technology-platform-3x3-c5}{5}
		{Constructs opportunity loss tables and applies the Minimax regret criterion}
	\end{enumerate}
	
	\subsubsection*{C4 - Applies the Maximin and Maximax criteria}
	\phantomsection
	\label{subsec:technology-platform-3x3-c4}
	
	\[
	\begin{aligned}
		\max(\text{Platform A}) &= \max\{100,50,50\}=100,\\
		\max(\text{Platform B}) &= \max\{70,49,90\}=90,\\
		\max(\text{Platform C}) &= \max\{60,40,45\}=60.
	\end{aligned}
	\]
	\[
	\max\{\max(\text{Platform A}),\max(\text{Platform B}),
	\max(\text{Platform C})\}
	=\max\{100,90,60\}
	=100.
	\]
	The \GTenToneHL{best}{Maximax} choice is Platform A because it has the highest possible \GTenToneHL{payoff}{payoff}.
	
	\[
	\begin{aligned}
		\min(\text{Platform A}) &= \min\{100,50,50\}=50,\\
		\min(\text{Platform B}) &= \min\{70,49,90\}=49,\\
		\min(\text{Platform C}) &= \min\{60,40,45\}=40.
	\end{aligned}
	\]
	\[
	\max\{\min(\text{Platform A}),\min(\text{Platform B}),
	\min(\text{Platform C})\}
	=\max\{50,49,40\}
	=50.
	\]
	The \GTenToneHL{worst}{Maximin} choice is Platform A because it has the largest worst-case \GTenToneHL{payoff}{payoff}.
	
	\subsubsection*{C5 - Constructs opportunity loss tables and applies the Minimax regret criterion}
	\phantomsection
	\label{subsec:technology-platform-3x3-c5}
	
	Best \GTenToneHL{payoff}{payoff} in each state:
	\[
	\begin{aligned}
		\text{Favourable conditions: }&\max\{100,70,60\}=100,\\
		\text{Stable conditions: }&\max\{50,49,40\}=50,\\
		\text{Disrupted conditions: }&\max\{50,90,45\}=90.
	\end{aligned}
	\]
	
	Regret table (best \GTenToneHL{payoff}{payoff} \( - \) \GTenToneHL{payoff}{payoff}):
	\begin{center}
		\begin{tabular}{l c c c c}
			\toprule
			Alternative & Favourable & Stable & Disrupted & \GTenToneHL{maxregret}{Maximum} \GTenToneHL{maxregret}{regret} \\
			\midrule
			Platform A
			& \(100-100=0\)
			& \(50-50=0\)
			& \(90-50=40\)
			& \GTenToneHL{maxregret}{40} \\
			Platform B
			& \(100-70=30\)
			& \(50-49=1\)
			& \(90-90=0\)
			& \GTenToneHL{maxregret}{30} \\
			Platform C
			& \(100-60=40\)
			& \(50-40=10\)
			& \(90-45=45\)
			& \GTenToneHL{maxregret}{45} \\
			\bottomrule
		\end{tabular}
	\end{center}
	
	The \GTenToneHL{maxregret}{Minimax} \GTenToneHL{maxregret}{Regret} choice is Platform B because it has the smallest maximum
	\GTenToneHL{regret}{regret}, \(30\).
	
	\subsubsection*{Decisions summary}
	
	\begin{center}
		\begin{tabular}{l c}
			\toprule
			Criterion & Final decision \\
			\midrule
			\GTenToneHL{best}{Maximax}        & Platform A \\
			\GTenToneHL{worst}{Maximin}        & Platform A \\
			\GTenToneHL{maxregret}{Minimax} \GTenToneHL{maxregret}{Regret} & Platform B \\
			\bottomrule
		\end{tabular}
	\end{center}
	
	\secInicio[sec:technology-platform-3x3][sec:problems-3x3]
	
	
	% ----------------------------------------------------------------
	\subsection{Service Model}
	\label{sec:service-model-3x3}
	
	A service company must choose Service Model A, Model B, or Model C.
	Demand can be high, moderate, or low. The final \GTenToneHL{payoff}{payoff} table, in thousands
	of dollars, is given below. Use \GTenToneHL{best}{Maximax}, \GTenToneHL{worst}{Maximin}, and \GTenToneHL{maxregret}{Minimax} \GTenToneHL{maxregret}{Regret} to
	select a service model.
	
	\begin{center}
		\textit{Payoff table} \\
		\begin{tabular}{l c c c}
			\toprule
			State of nature & Model A & Model B & Model C \\
			\midrule
			High demand     & \GTenToneHL{payoff}{100} & \GTenToneHL{payoff}{80} & \GTenToneHL{payoff}{70} \\
			Moderate demand & \GTenToneHL{payoff}{45}  & \GTenToneHL{payoff}{50} & \GTenToneHL{payoff}{35} \\
			Low demand      & \GTenToneHL{payoff}{40}  & \GTenToneHL{payoff}{50} & \GTenToneHL{payoff}{45} \\
			\bottomrule
		\end{tabular}
	\end{center}
	
	\textbf{Questions/Tasks.}
	\begin{enumerate}
		\CriterionTask
		{subsec:service-model-3x3-c4}{4}
		{Applies the Maximin and Maximax criteria}

		\CriterionTask
		{subsec:service-model-3x3-c5}{5}
		{Constructs opportunity loss tables and applies the Minimax regret criterion}
	\end{enumerate}
	
	\subsubsection*{C4 - Applies the Maximin and Maximax criteria}
	\phantomsection
	\label{subsec:service-model-3x3-c4}
	
	\[
	\begin{aligned}
		\max(\text{Model A}) &= \max\{100,45,40\}=100,\\
		\max(\text{Model B}) &= \max\{80,50,50\}=80,\\
		\max(\text{Model C}) &= \max\{70,35,45\}=70.
	\end{aligned}
	\]
	\[
	\max\{\max(\text{Model A}),\max(\text{Model B}),
	\max(\text{Model C})\}
	=\max\{100,80,70\}
	=100.
	\]
	The \GTenToneHL{best}{Maximax} choice is Model A because it has the highest possible \GTenToneHL{payoff}{payoff}.
	
	\[
	\begin{aligned}
		\min(\text{Model A}) &= \min\{100,45,40\}=40,\\
		\min(\text{Model B}) &= \min\{80,50,50\}=50,\\
		\min(\text{Model C}) &= \min\{70,35,45\}=35.
	\end{aligned}
	\]
	\[
	\max\{\min(\text{Model A}),\min(\text{Model B}),
	\min(\text{Model C})\}
	=\max\{40,50,35\}
	=50.
	\]
	The \GTenToneHL{worst}{Maximin} choice is Model B because it has the largest worst-case \GTenToneHL{payoff}{payoff}.
	
	\subsubsection*{C5 - Constructs opportunity loss tables and applies the Minimax regret criterion}
	\phantomsection
	\label{subsec:service-model-3x3-c5}
	
	Best \GTenToneHL{payoff}{payoff} in each state:
	\[
	\begin{aligned}
		\text{High demand: }&\max\{100,80,70\}=100,\\
		\text{Moderate demand: }&\max\{45,50,35\}=50,\\
		\text{Low demand: }&\max\{40,50,45\}=50.
	\end{aligned}
	\]
	
	Regret table (best \GTenToneHL{payoff}{payoff} \( - \) \GTenToneHL{payoff}{payoff}):
	\begin{center}
		\begin{tabular}{l c c c c}
			\toprule
			Alternative & High & Moderate & Low & \GTenToneHL{maxregret}{Maximum} \GTenToneHL{maxregret}{regret} \\
			\midrule
			Model A
			& \(100-100=0\)
			& \(50-45=5\)
			& \(50-40=10\)
			& \GTenToneHL{maxregret}{10} \\
			Model B
			& \(100-80=20\)
			& \(50-50=0\)
			& \(50-50=0\)
			& \GTenToneHL{maxregret}{20} \\
			Model C
			& \(100-70=30\)
			& \(50-35=15\)
			& \(50-45=5\)
			& \GTenToneHL{maxregret}{30} \\
			\bottomrule
		\end{tabular}
	\end{center}
	
	The \GTenToneHL{maxregret}{Minimax} \GTenToneHL{maxregret}{Regret} choice is Model A because it has the smallest maximum
	\GTenToneHL{regret}{regret}, \(10\).
	
	\subsubsection*{Decisions summary}
	
	\begin{center}
		\begin{tabular}{l c}
			\toprule
			Criterion & Final decision \\
			\midrule
			\GTenToneHL{best}{Maximax}        & Model A \\
			\GTenToneHL{worst}{Maximin}        & Model B \\
			\GTenToneHL{maxregret}{Minimax} \GTenToneHL{maxregret}{Regret} & Model A \\
			\bottomrule
		\end{tabular}
	\end{center}
	
	\secInicio[sec:service-model-3x3][sec:problems-3x3]
	
	
	% ----------------------------------------------------------------
	\subsection{Inventory Policy}
	\label{sec:inventory-policy-3x3}
	
	A retailer must choose Inventory Policy A, Policy B, or Policy C. Demand
	can be high, moderate, or low. The final \GTenToneHL{payoff}{payoff} table, in thousands of
	dollars, is given below. Use \GTenToneHL{best}{Maximax}, \GTenToneHL{worst}{Maximin}, and \GTenToneHL{maxregret}{Minimax} \GTenToneHL{maxregret}{Regret} to select
	an inventory policy.
	
	\begin{center}
		\textit{Payoff table} \\
		\begin{tabular}{l c c c}
			\toprule
			State of nature & Policy A & Policy B & Policy C \\
			\midrule
			High demand     & \GTenToneHL{payoff}{100} & \GTenToneHL{payoff}{70} & \GTenToneHL{payoff}{60} \\
			Moderate demand & \GTenToneHL{payoff}{20}  & \GTenToneHL{payoff}{60} & \GTenToneHL{payoff}{40} \\
			Low demand      & \GTenToneHL{payoff}{0}   & \GTenToneHL{payoff}{50} & \GTenToneHL{payoff}{30} \\
			\bottomrule
		\end{tabular}
	\end{center}
	
	\textbf{Questions/Tasks.}
	\begin{enumerate}
		\CriterionTask
		{subsec:inventory-policy-3x3-c4}{4}
		{Applies the Maximin and Maximax criteria}

		\CriterionTask
		{subsec:inventory-policy-3x3-c5}{5}
		{Constructs opportunity loss tables and applies the Minimax regret criterion}
	\end{enumerate}
	
	\subsubsection*{C4 - Applies the Maximin and Maximax criteria}
	\phantomsection
	\label{subsec:inventory-policy-3x3-c4}
	
	\[
	\begin{aligned}
		\max(\text{Policy A}) &= \max\{100,20,0\}=100,\\
		\max(\text{Policy B}) &= \max\{70,60,50\}=70,\\
		\max(\text{Policy C}) &= \max\{60,40,30\}=60.
	\end{aligned}
	\]
	\[
	\max\{\max(\text{Policy A}),\max(\text{Policy B}),
	\max(\text{Policy C})\}
	=\max\{100,70,60\}
	=100.
	\]
	The \GTenToneHL{best}{Maximax} choice is Policy A because it has the highest possible \GTenToneHL{payoff}{payoff}.
	
	\[
	\begin{aligned}
		\min(\text{Policy A}) &= \min\{100,20,0\}=0,\\
		\min(\text{Policy B}) &= \min\{70,60,50\}=50,\\
		\min(\text{Policy C}) &= \min\{60,40,30\}=30.
	\end{aligned}
	\]
	\[
	\max\{\min(\text{Policy A}),\min(\text{Policy B}),
	\min(\text{Policy C})\}
	=\max\{0,50,30\}
	=50.
	\]
	The \GTenToneHL{worst}{Maximin} choice is Policy B because it has the largest worst-case \GTenToneHL{payoff}{payoff}.
	
	\subsubsection*{C5 - Constructs opportunity loss tables and applies the Minimax regret criterion}
	\phantomsection
	\label{subsec:inventory-policy-3x3-c5}
	
	Best \GTenToneHL{payoff}{payoff} in each state:
	\[
	\begin{aligned}
		\text{High demand: }&\max\{100,70,60\}=100,\\
		\text{Moderate demand: }&\max\{20,60,40\}=60,\\
		\text{Low demand: }&\max\{0,50,30\}=50.
	\end{aligned}
	\]
	
	Regret table (best \GTenToneHL{payoff}{payoff} \( - \) \GTenToneHL{payoff}{payoff}):
	\begin{center}
		\begin{tabular}{l c c c c}
			\toprule
			Alternative & High & Moderate & Low & \GTenToneHL{maxregret}{Maximum} \GTenToneHL{maxregret}{regret} \\
			\midrule
			Policy A
			& \(100-100=0\)
			& \(60-20=40\)
			& \(50-0=50\)
			& \GTenToneHL{maxregret}{50} \\
			Policy B
			& \(100-70=30\)
			& \(60-60=0\)
			& \(50-50=0\)
			& \GTenToneHL{maxregret}{30} \\
			Policy C
			& \(100-60=40\)
			& \(60-40=20\)
			& \(50-30=20\)
			& \GTenToneHL{maxregret}{40} \\
			\bottomrule
		\end{tabular}
	\end{center}
	
	The \GTenToneHL{maxregret}{Minimax} \GTenToneHL{maxregret}{Regret} choice is Policy B because it has the smallest maximum
	\GTenToneHL{regret}{regret}, \(30\).
	
	\subsubsection*{Decisions summary}
	
	\begin{center}
		\begin{tabular}{l c}
			\toprule
			Criterion & Final decision \\
			\midrule
			\GTenToneHL{best}{Maximax}        & Policy A \\
			\GTenToneHL{worst}{Maximin}        & Policy B \\
			\GTenToneHL{maxregret}{Minimax} \GTenToneHL{maxregret}{Regret} & Policy B \\
			\bottomrule
		\end{tabular}
	\end{center}
	
	\secInicio[sec:inventory-policy-3x3][sec:problems-3x3]
	
	
	% ----------------------------------------------------------------
	\subsection{Event Venue}
	\label{sec:event-venue-3x3}
	
	An event organiser must choose Venue A, Venue B, or Venue C. Attendance
	can be high, moderate, or low. The final \GTenToneHL{payoff}{payoff} table, in thousands of
	dollars, is given below. Use \GTenToneHL{best}{Maximax}, \GTenToneHL{worst}{Maximin}, and \GTenToneHL{maxregret}{Minimax} \GTenToneHL{maxregret}{Regret} to select
	an event venue.
	
	\begin{center}
		\textit{Payoff table} \\
		\begin{tabular}{l c c c}
			\toprule
			State of nature & Venue A & Venue B & Venue C \\
			\midrule
			High attendance     & \GTenToneHL{payoff}{100} & \GTenToneHL{payoff}{70} & \GTenToneHL{payoff}{80} \\
			Moderate attendance & \GTenToneHL{payoff}{0}   & \GTenToneHL{payoff}{60} & \GTenToneHL{payoff}{50} \\
			Low attendance      & \GTenToneHL{payoff}{20}  & \GTenToneHL{payoff}{55} & \GTenToneHL{payoff}{50} \\
			\bottomrule
		\end{tabular}
	\end{center}
	
	\textbf{Questions/Tasks.}
	\begin{enumerate}
		\CriterionTask
		{subsec:event-venue-3x3-c4}{4}
		{Applies the Maximin and Maximax criteria}

		\CriterionTask
		{subsec:event-venue-3x3-c5}{5}
		{Constructs opportunity loss tables and applies the Minimax regret criterion}
	\end{enumerate}
	
	\subsubsection*{C4 - Applies the Maximin and Maximax criteria}
	\phantomsection
	\label{subsec:event-venue-3x3-c4}
	
	\[
	\begin{aligned}
		\max(\text{Venue A}) &= \max\{100,0,20\}=100,\\
		\max(\text{Venue B}) &= \max\{70,60,55\}=70,\\
		\max(\text{Venue C}) &= \max\{80,50,50\}=80.
	\end{aligned}
	\]
	\[
	\max\{\max(\text{Venue A}),\max(\text{Venue B}),\max(\text{Venue C})\}
	=\max\{100,70,80\}
	=100.
	\]
	The \GTenToneHL{best}{Maximax} choice is Venue A because it has the highest possible \GTenToneHL{payoff}{payoff}.
	
	\[
	\begin{aligned}
		\min(\text{Venue A}) &= \min\{100,0,20\}=0,\\
		\min(\text{Venue B}) &= \min\{70,60,55\}=55,\\
		\min(\text{Venue C}) &= \min\{80,50,50\}=50.
	\end{aligned}
	\]
	\[
	\max\{\min(\text{Venue A}),\min(\text{Venue B}),\min(\text{Venue C})\}
	=\max\{0,55,50\}
	=55.
	\]
	The \GTenToneHL{worst}{Maximin} choice is Venue B because it has the largest worst-case \GTenToneHL{payoff}{payoff}.
	
	\subsubsection*{C5 - Constructs opportunity loss tables and applies the Minimax regret criterion}
	\phantomsection
	\label{subsec:event-venue-3x3-c5}
	
	Best \GTenToneHL{payoff}{payoff} in each state:
	\[
	\begin{aligned}
		\text{High attendance: }&\max\{100,70,80\}=100,\\
		\text{Moderate attendance: }&\max\{0,60,50\}=60,\\
		\text{Low attendance: }&\max\{20,55,50\}=55.
	\end{aligned}
	\]
	
	Regret table (best \GTenToneHL{payoff}{payoff} \( - \) \GTenToneHL{payoff}{payoff}):
	\begin{center}
		\begin{tabular}{l c c c c}
			\toprule
			Alternative & High & Moderate & Low & \GTenToneHL{maxregret}{Maximum} \GTenToneHL{maxregret}{regret} \\
			\midrule
			Venue A
			& \(100-100=0\)
			& \(60-0=60\)
			& \(55-20=35\)
			& \GTenToneHL{maxregret}{60} \\
			Venue B
			& \(100-70=30\)
			& \(60-60=0\)
			& \(55-55=0\)
			& \GTenToneHL{maxregret}{30} \\
			Venue C
			& \(100-80=20\)
			& \(60-50=10\)
			& \(55-50=5\)
			& \GTenToneHL{maxregret}{20} \\
			\bottomrule
		\end{tabular}
	\end{center}
	
	The \GTenToneHL{maxregret}{Minimax} \GTenToneHL{maxregret}{Regret} choice is Venue C because it has the smallest maximum
	\GTenToneHL{regret}{regret}, \(20\).
	
	\subsubsection*{Decisions summary}
	
	\begin{center}
		\begin{tabular}{l c}
			\toprule
			Criterion & Final decision \\
			\midrule
			\GTenToneHL{best}{Maximax}        & Venue A \\
			\GTenToneHL{worst}{Maximin}        & Venue B \\
			\GTenToneHL{maxregret}{Minimax} \GTenToneHL{maxregret}{Regret} & Venue C \\
			\bottomrule
		\end{tabular}
	\end{center}
	
	\secInicio[sec:event-venue-3x3][sec:problems-3x3]
	
	
	% ----------------------------------------------------------------
	\subsection{Distribution Plan}
	\label{sec:distribution-plan-3x3}
	
	A company must choose Distribution Plan A, Plan B, or Plan C. Market
	conditions can be favourable, stable, or unfavourable. The final \GTenToneHL{payoff}{payoff}
	table, in thousands of dollars, is given below. Use \GTenToneHL{best}{Maximax}, \GTenToneHL{worst}{Maximin}, and
	\GTenToneHL{maxregret}{Minimax} \GTenToneHL{maxregret}{Regret} to select a distribution plan.
	
	\begin{center}
		\textit{Payoff table} \\
		\begin{tabular}{l c c c}
			\toprule
			State of nature & Plan A & Plan B & Plan C \\
			\midrule
			Favourable market   & \GTenToneHL{payoff}{90} & \GTenToneHL{payoff}{80} & \GTenToneHL{payoff}{70} \\
			Stable market       & \GTenToneHL{payoff}{70} & \GTenToneHL{payoff}{50} & \GTenToneHL{payoff}{60} \\
			Unfavourable market & \GTenToneHL{payoff}{60} & \GTenToneHL{payoff}{40} & \GTenToneHL{payoff}{30} \\
			\bottomrule
		\end{tabular}
	\end{center}
	
	\textbf{Questions/Tasks.}
	\begin{enumerate}
		\CriterionTask
		{subsec:distribution-plan-3x3-c4}{4}
		{Applies the Maximin and Maximax criteria}

		\CriterionTask
		{subsec:distribution-plan-3x3-c5}{5}
		{Constructs opportunity loss tables and applies the Minimax regret criterion}
	\end{enumerate}
	
	\subsubsection*{C4 - Applies the Maximin and Maximax criteria}
	\phantomsection
	\label{subsec:distribution-plan-3x3-c4}
	
	\[
	\begin{aligned}
		\max(\text{Plan A}) &= \max\{90,70,60\}=90,\\
		\max(\text{Plan B}) &= \max\{80,50,40\}=80,\\
		\max(\text{Plan C}) &= \max\{70,60,30\}=70.
	\end{aligned}
	\]
	\[
	\max\{\max(\text{Plan A}),\max(\text{Plan B}),\max(\text{Plan C})\}
	=\max\{90,80,70\}
	=90.
	\]
	The \GTenToneHL{best}{Maximax} choice is Plan A because it has the highest possible \GTenToneHL{payoff}{payoff}.
	
	\[
	\begin{aligned}
		\min(\text{Plan A}) &= \min\{90,70,60\}=60,\\
		\min(\text{Plan B}) &= \min\{80,50,40\}=40,\\
		\min(\text{Plan C}) &= \min\{70,60,30\}=30.
	\end{aligned}
	\]
	\[
	\max\{\min(\text{Plan A}),\min(\text{Plan B}),\min(\text{Plan C})\}
	=\max\{60,40,30\}
	=60.
	\]
	The \GTenToneHL{worst}{Maximin} choice is Plan A because it has the largest worst-case \GTenToneHL{payoff}{payoff}.
	
	\subsubsection*{C5 - Constructs opportunity loss tables and applies the Minimax regret criterion}
	\phantomsection
	\label{subsec:distribution-plan-3x3-c5}
	
	Best \GTenToneHL{payoff}{payoff} in each state:
	\[
	\begin{aligned}
		\text{Favourable market: }&\max\{90,80,70\}=90,\\
		\text{Stable market: }&\max\{70,50,60\}=70,\\
		\text{Unfavourable market: }&\max\{60,40,30\}=60.
	\end{aligned}
	\]
	
	Regret table (best \GTenToneHL{payoff}{payoff} \( - \) \GTenToneHL{payoff}{payoff}):
	\begin{center}
		\begin{tabular}{l c c c c}
			\toprule
			Alternative & Favourable & Stable & Unfavourable & \GTenToneHL{maxregret}{Maximum} \GTenToneHL{maxregret}{regret} \\
			\midrule
			Plan A
			& \(90-90=0\)
			& \(70-70=0\)
			& \(60-60=0\)
			& \GTenToneHL{maxregret}{0} \\
			Plan B
			& \(90-80=10\)
			& \(70-50=20\)
			& \(60-40=20\)
			& \GTenToneHL{maxregret}{20} \\
			Plan C
			& \(90-70=20\)
			& \(70-60=10\)
			& \(60-30=30\)
			& \GTenToneHL{maxregret}{30} \\
			\bottomrule
		\end{tabular}
	\end{center}
	
	The \GTenToneHL{maxregret}{Minimax} \GTenToneHL{maxregret}{Regret} choice is Plan A because it has the smallest maximum
	\GTenToneHL{regret}{regret}, \(0\).
	
	\subsubsection*{Decisions summary}
	
	\begin{center}
		\begin{tabular}{l c}
			\toprule
			Criterion & Final decision \\
			\midrule
			\GTenToneHL{best}{Maximax}        & Plan A \\
			\GTenToneHL{worst}{Maximin}        & Plan A \\
			\GTenToneHL{maxregret}{Minimax} \GTenToneHL{maxregret}{Regret} & Plan A \\
			\bottomrule
		\end{tabular}
	\end{center}
	
	\secInicio[sec:distribution-plan-3x3][sec:problems-3x3]
	
	
	% ================================================================
	\section{Problems \(4\times 2\)}
	\label{sec:problems-4x2}
	
	\SectionProblemsTableOfContents{
		\SectionProblemLink{sec:production-system-4x2}{Production System (AAB)} \\
		\SectionProblemLink{sec:service-network-4x2}{Service Network (ABA)} \\
		\SectionProblemLink{sec:inventory-strategy-4x2}{Inventory Strategy (ABB)} \\
		\SectionProblemLink{sec:supply-contract-4x2}{Supply Contract (AAA)}
	}
	
	
	% ----------------------------------------------------------------
	\subsection{Production System}
	\label{sec:production-system-4x2}
	
	A manufacturing company must choose between Production System A and
	Production System B. Four \GTenToneHL{state}{states} \GTenToneHL{state}{of} \GTenToneHL{state}{nature} represent different operating
	conditions. The final \GTenToneHL{payoff}{payoff} table, in thousands of dollars, is given below.
	Use \GTenToneHL{best}{Maximax}, \GTenToneHL{worst}{Maximin}, and \GTenToneHL{maxregret}{Minimax} \GTenToneHL{maxregret}{Regret} to select a production system.
	
	\begin{center}
		\textit{Payoff table} \\
		\begin{tabular}{l c c}
			\toprule
			State of nature & System A & System B \\
			\midrule
			State \(\GTenToneHL{state}{S_1}\) & \GTenToneHL{payoff}{100} & \GTenToneHL{payoff}{70} \\
			State \(\GTenToneHL{state}{S_2}\) & \GTenToneHL{payoff}{50}  & \GTenToneHL{payoff}{49} \\
			State \(\GTenToneHL{state}{S_3}\) & \GTenToneHL{payoff}{50}  & \GTenToneHL{payoff}{90} \\
			State \(\GTenToneHL{state}{S_4}\) & \GTenToneHL{payoff}{55}  & \GTenToneHL{payoff}{52} \\
			\bottomrule
		\end{tabular}
	\end{center}
	
	\textbf{Questions/Tasks.}
	\begin{enumerate}
		\CriterionTask
		{subsec:production-system-4x2-c4}{4}
		{Applies the Maximin and Maximax criteria}

		\CriterionTask
		{subsec:production-system-4x2-c5}{5}
		{Constructs opportunity loss tables and applies the Minimax regret criterion}
	\end{enumerate}
	
	\subsubsection*{C4 - Applies the Maximin and Maximax criteria}
	\phantomsection
	\label{subsec:production-system-4x2-c4}
	
	\[
	\begin{aligned}
		\max(\text{System A}) &= \max\{100,50,50,55\}=100,\\
		\max(\text{System B}) &= \max\{70,49,90,52\}=90.
	\end{aligned}
	\]
	\[
	\max\{\max(\text{System A}),\max(\text{System B})\}
	=\max\{100,90\}
	=100.
	\]
	The \GTenToneHL{best}{Maximax} choice is System A because it has the highest possible \GTenToneHL{payoff}{payoff}.
	
	\[
	\begin{aligned}
		\min(\text{System A}) &= \min\{100,50,50,55\}=50,\\
		\min(\text{System B}) &= \min\{70,49,90,52\}=49.
	\end{aligned}
	\]
	\[
	\max\{\min(\text{System A}),\min(\text{System B})\}
	=\max\{50,49\}
	=50.
	\]
	The \GTenToneHL{worst}{Maximin} choice is System A because it has the largest worst-case \GTenToneHL{payoff}{payoff}.
	
	\subsubsection*{C5 - Constructs opportunity loss tables and applies the Minimax regret criterion}
	\phantomsection
	\label{subsec:production-system-4x2-c5}
	
	Best \GTenToneHL{payoff}{payoff} in each state:
	\[
	\begin{aligned}
		\text{State }\GTenToneHL{state}{S_1}:\;&\max\{100,70\}=100,\\
		\text{State }\GTenToneHL{state}{S_2}:\;&\max\{50,49\}=50,\\
		\text{State }\GTenToneHL{state}{S_3}:\;&\max\{50,90\}=90,\\
		\text{State }\GTenToneHL{state}{S_4}:\;&\max\{55,52\}=55.
	\end{aligned}
	\]
	
	Regret table (best \GTenToneHL{payoff}{payoff} \( - \) \GTenToneHL{payoff}{payoff}):
	\begin{center}
		\begin{tabular}{l c c c c c}
			\toprule
			Alternative & \(\GTenToneHL{state}{S_1}\) & \(\GTenToneHL{state}{S_2}\) & \(\GTenToneHL{state}{S_3}\) & \(\GTenToneHL{state}{S_4}\)
			& \GTenToneHL{maxregret}{Maximum} \GTenToneHL{maxregret}{regret} \\
			\midrule
			System A
			& \(100-100=0\)
			& \(50-50=0\)
			& \(90-50=40\)
			& \(55-55=0\)
			& \GTenToneHL{maxregret}{40} \\
			System B
			& \(100-70=30\)
			& \(50-49=1\)
			& \(90-90=0\)
			& \(55-52=3\)
			& \GTenToneHL{maxregret}{30} \\
			\bottomrule
		\end{tabular}
	\end{center}
	
	The \GTenToneHL{maxregret}{Minimax} \GTenToneHL{maxregret}{Regret} choice is System B because it has the smallest maximum
	\GTenToneHL{regret}{regret}, \(30\).
	
	\subsubsection*{Decisions summary}
	
	\begin{center}
		\begin{tabular}{l c}
			\toprule
			Criterion & Final decision \\
			\midrule
			\GTenToneHL{best}{Maximax}        & System A \\
			\GTenToneHL{worst}{Maximin}        & System A \\
			\GTenToneHL{maxregret}{Minimax} \GTenToneHL{maxregret}{Regret} & System B \\
			\bottomrule
		\end{tabular}
	\end{center}
	
	\secInicio[sec:production-system-4x2][sec:problems-4x2]
	
	
	% ----------------------------------------------------------------
	\subsection{Service Network}
	\label{sec:service-network-4x2}
	
	A service company must choose between Service Network A and Service Network B.
	Four \GTenToneHL{state}{states} \GTenToneHL{state}{of} \GTenToneHL{state}{nature} represent different demand conditions. The final \GTenToneHL{payoff}{payoff}
	table, in thousands of dollars, is given below. Use \GTenToneHL{best}{Maximax}, \GTenToneHL{worst}{Maximin}, and
	\GTenToneHL{maxregret}{Minimax} \GTenToneHL{maxregret}{Regret} to select a service network.
	
	\begin{center}
		\textit{Payoff table} \\
		\begin{tabular}{l c c}
			\toprule
			State of nature & Network A & Network B \\
			\midrule
			State \(\GTenToneHL{state}{S_1}\) & \GTenToneHL{payoff}{100} & \GTenToneHL{payoff}{80} \\
			State \(\GTenToneHL{state}{S_2}\) & \GTenToneHL{payoff}{45}  & \GTenToneHL{payoff}{50} \\
			State \(\GTenToneHL{state}{S_3}\) & \GTenToneHL{payoff}{40}  & \GTenToneHL{payoff}{50} \\
			State \(\GTenToneHL{state}{S_4}\) & \GTenToneHL{payoff}{48}  & \GTenToneHL{payoff}{52} \\
			\bottomrule
		\end{tabular}
	\end{center}
	
	\textbf{Questions/Tasks.}
	\begin{enumerate}
		\CriterionTask
		{subsec:service-network-4x2-c4}{4}
		{Applies the Maximin and Maximax criteria}

		\CriterionTask
		{subsec:service-network-4x2-c5}{5}
		{Constructs opportunity loss tables and applies the Minimax regret criterion}
	\end{enumerate}
	
	\subsubsection*{C4 - Applies the Maximin and Maximax criteria}
	\phantomsection
	\label{subsec:service-network-4x2-c4}
	
	\[
	\begin{aligned}
		\max(\text{Network A}) &= \max\{100,45,40,48\}=100,\\
		\max(\text{Network B}) &= \max\{80,50,50,52\}=80.
	\end{aligned}
	\]
	\[
	\max\{\max(\text{Network A}),\max(\text{Network B})\}
	=\max\{100,80\}
	=100.
	\]
	The \GTenToneHL{best}{Maximax} choice is Network A because it has the highest possible \GTenToneHL{payoff}{payoff}.
	
	\[
	\begin{aligned}
		\min(\text{Network A}) &= \min\{100,45,40,48\}=40,\\
		\min(\text{Network B}) &= \min\{80,50,50,52\}=50.
	\end{aligned}
	\]
	\[
	\max\{\min(\text{Network A}),\min(\text{Network B})\}
	=\max\{40,50\}
	=50.
	\]
	The \GTenToneHL{worst}{Maximin} choice is Network B because it has the largest worst-case \GTenToneHL{payoff}{payoff}.
	
	\subsubsection*{C5 - Constructs opportunity loss tables and applies the Minimax regret criterion}
	\phantomsection
	\label{subsec:service-network-4x2-c5}
	
	Best \GTenToneHL{payoff}{payoff} in each state:
	\[
	\begin{aligned}
		\text{State }\GTenToneHL{state}{S_1}:\;&\max\{100,80\}=100,\\
		\text{State }\GTenToneHL{state}{S_2}:\;&\max\{45,50\}=50,\\
		\text{State }\GTenToneHL{state}{S_3}:\;&\max\{40,50\}=50,\\
		\text{State }\GTenToneHL{state}{S_4}:\;&\max\{48,52\}=52.
	\end{aligned}
	\]
	
	Regret table (best \GTenToneHL{payoff}{payoff} \( - \) \GTenToneHL{payoff}{payoff}):
	\begin{center}
		\begin{tabular}{l c c c c c}
			\toprule
			Alternative & \(\GTenToneHL{state}{S_1}\) & \(\GTenToneHL{state}{S_2}\) & \(\GTenToneHL{state}{S_3}\) & \(\GTenToneHL{state}{S_4}\)
			& \GTenToneHL{maxregret}{Maximum} \GTenToneHL{maxregret}{regret} \\
			\midrule
			Network A
			& \(100-100=0\)
			& \(50-45=5\)
			& \(50-40=10\)
			& \(52-48=4\)
			& \GTenToneHL{maxregret}{10} \\
			Network B
			& \(100-80=20\)
			& \(50-50=0\)
			& \(50-50=0\)
			& \(52-52=0\)
			& \GTenToneHL{maxregret}{20} \\
			\bottomrule
		\end{tabular}
	\end{center}
	
	The \GTenToneHL{maxregret}{Minimax} \GTenToneHL{maxregret}{Regret} choice is Network A because it has the smallest maximum
	\GTenToneHL{regret}{regret}, \(10\).
	
	\subsubsection*{Decisions summary}
	
	\begin{center}
		\begin{tabular}{l c}
			\toprule
			Criterion & Final decision \\
			\midrule
			\GTenToneHL{best}{Maximax}        & Network A \\
			\GTenToneHL{worst}{Maximin}        & Network B \\
			\GTenToneHL{maxregret}{Minimax} \GTenToneHL{maxregret}{Regret} & Network A \\
			\bottomrule
		\end{tabular}
	\end{center}
	
	\secInicio[sec:service-network-4x2][sec:problems-4x2]
	
	
	% ----------------------------------------------------------------
	\subsection{Inventory Strategy}
	\label{sec:inventory-strategy-4x2}
	
	A retailer must choose between Inventory Strategy A and Inventory Strategy B.
	Four \GTenToneHL{state}{states} \GTenToneHL{state}{of} \GTenToneHL{state}{nature} represent different demand conditions. The final \GTenToneHL{payoff}{payoff}
	table, in thousands of dollars, is given below. Use \GTenToneHL{best}{Maximax}, \GTenToneHL{worst}{Maximin}, and
	\GTenToneHL{maxregret}{Minimax} \GTenToneHL{maxregret}{Regret} to select an inventory strategy.
	
	\begin{center}
		\textit{Payoff table} \\
		\begin{tabular}{l c c}
			\toprule
			State of nature & Strategy A & Strategy B \\
			\midrule
			State \(\GTenToneHL{state}{S_1}\) & \GTenToneHL{rate}{100} & \GTenToneHL{rate}{70} \\
			State \(\GTenToneHL{state}{S_2}\) & \GTenToneHL{rate}{20}  & \GTenToneHL{rate}{60} \\
			State \(\GTenToneHL{state}{S_3}\) & \GTenToneHL{rate}{0}   & \GTenToneHL{rate}{50} \\
			State \(\GTenToneHL{state}{S_4}\) & \GTenToneHL{rate}{25}  & \GTenToneHL{rate}{55} \\
			\bottomrule
		\end{tabular}
	\end{center}
	
	\textbf{Questions/Tasks.}
	\begin{enumerate}
		\CriterionTask
		{subsec:inventory-strategy-4x2-c4}{4}
		{Applies the Maximin and Maximax criteria}

		\CriterionTask
		{subsec:inventory-strategy-4x2-c5}{5}
		{Constructs opportunity loss tables and applies the Minimax regret criterion}
	\end{enumerate}
	
	\subsubsection*{C4 - Applies the Maximin and Maximax criteria}
	\phantomsection
	\label{subsec:inventory-strategy-4x2-c4}
	
	\[
	\begin{aligned}
		\max(\text{Strategy A}) &= \max\{100,20,0,25\}=100,\\
		\max(\text{Strategy B}) &= \max\{70,60,50,55\}=70.
	\end{aligned}
	\]
	\[
	\max\{\max(\text{Strategy A}),\max(\text{Strategy B})\}
	=\max\{100,70\}
	=100.
	\]
	The \GTenToneHL{best}{Maximax} choice is Strategy A because it has the highest possible \GTenToneHL{payoff}{payoff}.
	
	\[
	\begin{aligned}
		\min(\text{Strategy A}) &= \min\{100,20,0,25\}=0,\\
		\min(\text{Strategy B}) &= \min\{70,60,50,55\}=50.
	\end{aligned}
	\]
	\[
	\max\{\min(\text{Strategy A}),\min(\text{Strategy B})\}
	=\max\{0,50\}
	=50.
	\]
	The \GTenToneHL{worst}{Maximin} choice is Strategy B because it has the largest worst-case \GTenToneHL{payoff}{payoff}.
	
	\subsubsection*{C5 - Constructs opportunity loss tables and applies the Minimax regret criterion}
	\phantomsection
	\label{subsec:inventory-strategy-4x2-c5}
	
	Best \GTenToneHL{payoff}{payoff} in each state:
	\[
	\begin{aligned}
		\text{State }\GTenToneHL{state}{S_1}:\;&\max\{100,70\}=100,\\
		\text{State }\GTenToneHL{state}{S_2}:\;&\max\{20,60\}=60,\\
		\text{State }\GTenToneHL{state}{S_3}:\;&\max\{0,50\}=50,\\
		\text{State }\GTenToneHL{state}{S_4}:\;&\max\{25,55\}=55.
	\end{aligned}
	\]
	
	Regret table (best \GTenToneHL{payoff}{payoff} \( - \) \GTenToneHL{payoff}{payoff}):
	\begin{center}
		\begin{tabular}{l c c c c c}
			\toprule
			Alternative & \(\GTenToneHL{state}{S_1}\) & \(\GTenToneHL{state}{S_2}\) & \(\GTenToneHL{state}{S_3}\) & \(\GTenToneHL{state}{S_4}\)
			& \GTenToneHL{maxregret}{Maximum} \GTenToneHL{maxregret}{regret} \\
			\midrule
			Strategy A
			& \(100-100=0\)
			& \(60-20=40\)
			& \(50-0=50\)
			& \(55-25=30\)
			& \GTenToneHL{maxregret}{50} \\
			Strategy B
			& \(100-70=30\)
			& \(60-60=0\)
			& \(50-50=0\)
			& \(55-55=0\)
			& \GTenToneHL{maxregret}{30} \\
			\bottomrule
		\end{tabular}
	\end{center}
	
	The \GTenToneHL{maxregret}{Minimax} \GTenToneHL{maxregret}{Regret} choice is Strategy B because it has the smallest maximum
	\GTenToneHL{regret}{regret}, \(30\).
	
	\subsubsection*{Decisions summary}
	
	\begin{center}
		\begin{tabular}{l c}
			\toprule
			Criterion & Final decision \\
			\midrule
			\GTenToneHL{best}{Maximax}        & Strategy A \\
			\GTenToneHL{worst}{Maximin}        & Strategy B \\
			\GTenToneHL{maxregret}{Minimax} \GTenToneHL{maxregret}{Regret} & Strategy B \\
			\bottomrule
		\end{tabular}
	\end{center}
	
	\secInicio[sec:inventory-strategy-4x2][sec:problems-4x2]
	
	
	% ----------------------------------------------------------------
	\subsection{Supply Contract}
	\label{sec:supply-contract-4x2}
	
	A company must choose between Supply Contract A and Supply Contract B.
	Four \GTenToneHL{state}{states} \GTenToneHL{state}{of} \GTenToneHL{state}{nature} represent different market and supply conditions.
	The final \GTenToneHL{payoff}{payoff} table, in thousands of dollars, is given below. Use \GTenToneHL{best}{Maximax},
	\GTenToneHL{worst}{Maximin}, and \GTenToneHL{maxregret}{Minimax} \GTenToneHL{maxregret}{Regret} to select a supply contract.
	
	\begin{center}
		\textit{Payoff table} \\
		\begin{tabular}{l c c}
			\toprule
			State of nature & Contract A & Contract B \\
			\midrule
			State \(\GTenToneHL{state}{S_1}\) & \GTenToneHL{payoff}{90} & \GTenToneHL{payoff}{80} \\
			State \(\GTenToneHL{state}{S_2}\) & \GTenToneHL{payoff}{70} & \GTenToneHL{payoff}{50} \\
			State \(\GTenToneHL{state}{S_3}\) & \GTenToneHL{payoff}{60} & \GTenToneHL{payoff}{40} \\
			State \(\GTenToneHL{state}{S_4}\) & \GTenToneHL{payoff}{65} & \GTenToneHL{payoff}{55} \\
			\bottomrule
		\end{tabular}
	\end{center}
	
	\textbf{Questions/Tasks.}
	\begin{enumerate}
		\CriterionTask
		{subsec:supply-contract-4x2-c4}{4}
		{Applies the Maximin and Maximax criteria}

		\CriterionTask
		{subsec:supply-contract-4x2-c5}{5}
		{Constructs opportunity loss tables and applies the Minimax regret criterion}
	\end{enumerate}
	
	\subsubsection*{C4 - Applies the Maximin and Maximax criteria}
	\phantomsection
	\label{subsec:supply-contract-4x2-c4}
	
	\[
	\begin{aligned}
		\max(\text{Contract A}) &= \max\{90,70,60,65\}=90,\\
		\max(\text{Contract B}) &= \max\{80,50,40,55\}=80.
	\end{aligned}
	\]
	\[
	\max\{\max(\text{Contract A}),\max(\text{Contract B})\}
	=\max\{90,80\}
	=90.
	\]
	The \GTenToneHL{best}{Maximax} choice is Contract A because it has the highest possible \GTenToneHL{payoff}{payoff}.
	
	\[
	\begin{aligned}
		\min(\text{Contract A}) &= \min\{90,70,60,65\}=60,\\
		\min(\text{Contract B}) &= \min\{80,50,40,55\}=40.
	\end{aligned}
	\]
	\[
	\max\{\min(\text{Contract A}),\min(\text{Contract B})\}
	=\max\{60,40\}
	=60.
	\]
	The \GTenToneHL{worst}{Maximin} choice is Contract A because it has the largest worst-case \GTenToneHL{payoff}{payoff}.
	
	\subsubsection*{C5 - Constructs opportunity loss tables and applies the Minimax regret criterion}
	\phantomsection
	\label{subsec:supply-contract-4x2-c5}
	
	Best \GTenToneHL{payoff}{payoff} in each state:
	\[
	\begin{aligned}
		\text{State }\GTenToneHL{state}{S_1}:\;&\max\{90,80\}=90,\\
		\text{State }\GTenToneHL{state}{S_2}:\;&\max\{70,50\}=70,\\
		\text{State }\GTenToneHL{state}{S_3}:\;&\max\{60,40\}=60,\\
		\text{State }\GTenToneHL{state}{S_4}:\;&\max\{65,55\}=65.
	\end{aligned}
	\]
	
	Regret table (best \GTenToneHL{payoff}{payoff} \( - \) \GTenToneHL{payoff}{payoff}):
	\begin{center}
		\begin{tabular}{l c c c c c}
			\toprule
			Alternative & \(\GTenToneHL{state}{S_1}\) & \(\GTenToneHL{state}{S_2}\) & \(\GTenToneHL{state}{S_3}\) & \(\GTenToneHL{state}{S_4}\)
			& \GTenToneHL{maxregret}{Maximum} \GTenToneHL{maxregret}{regret} \\
			\midrule
			Contract A
			& \(90-90=0\)
			& \(70-70=0\)
			& \(60-60=0\)
			& \(65-65=0\)
			& \GTenToneHL{maxregret}{0} \\
			Contract B
			& \(90-80=10\)
			& \(70-50=20\)
			& \(60-40=20\)
			& \(65-55=10\)
			& \GTenToneHL{maxregret}{20} \\
			\bottomrule
		\end{tabular}
	\end{center}
	
	The \GTenToneHL{maxregret}{Minimax} \GTenToneHL{maxregret}{Regret} choice is Contract A because it has the smallest maximum
	\GTenToneHL{regret}{regret}, \(0\).
	
	\subsubsection*{Decisions summary}
	
	\begin{center}
		\begin{tabular}{l c}
			\toprule
			Criterion & Final decision \\
			\midrule
			\GTenToneHL{best}{Maximax}        & Contract A \\
			\GTenToneHL{worst}{Maximin}        & Contract A \\
			\GTenToneHL{maxregret}{Minimax} \GTenToneHL{maxregret}{Regret} & Contract A \\
			\bottomrule
		\end{tabular}
	\end{center}
	
	\secInicio[sec:supply-contract-4x2][sec:problems-4x2]
	
	
	% ================================================================
	\section{Problems \(2\times 4\)}
	\label{sec:problems-2x4}
	
	\SectionProblemsTableOfContents{
		\SectionProblemLink{sec:market-entry-strategy-2x4}{Market Entry Strategy (ABA)} \\
		\SectionProblemLink{sec:production-plan-2x4}{Production Plan (ABB)} \\
		\SectionProblemLink{sec:facility-location-2x4}{Facility Location (ABC)} \\
		\SectionProblemLink{sec:supply-agreement-2x4}{Supply Agreement (AAA)}
	}
	
	% ----------------------------------------------------------------
	\subsection{Market Entry Strategy}
	\label{sec:market-entry-strategy-2x4}
	
	A company must choose among Market Entry Strategy A, Strategy B, Strategy C,
	and Strategy D. Market conditions can be favourable or unfavourable. The final
	\GTenToneHL{payoff}{payoff} table, in thousands of dollars, is given below. Use \GTenToneHL{best}{Maximax}, \GTenToneHL{worst}{Maximin},
	and \GTenToneHL{maxregret}{Minimax} \GTenToneHL{maxregret}{Regret} to select a market entry strategy.
	
	\begin{center}
		\textit{Payoff table} \\
		\begin{tabular}{l c c c c}
			\toprule
			State of nature
			& Strategy A & Strategy B & Strategy C & Strategy D \\
			\midrule
			Favourable market   & \GTenToneHL{rate}{30}  & \GTenToneHL{rate}{60} & \GTenToneHL{rate}{50} & \GTenToneHL{rate}{0}  \\
			Unfavourable market & \GTenToneHL{rate}{100} & \GTenToneHL{rate}{40} & \GTenToneHL{rate}{10} & \GTenToneHL{rate}{30} \\
			\bottomrule
		\end{tabular}
	\end{center}
	
	\textbf{Questions/Tasks.}
	\begin{enumerate}
		\CriterionTask
		{subsec:market-entry-strategy-2x4-c4}{4}
		{Applies the Maximin and Maximax criteria}

		\CriterionTask
		{subsec:market-entry-strategy-2x4-c5}{5}
		{Constructs opportunity loss tables and applies the Minimax regret criterion}
	\end{enumerate}
	
	\subsubsection*{C4 - Applies the Maximin and Maximax criteria}
	\phantomsection
	\label{subsec:market-entry-strategy-2x4-c4}
	
	\[
	\begin{aligned}
		\max(\text{Strategy A}) &= \max\{30,100\}=100,\\
		\max(\text{Strategy B}) &= \max\{60,40\}=60,\\
		\max(\text{Strategy C}) &= \max\{50,10\}=50,\\
		\max(\text{Strategy D}) &= \max\{0,30\}=30.
	\end{aligned}
	\]
	\[
	\begin{aligned}
		&\max\{\max(\text{Strategy A}),\max(\text{Strategy B}),
		\max(\text{Strategy C}),\max(\text{Strategy D})\}\\
		&\qquad=\max\{100,60,50,30\}
		=100.
	\end{aligned}
	\]
	The \GTenToneHL{best}{Maximax} choice is Strategy A because it has the highest possible \GTenToneHL{payoff}{payoff}.
	
	\[
	\begin{aligned}
		\min(\text{Strategy A}) &= \min\{30,100\}=30,\\
		\min(\text{Strategy B}) &= \min\{60,40\}=40,\\
		\min(\text{Strategy C}) &= \min\{50,10\}=10,\\
		\min(\text{Strategy D}) &= \min\{0,30\}=0.
	\end{aligned}
	\]
	\[
	\begin{aligned}
		&\max\{\min(\text{Strategy A}),\min(\text{Strategy B}),
		\min(\text{Strategy C}),\min(\text{Strategy D})\}\\
		&\qquad=\max\{30,40,10,0\}
		=40.
	\end{aligned}
	\]
	The \GTenToneHL{worst}{Maximin} choice is Strategy B because it has the largest worst-case \GTenToneHL{payoff}{payoff}.
	
	\subsubsection*{C5 - Constructs opportunity loss tables and applies the Minimax regret criterion}
	\phantomsection
	\label{subsec:market-entry-strategy-2x4-c5}
	
	Best \GTenToneHL{payoff}{payoff} in each state:
	\[
	\begin{aligned}
		\text{Favourable market: }&
		\max\{30,60,50,0\}=60,\\
		\text{Unfavourable market: }&
		\max\{100,40,10,30\}=100.
	\end{aligned}
	\]
	
	Regret table (best \GTenToneHL{payoff}{payoff} \( - \) \GTenToneHL{payoff}{payoff}):
	\begin{center}
		\begin{tabular}{l c c c}
			\toprule
			Alternative & Favourable & Unfavourable & \GTenToneHL{maxregret}{Maximum} \GTenToneHL{maxregret}{regret} \\
			\midrule
			Strategy A
			& \(60-30=30\)
			& \(100-100=0\)
			& \GTenToneHL{maxregret}{30} \\
			Strategy B
			& \(60-60=0\)
			& \(100-40=60\)
			& \GTenToneHL{maxregret}{60} \\
			Strategy C
			& \(60-50=10\)
			& \(100-10=90\)
			& \GTenToneHL{maxregret}{90} \\
			Strategy D
			& \(60-0=60\)
			& \(100-30=70\)
			& \GTenToneHL{maxregret}{70} \\
			\bottomrule
		\end{tabular}
	\end{center}
	
	The \GTenToneHL{maxregret}{Minimax} \GTenToneHL{maxregret}{Regret} choice is Strategy A because it has the smallest maximum
	\GTenToneHL{regret}{regret}, \(30\).
	
	\subsubsection*{Decisions summary}
	
	\begin{center}
		\begin{tabular}{l c}
			\toprule
			Criterion & Final decision \\
			\midrule
			\GTenToneHL{best}{Maximax}        & Strategy A \\
			\GTenToneHL{worst}{Maximin}        & Strategy B \\
			\GTenToneHL{maxregret}{Minimax} \GTenToneHL{maxregret}{Regret} & Strategy A \\
			\bottomrule
		\end{tabular}
	\end{center}
	
	\secInicio[sec:market-entry-strategy-2x4][sec:problems-2x4]
	
	
	% ----------------------------------------------------------------
	\subsection{Production Plan}
	\label{sec:production-plan-2x4}
	
	A manufacturing company must choose among Production Plan A, Plan B, Plan C,
	and Plan D. Demand can be high or low. The final \GTenToneHL{payoff}{payoff} table, in thousands
	of dollars, is given below. Use \GTenToneHL{best}{Maximax}, \GTenToneHL{worst}{Maximin}, and \GTenToneHL{maxregret}{Minimax} \GTenToneHL{maxregret}{Regret} to select
	a production plan.
	
	\begin{center}
		\textit{Payoff table} \\
		\begin{tabular}{l c c c c}
			\toprule
			State of nature & Plan A & Plan B & Plan C & Plan D \\
			\midrule
			High demand & \GTenToneHL{payoff}{100} & \GTenToneHL{payoff}{90} & \GTenToneHL{payoff}{40} & \GTenToneHL{payoff}{30} \\
			Low demand  & \GTenToneHL{payoff}{70}  & \GTenToneHL{payoff}{90} & \GTenToneHL{payoff}{80} & \GTenToneHL{payoff}{40} \\
			\bottomrule
		\end{tabular}
	\end{center}
	
	\textbf{Questions/Tasks.}
	\begin{enumerate}
		\CriterionTask
		{subsec:production-plan-2x4-c4}{4}
		{Applies the Maximin and Maximax criteria}

		\CriterionTask
		{subsec:production-plan-2x4-c5}{5}
		{Constructs opportunity loss tables and applies the Minimax regret criterion}
	\end{enumerate}
	
	\subsubsection*{C4 - Applies the Maximin and Maximax criteria}
	\phantomsection
	\label{subsec:production-plan-2x4-c4}
	
	\[
	\begin{aligned}
		\max(\text{Plan A}) &= \max\{100,70\}=100,\\
		\max(\text{Plan B}) &= \max\{90,90\}=90,\\
		\max(\text{Plan C}) &= \max\{40,80\}=80,\\
		\max(\text{Plan D}) &= \max\{30,40\}=40.
	\end{aligned}
	\]
	\[
	\begin{aligned}
		&\max\{\max(\text{Plan A}),\max(\text{Plan B}),
		\max(\text{Plan C}),\max(\text{Plan D})\}\\
		&\qquad=\max\{100,90,80,40\}
		=100.
	\end{aligned}
	\]
	The \GTenToneHL{best}{Maximax} choice is Plan A because it has the highest possible \GTenToneHL{payoff}{payoff}.
	
	\[
	\begin{aligned}
		\min(\text{Plan A}) &= \min\{100,70\}=70,\\
		\min(\text{Plan B}) &= \min\{90,90\}=90,\\
		\min(\text{Plan C}) &= \min\{40,80\}=40,\\
		\min(\text{Plan D}) &= \min\{30,40\}=30.
	\end{aligned}
	\]
	\[
	\begin{aligned}
		&\max\{\min(\text{Plan A}),\min(\text{Plan B}),
		\min(\text{Plan C}),\min(\text{Plan D})\}\\
		&\qquad=\max\{70,90,40,30\}
		=90.
	\end{aligned}
	\]
	The \GTenToneHL{worst}{Maximin} choice is Plan B because it has the largest worst-case \GTenToneHL{payoff}{payoff}.
	
	\subsubsection*{C5 - Constructs opportunity loss tables and applies the Minimax regret criterion}
	\phantomsection
	\label{subsec:production-plan-2x4-c5}
	
	Best \GTenToneHL{payoff}{payoff} in each state:
	\[
	\begin{aligned}
		\text{High demand: }&\max\{100,90,40,30\}=100,\\
		\text{Low demand: }&\max\{70,90,80,40\}=90.
	\end{aligned}
	\]
	
	Regret table (best \GTenToneHL{payoff}{payoff} \( - \) \GTenToneHL{payoff}{payoff}):
	\begin{center}
		\begin{tabular}{l c c c}
			\toprule
			Alternative & High demand & Low demand & \GTenToneHL{maxregret}{Maximum} \GTenToneHL{maxregret}{regret} \\
			\midrule
			Plan A
			& \(100-100=0\)
			& \(90-70=20\)
			& \GTenToneHL{maxregret}{20} \\
			Plan B
			& \(100-90=10\)
			& \(90-90=0\)
			& \GTenToneHL{maxregret}{10} \\
			Plan C
			& \(100-40=60\)
			& \(90-80=10\)
			& \GTenToneHL{maxregret}{60} \\
			Plan D
			& \(100-30=70\)
			& \(90-40=50\)
			& \GTenToneHL{maxregret}{70} \\
			\bottomrule
		\end{tabular}
	\end{center}
	
	The \GTenToneHL{maxregret}{Minimax} \GTenToneHL{maxregret}{Regret} choice is Plan B because it has the smallest maximum
	\GTenToneHL{regret}{regret}, \(10\).
	
	\subsubsection*{Decisions summary}
	
	\begin{center}
		\begin{tabular}{l c}
			\toprule
			Criterion & Final decision \\
			\midrule
			\GTenToneHL{best}{Maximax}        & Plan A \\
			\GTenToneHL{worst}{Maximin}        & Plan B \\
			\GTenToneHL{maxregret}{Minimax} \GTenToneHL{maxregret}{Regret} & Plan B \\
			\bottomrule
		\end{tabular}
	\end{center}
	
	\secInicio[sec:production-plan-2x4][sec:problems-2x4]
	
	
	% ----------------------------------------------------------------
	\subsection{Facility Location}
	\label{sec:facility-location-2x4}
	
	A company must choose among Facility Location A, Location B, Location C, and
	Location D. Regional demand can be high or low. The final \GTenToneHL{payoff}{payoff} table, in
	thousands of dollars, is given below. Use \GTenToneHL{best}{Maximax}, \GTenToneHL{worst}{Maximin}, and \GTenToneHL{maxregret}{Minimax} \GTenToneHL{maxregret}{Regret}
	to select a facility location.
	
	\begin{center}
		\textit{Payoff table} \\
		\begin{tabular}{l c c c c}
			\toprule
			State of nature
			& Location A & Location B & Location C & Location D \\
			\midrule
			High regional demand & \GTenToneHL{payoff}{10}  & \GTenToneHL{payoff}{60} & \GTenToneHL{payoff}{50} & \GTenToneHL{payoff}{20} \\
			Low regional demand  & \GTenToneHL{payoff}{100} & \GTenToneHL{payoff}{70} & \GTenToneHL{payoff}{80} & \GTenToneHL{payoff}{40} \\
			\bottomrule
		\end{tabular}
	\end{center}
	
	\textbf{Questions/Tasks.}
	\begin{enumerate}
		\CriterionTask
		{subsec:facility-location-2x4-c4}{4}
		{Applies the Maximin and Maximax criteria}

		\CriterionTask
		{subsec:facility-location-2x4-c5}{5}
		{Constructs opportunity loss tables and applies the Minimax regret criterion}
	\end{enumerate}
	
	\subsubsection*{C4 - Applies the Maximin and Maximax criteria}
	\phantomsection
	\label{subsec:facility-location-2x4-c4}
	
	\[
	\begin{aligned}
		\max(\text{Location A}) &= \max\{10,100\}=100,\\
		\max(\text{Location B}) &= \max\{60,70\}=70,\\
		\max(\text{Location C}) &= \max\{50,80\}=80,\\
		\max(\text{Location D}) &= \max\{20,40\}=40.
	\end{aligned}
	\]
	\[
	\begin{aligned}
		&\max\{\max(\text{Location A}),\max(\text{Location B}),
		\max(\text{Location C}),\max(\text{Location D})\}\\
		&\qquad=\max\{100,70,80,40\}
		=100.
	\end{aligned}
	\]
	The \GTenToneHL{best}{Maximax} choice is Location A because it has the highest possible \GTenToneHL{payoff}{payoff}.
	
	\[
	\begin{aligned}
		\min(\text{Location A}) &= \min\{10,100\}=10,\\
		\min(\text{Location B}) &= \min\{60,70\}=60,\\
		\min(\text{Location C}) &= \min\{50,80\}=50,\\
		\min(\text{Location D}) &= \min\{20,40\}=20.
	\end{aligned}
	\]
	\[
	\begin{aligned}
		&\max\{\min(\text{Location A}),\min(\text{Location B}),
		\min(\text{Location C}),\min(\text{Location D})\}\\
		&\qquad=\max\{10,60,50,20\}
		=60.
	\end{aligned}
	\]
	The \GTenToneHL{worst}{Maximin} choice is Location B because it has the largest worst-case \GTenToneHL{payoff}{payoff}.
	
	\subsubsection*{C5 - Constructs opportunity loss tables and applies the Minimax regret criterion}
	\phantomsection
	\label{subsec:facility-location-2x4-c5}
	
	Best \GTenToneHL{payoff}{payoff} in each state:
	\[
	\begin{aligned}
		\text{High regional demand: }&
		\max\{10,60,50,20\}=60,\\
		\text{Low regional demand: }&
		\max\{100,70,80,40\}=100.
	\end{aligned}
	\]
	
	Regret table (best \GTenToneHL{payoff}{payoff} \( - \) \GTenToneHL{payoff}{payoff}):
	\begin{center}
		\begin{tabular}{l c c c}
			\toprule
			Alternative & High demand & Low demand & \GTenToneHL{maxregret}{Maximum} \GTenToneHL{maxregret}{regret} \\
			\midrule
			Location A
			& \(60-10=50\)
			& \(100-100=0\)
			& \GTenToneHL{maxregret}{50} \\
			Location B
			& \(60-60=0\)
			& \(100-70=30\)
			& \GTenToneHL{maxregret}{30} \\
			Location C
			& \(60-50=10\)
			& \(100-80=20\)
			& \GTenToneHL{maxregret}{20} \\
			Location D
			& \(60-20=40\)
			& \(100-40=60\)
			& \GTenToneHL{maxregret}{60} \\
			\bottomrule
		\end{tabular}
	\end{center}
	
	The \GTenToneHL{maxregret}{Minimax} \GTenToneHL{maxregret}{Regret} choice is Location C because it has the smallest maximum
	\GTenToneHL{regret}{regret}, \(20\).
	
	\subsubsection*{Decisions summary}
	
	\begin{center}
		\begin{tabular}{l c}
			\toprule
			Criterion & Final decision \\
			\midrule
			\GTenToneHL{best}{Maximax}        & Location A \\
			\GTenToneHL{worst}{Maximin}        & Location B \\
			\GTenToneHL{maxregret}{Minimax} \GTenToneHL{maxregret}{Regret} & Location C \\
			\bottomrule
		\end{tabular}
	\end{center}
	
	\secInicio[sec:facility-location-2x4][sec:problems-2x4]
	
	
	% ----------------------------------------------------------------
	\subsection{Supply Agreement}
	\label{sec:supply-agreement-2x4}
	
	A company must choose among Supply Agreement A, Agreement B, Agreement C,
	and Agreement D. Supply conditions can be stable or disrupted. The final
	\GTenToneHL{payoff}{payoff} table, in thousands of dollars, is given below. Use \GTenToneHL{best}{Maximax}, \GTenToneHL{worst}{Maximin},
	and \GTenToneHL{maxregret}{Minimax} \GTenToneHL{maxregret}{Regret} to select a supply agreement.
	
	\begin{center}
		\textit{Payoff table} \\
		\begin{tabular}{l c c c c}
			\toprule
			State of nature
			& Agreement A & Agreement B & Agreement C & Agreement D \\
			\midrule
			Stable supply    & \GTenToneHL{payoff}{90}  & \GTenToneHL{payoff}{70} & \GTenToneHL{payoff}{60} & \GTenToneHL{payoff}{30} \\
			Disrupted supply & \GTenToneHL{payoff}{100} & \GTenToneHL{payoff}{80} & \GTenToneHL{payoff}{50} & \GTenToneHL{payoff}{70} \\
			\bottomrule
		\end{tabular}
	\end{center}
	
	\textbf{Questions/Tasks.}
	\begin{enumerate}
		\CriterionTask
		{subsec:supply-agreement-2x4-c4}{4}
		{Applies the Maximin and Maximax criteria}

		\CriterionTask
		{subsec:supply-agreement-2x4-c5}{5}
		{Constructs opportunity loss tables and applies the Minimax regret criterion}
	\end{enumerate}
	
	\subsubsection*{C4 - Applies the Maximin and Maximax criteria}
	\phantomsection
	\label{subsec:supply-agreement-2x4-c4}
	
	\[
	\begin{aligned}
		\max(\text{Agreement A}) &= \max\{90,100\}=100,\\
		\max(\text{Agreement B}) &= \max\{70,80\}=80,\\
		\max(\text{Agreement C}) &= \max\{60,50\}=60,\\
		\max(\text{Agreement D}) &= \max\{30,70\}=70.
	\end{aligned}
	\]
	\[
	\begin{aligned}
		&\max\{\max(\text{Agreement A}),\max(\text{Agreement B}),
		\max(\text{Agreement C}),\max(\text{Agreement D})\}\\
		&\qquad=\max\{100,80,60,70\}
		=100.
	\end{aligned}
	\]
	The \GTenToneHL{best}{Maximax} choice is Agreement A because it has the highest possible \GTenToneHL{payoff}{payoff}.
	
	\[
	\begin{aligned}
		\min(\text{Agreement A}) &= \min\{90,100\}=90,\\
		\min(\text{Agreement B}) &= \min\{70,80\}=70,\\
		\min(\text{Agreement C}) &= \min\{60,50\}=50,\\
		\min(\text{Agreement D}) &= \min\{30,70\}=30.
	\end{aligned}
	\]
	\[
	\begin{aligned}
		&\max\{\min(\text{Agreement A}),\min(\text{Agreement B}),
		\min(\text{Agreement C}),\min(\text{Agreement D})\}\\
		&\qquad=\max\{90,70,50,30\}
		=90.
	\end{aligned}
	\]
	The \GTenToneHL{worst}{Maximin} choice is Agreement A because it has the largest worst-case \GTenToneHL{payoff}{payoff}.
	
	\subsubsection*{C5 - Constructs opportunity loss tables and applies the Minimax regret criterion}
	\phantomsection
	\label{subsec:supply-agreement-2x4-c5}
	
	Best \GTenToneHL{payoff}{payoff} in each state:
	\[
	\begin{aligned}
		\text{Stable supply: }&
		\max\{90,70,60,30\}=90,\\
		\text{Disrupted supply: }&
		\max\{100,80,50,70\}=100.
	\end{aligned}
	\]
	
	Regret table (best \GTenToneHL{payoff}{payoff} \( - \) \GTenToneHL{payoff}{payoff}):
	\begin{center}
		\begin{tabular}{l c c c}
			\toprule
			Alternative & Stable supply & Disrupted supply & \GTenToneHL{maxregret}{Maximum} \GTenToneHL{maxregret}{regret} \\
			\midrule
			Agreement A
			& \(90-90=0\)
			& \(100-100=0\)
			& \GTenToneHL{maxregret}{0} \\
			Agreement B
			& \(90-70=20\)
			& \(100-80=20\)
			& \GTenToneHL{maxregret}{20} \\
			Agreement C
			& \(90-60=30\)
			& \(100-50=50\)
			& \GTenToneHL{maxregret}{50} \\
			Agreement D
			& \(90-30=60\)
			& \(100-70=30\)
			& \GTenToneHL{maxregret}{60} \\
			\bottomrule
		\end{tabular}
	\end{center}
	
	The \GTenToneHL{maxregret}{Minimax} \GTenToneHL{maxregret}{Regret} choice is Agreement A because it has the smallest maximum
	\GTenToneHL{regret}{regret}, \(0\).
	
	\subsubsection*{Decisions summary}
	
	\begin{center}
		\begin{tabular}{l c}
			\toprule
			Criterion & Final decision \\
			\midrule
			\GTenToneHL{best}{Maximax}        & Agreement A \\
			\GTenToneHL{worst}{Maximin}        & Agreement A \\
			\GTenToneHL{maxregret}{Minimax} \GTenToneHL{maxregret}{Regret} & Agreement A \\
			\bottomrule
		\end{tabular}
	\end{center}
	
	\secInicio[sec:supply-agreement-2x4][sec:problems-2x4]
	
	
	% ================================================================
	\section{Problems \(4\times 3\)}
	\label{sec:problems-4x3}
	
	\SectionProblemsTableOfContents{
		\SectionProblemLink{sec:production-system-4x3}{Production System (AAB)} \\
		\SectionProblemLink{sec:service-network-4x3}{Service Network (ABA)} \\
		\SectionProblemLink{sec:inventory-strategy-4x3}{Inventory Strategy (ABB)} \\
		\SectionProblemLink{sec:event-venue-4x3}{Event Venue (ABC)} \\
		\SectionProblemLink{sec:supply-contract-4x3}{Supply Contract (AAA)}
	}
	
	
	% ----------------------------------------------------------------
	\subsection{Production System}
	\label{sec:production-system-4x3}
	
	A manufacturing company must choose among Production System A, System B, and
	System C. Four \GTenToneHL{state}{states} \GTenToneHL{state}{of} \GTenToneHL{state}{nature} represent different operating conditions. The
	final \GTenToneHL{payoff}{payoff} table, in thousands of dollars, is given below. Use \GTenToneHL{best}{Maximax},
	\GTenToneHL{worst}{Maximin}, and \GTenToneHL{maxregret}{Minimax} \GTenToneHL{maxregret}{Regret} to select a production system.
	
	\begin{center}
		\textit{Payoff table} \\
		\begin{tabular}{l c c c}
			\toprule
			State of nature & System A & System B & System C \\
			\midrule
			State \(\GTenToneHL{state}{S_1}\) & \GTenToneHL{payoff}{100} & \GTenToneHL{payoff}{70} & \GTenToneHL{payoff}{60} \\
			State \(\GTenToneHL{state}{S_2}\) & \GTenToneHL{payoff}{50}  & \GTenToneHL{payoff}{49} & \GTenToneHL{payoff}{30} \\
			State \(\GTenToneHL{state}{S_3}\) & \GTenToneHL{payoff}{50}  & \GTenToneHL{payoff}{90} & \GTenToneHL{payoff}{40} \\
			State \(\GTenToneHL{state}{S_4}\) & \GTenToneHL{payoff}{55}  & \GTenToneHL{payoff}{52} & \GTenToneHL{payoff}{35} \\
			\bottomrule
		\end{tabular}
	\end{center}
	
	\textbf{Questions/Tasks.}
	\begin{enumerate}
		\CriterionTask
		{subsec:production-system-4x3-c4}{4}
		{Applies the Maximin and Maximax criteria}

		\CriterionTask
		{subsec:production-system-4x3-c5}{5}
		{Constructs opportunity loss tables and applies the Minimax regret criterion}
	\end{enumerate}
	
	\subsubsection*{C4 - Applies the Maximin and Maximax criteria}
	\phantomsection
	\label{subsec:production-system-4x3-c4}
	
	\[
	\begin{aligned}
		\max(\text{System A}) &= \max\{100,50,50,55\}=100,\\
		\max(\text{System B}) &= \max\{70,49,90,52\}=90,\\
		\max(\text{System C}) &= \max\{60,30,40,35\}=60.
	\end{aligned}
	\]
	\[
	\begin{aligned}
		&\max\{\max(\text{System A}),\max(\text{System B}),
		\max(\text{System C})\}\\
		&\qquad=\max\{100,90,60\}
		=100.
	\end{aligned}
	\]
	The \GTenToneHL{best}{Maximax} choice is System A because it has the highest possible \GTenToneHL{payoff}{payoff}.
	
	\[
	\begin{aligned}
		\min(\text{System A}) &= \min\{100,50,50,55\}=50,\\
		\min(\text{System B}) &= \min\{70,49,90,52\}=49,\\
		\min(\text{System C}) &= \min\{60,30,40,35\}=30.
	\end{aligned}
	\]
	\[
	\begin{aligned}
		&\max\{\min(\text{System A}),\min(\text{System B}),
		\min(\text{System C})\}\\
		&\qquad=\max\{50,49,30\}
		=50.
	\end{aligned}
	\]
	The \GTenToneHL{worst}{Maximin} choice is System A because it has the largest worst-case \GTenToneHL{payoff}{payoff}.
	
	\subsubsection*{C5 - Constructs opportunity loss tables and applies the Minimax regret criterion}
	\phantomsection
	\label{subsec:production-system-4x3-c5}
	
	Best \GTenToneHL{payoff}{payoff} in each state:
	\[
	\begin{aligned}
		\text{State }\GTenToneHL{state}{S_1}:\;&\max\{100,70,60\}=100,\\
		\text{State }\GTenToneHL{state}{S_2}:\;&\max\{50,49,30\}=50,\\
		\text{State }\GTenToneHL{state}{S_3}:\;&\max\{50,90,40\}=90,\\
		\text{State }\GTenToneHL{state}{S_4}:\;&\max\{55,52,35\}=55.
	\end{aligned}
	\]
	
	Regret table (best \GTenToneHL{payoff}{payoff} \( - \) \GTenToneHL{payoff}{payoff}):
	\begin{center}
		\begin{tabular}{l c c c c c}
			\toprule
			Alternative & \(\GTenToneHL{state}{S_1}\) & \(\GTenToneHL{state}{S_2}\) & \(\GTenToneHL{state}{S_3}\) & \(\GTenToneHL{state}{S_4}\)
			& \GTenToneHL{maxregret}{Maximum} \GTenToneHL{maxregret}{regret} \\
			\midrule
			System A
			& \(100-100=0\)
			& \(50-50=0\)
			& \(90-50=40\)
			& \(55-55=0\)
			& \GTenToneHL{maxregret}{40} \\
			System B
			& \(100-70=30\)
			& \(50-49=1\)
			& \(90-90=0\)
			& \(55-52=3\)
			& \GTenToneHL{maxregret}{30} \\
			System C
			& \(100-60=40\)
			& \(50-30=20\)
			& \(90-40=50\)
			& \(55-35=20\)
			& \GTenToneHL{maxregret}{50} \\
			\bottomrule
		\end{tabular}
	\end{center}
	
	The \GTenToneHL{maxregret}{Minimax} \GTenToneHL{maxregret}{Regret} choice is System B because it has the smallest maximum
	\GTenToneHL{regret}{regret}, \(30\).
	
	\subsubsection*{Decisions summary}
	
	\begin{center}
		\begin{tabular}{l c}
			\toprule
			Criterion & Final decision \\
			\midrule
			\GTenToneHL{best}{Maximax}        & System A \\
			\GTenToneHL{worst}{Maximin}        & System A \\
			\GTenToneHL{maxregret}{Minimax} \GTenToneHL{maxregret}{Regret} & System B \\
			\bottomrule
		\end{tabular}
	\end{center}
	
	\secInicio[sec:production-system-4x3][sec:problems-4x3]
	
	
	% ----------------------------------------------------------------
	\subsection{Service Network}
	\label{sec:service-network-4x3}
	
	A service company must choose among Service Network A, Network B, and Network C.
	Four \GTenToneHL{state}{states} \GTenToneHL{state}{of} \GTenToneHL{state}{nature} represent different demand conditions. The final \GTenToneHL{payoff}{payoff}
	table, in thousands of dollars, is given below. Use \GTenToneHL{best}{Maximax}, \GTenToneHL{worst}{Maximin}, and
	\GTenToneHL{maxregret}{Minimax} \GTenToneHL{maxregret}{Regret} to select a service network.
	
	\begin{center}
		\textit{Payoff table} \\
		\begin{tabular}{l c c c}
			\toprule
			State of nature & Network A & Network B & Network C \\
			\midrule
			State \(\GTenToneHL{state}{S_1}\) & \GTenToneHL{payoff}{100} & \GTenToneHL{payoff}{80} & \GTenToneHL{payoff}{60} \\
			State \(\GTenToneHL{state}{S_2}\) & \GTenToneHL{payoff}{45}  & \GTenToneHL{payoff}{50} & \GTenToneHL{payoff}{30} \\
			State \(\GTenToneHL{state}{S_3}\) & \GTenToneHL{payoff}{40}  & \GTenToneHL{payoff}{50} & \GTenToneHL{payoff}{35} \\
			State \(\GTenToneHL{state}{S_4}\) & \GTenToneHL{payoff}{48}  & \GTenToneHL{payoff}{52} & \GTenToneHL{payoff}{40} \\
			\bottomrule
		\end{tabular}
	\end{center}
	
	\textbf{Questions/Tasks.}
	\begin{enumerate}
		\CriterionTask
		{subsec:service-network-4x3-c4}{4}
		{Applies the Maximin and Maximax criteria}

		\CriterionTask
		{subsec:service-network-4x3-c5}{5}
		{Constructs opportunity loss tables and applies the Minimax regret criterion}
	\end{enumerate}
	
	\subsubsection*{C4 - Applies the Maximin and Maximax criteria}
	\phantomsection
	\label{subsec:service-network-4x3-c4}
	
	\[
	\begin{aligned}
		\max(\text{Network A}) &= \max\{100,45,40,48\}=100,\\
		\max(\text{Network B}) &= \max\{80,50,50,52\}=80,\\
		\max(\text{Network C}) &= \max\{60,30,35,40\}=60.
	\end{aligned}
	\]
	\[
	\begin{aligned}
		&\max\{\max(\text{Network A}),\max(\text{Network B}),
		\max(\text{Network C})\}\\
		&\qquad=\max\{100,80,60\}
		=100.
	\end{aligned}
	\]
	The \GTenToneHL{best}{Maximax} choice is Network A because it has the highest possible \GTenToneHL{payoff}{payoff}.
	
	\[
	\begin{aligned}
		\min(\text{Network A}) &= \min\{100,45,40,48\}=40,\\
		\min(\text{Network B}) &= \min\{80,50,50,52\}=50,\\
		\min(\text{Network C}) &= \min\{60,30,35,40\}=30.
	\end{aligned}
	\]
	\[
	\begin{aligned}
		&\max\{\min(\text{Network A}),\min(\text{Network B}),
		\min(\text{Network C})\}\\
		&\qquad=\max\{40,50,30\}
		=50.
	\end{aligned}
	\]
	The \GTenToneHL{worst}{Maximin} choice is Network B because it has the largest worst-case \GTenToneHL{payoff}{payoff}.
	
	\subsubsection*{C5 - Constructs opportunity loss tables and applies the Minimax regret criterion}
	\phantomsection
	\label{subsec:service-network-4x3-c5}
	
	Best \GTenToneHL{payoff}{payoff} in each state:
	\[
	\begin{aligned}
		\text{State }\GTenToneHL{state}{S_1}:\;&\max\{100,80,60\}=100,\\
		\text{State }\GTenToneHL{state}{S_2}:\;&\max\{45,50,30\}=50,\\
		\text{State }\GTenToneHL{state}{S_3}:\;&\max\{40,50,35\}=50,\\
		\text{State }\GTenToneHL{state}{S_4}:\;&\max\{48,52,40\}=52.
	\end{aligned}
	\]
	
	Regret table (best \GTenToneHL{payoff}{payoff} \( - \) \GTenToneHL{payoff}{payoff}):
	\begin{center}
		\begin{tabular}{l c c c c c}
			\toprule
			Alternative & \(\GTenToneHL{state}{S_1}\) & \(\GTenToneHL{state}{S_2}\) & \(\GTenToneHL{state}{S_3}\) & \(\GTenToneHL{state}{S_4}\)
			& \GTenToneHL{maxregret}{Maximum} \GTenToneHL{maxregret}{regret} \\
			\midrule
			Network A
			& \(100-100=0\)
			& \(50-45=5\)
			& \(50-40=10\)
			& \(52-48=4\)
			& \GTenToneHL{maxregret}{10} \\
			Network B
			& \(100-80=20\)
			& \(50-50=0\)
			& \(50-50=0\)
			& \(52-52=0\)
			& \GTenToneHL{maxregret}{20} \\
			Network C
			& \(100-60=40\)
			& \(50-30=20\)
			& \(50-35=15\)
			& \(52-40=12\)
			& \GTenToneHL{maxregret}{40} \\
			\bottomrule
		\end{tabular}
	\end{center}
	
	The \GTenToneHL{maxregret}{Minimax} \GTenToneHL{maxregret}{Regret} choice is Network A because it has the smallest maximum
	\GTenToneHL{regret}{regret}, \(10\).
	
	\subsubsection*{Decisions summary}
	
	\begin{center}
		\begin{tabular}{l c}
			\toprule
			Criterion & Final decision \\
			\midrule
			\GTenToneHL{best}{Maximax}        & Network A \\
			\GTenToneHL{worst}{Maximin}        & Network B \\
			\GTenToneHL{maxregret}{Minimax} \GTenToneHL{maxregret}{Regret} & Network A \\
			\bottomrule
		\end{tabular}
	\end{center}
	
	\secInicio[sec:service-network-4x3][sec:problems-4x3]
	
	
	% ----------------------------------------------------------------
	\subsection{Inventory Strategy}
	\label{sec:inventory-strategy-4x3}
	
	A retailer must choose among Inventory Strategy A, Strategy B, and Strategy C.
	Four \GTenToneHL{state}{states} \GTenToneHL{state}{of} \GTenToneHL{state}{nature} represent different demand conditions. The final \GTenToneHL{payoff}{payoff}
	table, in thousands of dollars, is given below. Use \GTenToneHL{best}{Maximax}, \GTenToneHL{worst}{Maximin}, and
	\GTenToneHL{maxregret}{Minimax} \GTenToneHL{maxregret}{Regret} to select an inventory strategy.
	
	\begin{center}
		\textit{Payoff table} \\
		\begin{tabular}{l c c c}
			\toprule
			State of nature & Strategy A & Strategy B & Strategy C \\
			\midrule
			State \(\GTenToneHL{state}{S_1}\) & \GTenToneHL{rate}{100} & \GTenToneHL{rate}{70} & \GTenToneHL{rate}{50} \\
			State \(\GTenToneHL{state}{S_2}\) & \GTenToneHL{rate}{20}  & \GTenToneHL{rate}{60} & \GTenToneHL{rate}{40} \\
			State \(\GTenToneHL{state}{S_3}\) & \GTenToneHL{rate}{0}   & \GTenToneHL{rate}{50} & \GTenToneHL{rate}{30} \\
			State \(\GTenToneHL{state}{S_4}\) & \GTenToneHL{rate}{25}  & \GTenToneHL{rate}{55} & \GTenToneHL{rate}{45} \\
			\bottomrule
		\end{tabular}
	\end{center}
	
	\textbf{Questions/Tasks.}
	\begin{enumerate}
		\CriterionTask
		{subsec:inventory-strategy-4x3-c4}{4}
		{Applies the Maximin and Maximax criteria}

		\CriterionTask
		{subsec:inventory-strategy-4x3-c5}{5}
		{Constructs opportunity loss tables and applies the Minimax regret criterion}
	\end{enumerate}
	
	\subsubsection*{C4 - Applies the Maximin and Maximax criteria}
	\phantomsection
	\label{subsec:inventory-strategy-4x3-c4}
	
	\[
	\begin{aligned}
		\max(\text{Strategy A}) &= \max\{100,20,0,25\}=100,\\
		\max(\text{Strategy B}) &= \max\{70,60,50,55\}=70,\\
		\max(\text{Strategy C}) &= \max\{50,40,30,45\}=50.
	\end{aligned}
	\]
	\[
	\begin{aligned}
		&\max\{\max(\text{Strategy A}),\max(\text{Strategy B}),
		\max(\text{Strategy C})\}\\
		&\qquad=\max\{100,70,50\}
		=100.
	\end{aligned}
	\]
	The \GTenToneHL{best}{Maximax} choice is Strategy A because it has the highest possible \GTenToneHL{payoff}{payoff}.
	
	\[
	\begin{aligned}
		\min(\text{Strategy A}) &= \min\{100,20,0,25\}=0,\\
		\min(\text{Strategy B}) &= \min\{70,60,50,55\}=50,\\
		\min(\text{Strategy C}) &= \min\{50,40,30,45\}=30.
	\end{aligned}
	\]
	\[
	\begin{aligned}
		&\max\{\min(\text{Strategy A}),\min(\text{Strategy B}),
		\min(\text{Strategy C})\}\\
		&\qquad=\max\{0,50,30\}
		=50.
	\end{aligned}
	\]
	The \GTenToneHL{worst}{Maximin} choice is Strategy B because it has the largest worst-case \GTenToneHL{payoff}{payoff}.
	
	\subsubsection*{C5 - Constructs opportunity loss tables and applies the Minimax regret criterion}
	\phantomsection
	\label{subsec:inventory-strategy-4x3-c5}
	
	Best \GTenToneHL{payoff}{payoff} in each state:
	\[
	\begin{aligned}
		\text{State }\GTenToneHL{state}{S_1}:\;&\max\{100,70,50\}=100,\\
		\text{State }\GTenToneHL{state}{S_2}:\;&\max\{20,60,40\}=60,\\
		\text{State }\GTenToneHL{state}{S_3}:\;&\max\{0,50,30\}=50,\\
		\text{State }\GTenToneHL{state}{S_4}:\;&\max\{25,55,45\}=55.
	\end{aligned}
	\]
	
	Regret table (best \GTenToneHL{payoff}{payoff} \( - \) \GTenToneHL{payoff}{payoff}):
	\begin{center}
		\begin{tabular}{l c c c c c}
			\toprule
			Alternative & \(\GTenToneHL{state}{S_1}\) & \(\GTenToneHL{state}{S_2}\) & \(\GTenToneHL{state}{S_3}\) & \(\GTenToneHL{state}{S_4}\)
			& \GTenToneHL{maxregret}{Maximum} \GTenToneHL{maxregret}{regret} \\
			\midrule
			Strategy A
			& \(100-100=0\)
			& \(60-20=40\)
			& \(50-0=50\)
			& \(55-25=30\)
			& \GTenToneHL{maxregret}{50} \\
			Strategy B
			& \(100-70=30\)
			& \(60-60=0\)
			& \(50-50=0\)
			& \(55-55=0\)
			& \GTenToneHL{maxregret}{30} \\
			Strategy C
			& \(100-50=50\)
			& \(60-40=20\)
			& \(50-30=20\)
			& \(55-45=10\)
			& \GTenToneHL{maxregret}{50} \\
			\bottomrule
		\end{tabular}
	\end{center}
	
	The \GTenToneHL{maxregret}{Minimax} \GTenToneHL{maxregret}{Regret} choice is Strategy B because it has the smallest maximum
	\GTenToneHL{regret}{regret}, \(30\).
	
	\subsubsection*{Decisions summary}
	
	\begin{center}
		\begin{tabular}{l c}
			\toprule
			Criterion & Final decision \\
			\midrule
			\GTenToneHL{best}{Maximax}        & Strategy A \\
			\GTenToneHL{worst}{Maximin}        & Strategy B \\
			\GTenToneHL{maxregret}{Minimax} \GTenToneHL{maxregret}{Regret} & Strategy B \\
			\bottomrule
		\end{tabular}
	\end{center}
	
	\secInicio[sec:inventory-strategy-4x3][sec:problems-4x3]
	
	
	% ----------------------------------------------------------------
	\subsection{Event Venue}
	\label{sec:event-venue-4x3}
	
	An event company must choose among Venue A, Venue B, and Venue C. Four states
	of nature represent different attendance levels. The final \GTenToneHL{payoff}{payoff} table, in
	thousands of dollars, is given below. Use \GTenToneHL{best}{Maximax}, \GTenToneHL{worst}{Maximin}, and \GTenToneHL{maxregret}{Minimax} \GTenToneHL{maxregret}{Regret}
	to select an event venue.
	
	\begin{center}
		\textit{Payoff table} \\
		\begin{tabular}{l c c c}
			\toprule
			State of nature & Venue A & Venue B & Venue C \\
			\midrule
			State \(\GTenToneHL{state}{S_1}\) & \GTenToneHL{payoff}{100} & \GTenToneHL{payoff}{70} & \GTenToneHL{payoff}{80} \\
			State \(\GTenToneHL{state}{S_2}\) & \GTenToneHL{payoff}{0}   & \GTenToneHL{payoff}{60} & \GTenToneHL{payoff}{50} \\
			State \(\GTenToneHL{state}{S_3}\) & \GTenToneHL{payoff}{20}  & \GTenToneHL{payoff}{55} & \GTenToneHL{payoff}{50} \\
			State \(\GTenToneHL{state}{S_4}\) & \GTenToneHL{payoff}{40}  & \GTenToneHL{payoff}{58} & \GTenToneHL{payoff}{52} \\
			\bottomrule
		\end{tabular}
	\end{center}
	
	\textbf{Questions/Tasks.}
	\begin{enumerate}
		\CriterionTask
		{subsec:event-venue-4x3-c4}{4}
		{Applies the Maximin and Maximax criteria}

		\CriterionTask
		{subsec:event-venue-4x3-c5}{5}
		{Constructs opportunity loss tables and applies the Minimax regret criterion}
	\end{enumerate}
	
	\subsubsection*{C4 - Applies the Maximin and Maximax criteria}
	\phantomsection
	\label{subsec:event-venue-4x3-c4}
	
	\[
	\begin{aligned}
		\max(\text{Venue A}) &= \max\{100,0,20,40\}=100,\\
		\max(\text{Venue B}) &= \max\{70,60,55,58\}=70,\\
		\max(\text{Venue C}) &= \max\{80,50,50,52\}=80.
	\end{aligned}
	\]
	\[
	\begin{aligned}
		&\max\{\max(\text{Venue A}),\max(\text{Venue B}),
		\max(\text{Venue C})\}\\
		&\qquad=\max\{100,70,80\}
		=100.
	\end{aligned}
	\]
	The \GTenToneHL{best}{Maximax} choice is Venue A because it has the highest possible \GTenToneHL{payoff}{payoff}.
	
	\[
	\begin{aligned}
		\min(\text{Venue A}) &= \min\{100,0,20,40\}=0,\\
		\min(\text{Venue B}) &= \min\{70,60,55,58\}=55,\\
		\min(\text{Venue C}) &= \min\{80,50,50,52\}=50.
	\end{aligned}
	\]
	\[
	\begin{aligned}
		&\max\{\min(\text{Venue A}),\min(\text{Venue B}),
		\min(\text{Venue C})\}\\
		&\qquad=\max\{0,55,50\}
		=55.
	\end{aligned}
	\]
	The \GTenToneHL{worst}{Maximin} choice is Venue B because it has the largest worst-case \GTenToneHL{payoff}{payoff}.
	
	\subsubsection*{C5 - Constructs opportunity loss tables and applies the Minimax regret criterion}
	\phantomsection
	\label{subsec:event-venue-4x3-c5}
	
	Best \GTenToneHL{payoff}{payoff} in each state:
	\[
	\begin{aligned}
		\text{State }\GTenToneHL{state}{S_1}:\;&\max\{100,70,80\}=100,\\
		\text{State }\GTenToneHL{state}{S_2}:\;&\max\{0,60,50\}=60,\\
		\text{State }\GTenToneHL{state}{S_3}:\;&\max\{20,55,50\}=55,\\
		\text{State }\GTenToneHL{state}{S_4}:\;&\max\{40,58,52\}=58.
	\end{aligned}
	\]
	
	Regret table (best \GTenToneHL{payoff}{payoff} \( - \) \GTenToneHL{payoff}{payoff}):
	\begin{center}
		\begin{tabular}{l c c c c c}
			\toprule
			Alternative & \(\GTenToneHL{state}{S_1}\) & \(\GTenToneHL{state}{S_2}\) & \(\GTenToneHL{state}{S_3}\) & \(\GTenToneHL{state}{S_4}\)
			& \GTenToneHL{maxregret}{Maximum} \GTenToneHL{maxregret}{regret} \\
			\midrule
			Venue A
			& \(100-100=0\)
			& \(60-0=60\)
			& \(55-20=35\)
			& \(58-40=18\)
			& \GTenToneHL{maxregret}{60} \\
			Venue B
			& \(100-70=30\)
			& \(60-60=0\)
			& \(55-55=0\)
			& \(58-58=0\)
			& \GTenToneHL{maxregret}{30} \\
			Venue C
			& \(100-80=20\)
			& \(60-50=10\)
			& \(55-50=5\)
			& \(58-52=6\)
			& \GTenToneHL{maxregret}{20} \\
			\bottomrule
		\end{tabular}
	\end{center}
	
	The \GTenToneHL{maxregret}{Minimax} \GTenToneHL{maxregret}{Regret} choice is Venue C because it has the smallest maximum
	\GTenToneHL{regret}{regret}, \(20\).
	
	\subsubsection*{Decisions summary}
	
	\begin{center}
		\begin{tabular}{l c}
			\toprule
			Criterion & Final decision \\
			\midrule
			\GTenToneHL{best}{Maximax}        & Venue A \\
			\GTenToneHL{worst}{Maximin}        & Venue B \\
			\GTenToneHL{maxregret}{Minimax} \GTenToneHL{maxregret}{Regret} & Venue C \\
			\bottomrule
		\end{tabular}
	\end{center}
	
	\secInicio[sec:event-venue-4x3][sec:problems-4x3]
	
	
	% ----------------------------------------------------------------
	\subsection{Supply Contract}
	\label{sec:supply-contract-4x3}
	
	A company must choose among Supply Contract A, Contract B, and Contract C.
	Four \GTenToneHL{state}{states} \GTenToneHL{state}{of} \GTenToneHL{state}{nature} represent different market and supply conditions. The
	final \GTenToneHL{payoff}{payoff} table, in thousands of dollars, is given below. Use \GTenToneHL{best}{Maximax},
	\GTenToneHL{worst}{Maximin}, and \GTenToneHL{maxregret}{Minimax} \GTenToneHL{maxregret}{Regret} to select a supply contract.
	
	\begin{center}
		\textit{Payoff table} \\
		\begin{tabular}{l c c c}
			\toprule
			State of nature & Contract A & Contract B & Contract C \\
			\midrule
			State \(\GTenToneHL{state}{S_1}\) & \GTenToneHL{payoff}{90} & \GTenToneHL{payoff}{80} & \GTenToneHL{payoff}{70} \\
			State \(\GTenToneHL{state}{S_2}\) & \GTenToneHL{payoff}{70} & \GTenToneHL{payoff}{50} & \GTenToneHL{payoff}{45} \\
			State \(\GTenToneHL{state}{S_3}\) & \GTenToneHL{payoff}{60} & \GTenToneHL{payoff}{40} & \GTenToneHL{payoff}{35} \\
			State \(\GTenToneHL{state}{S_4}\) & \GTenToneHL{payoff}{65} & \GTenToneHL{payoff}{55} & \GTenToneHL{payoff}{50} \\
			\bottomrule
		\end{tabular}
	\end{center}
	
	\textbf{Questions/Tasks.}
	\begin{enumerate}
		\CriterionTask
		{subsec:supply-contract-4x3-c4}{4}
		{Applies the Maximin and Maximax criteria}

		\CriterionTask
		{subsec:supply-contract-4x3-c5}{5}
		{Constructs opportunity loss tables and applies the Minimax regret criterion}
	\end{enumerate}
	
	\subsubsection*{C4 - Applies the Maximin and Maximax criteria}
	\phantomsection
	\label{subsec:supply-contract-4x3-c4}
	
	\[
	\begin{aligned}
		\max(\text{Contract A}) &= \max\{90,70,60,65\}=90,\\
		\max(\text{Contract B}) &= \max\{80,50,40,55\}=80,\\
		\max(\text{Contract C}) &= \max\{70,45,35,50\}=70.
	\end{aligned}
	\]
	\[
	\begin{aligned}
		&\max\{\max(\text{Contract A}),\max(\text{Contract B}),
		\max(\text{Contract C})\}\\
		&\qquad=\max\{90,80,70\}
		=90.
	\end{aligned}
	\]
	The \GTenToneHL{best}{Maximax} choice is Contract A because it has the highest possible \GTenToneHL{payoff}{payoff}.
	
	\[
	\begin{aligned}
		\min(\text{Contract A}) &= \min\{90,70,60,65\}=60,\\
		\min(\text{Contract B}) &= \min\{80,50,40,55\}=40,\\
		\min(\text{Contract C}) &= \min\{70,45,35,50\}=35.
	\end{aligned}
	\]
	\[
	\begin{aligned}
		&\max\{\min(\text{Contract A}),\min(\text{Contract B}),
		\min(\text{Contract C})\}\\
		&\qquad=\max\{60,40,35\}
		=60.
	\end{aligned}
	\]
	The \GTenToneHL{worst}{Maximin} choice is Contract A because it has the largest worst-case \GTenToneHL{payoff}{payoff}.
	
	\subsubsection*{C5 - Constructs opportunity loss tables and applies the Minimax regret criterion}
	\phantomsection
	\label{subsec:supply-contract-4x3-c5}
	
	Best \GTenToneHL{payoff}{payoff} in each state:
	\[
	\begin{aligned}
		\text{State }\GTenToneHL{state}{S_1}:\;&\max\{90,80,70\}=90,\\
		\text{State }\GTenToneHL{state}{S_2}:\;&\max\{70,50,45\}=70,\\
		\text{State }\GTenToneHL{state}{S_3}:\;&\max\{60,40,35\}=60,\\
		\text{State }\GTenToneHL{state}{S_4}:\;&\max\{65,55,50\}=65.
	\end{aligned}
	\]
	
	Regret table (best \GTenToneHL{payoff}{payoff} \( - \) \GTenToneHL{payoff}{payoff}):
	\begin{center}
		\begin{tabular}{l c c c c c}
			\toprule
			Alternative & \(\GTenToneHL{state}{S_1}\) & \(\GTenToneHL{state}{S_2}\) & \(\GTenToneHL{state}{S_3}\) & \(\GTenToneHL{state}{S_4}\)
			& \GTenToneHL{maxregret}{Maximum} \GTenToneHL{maxregret}{regret} \\
			\midrule
			Contract A
			& \(90-90=0\)
			& \(70-70=0\)
			& \(60-60=0\)
			& \(65-65=0\)
			& \GTenToneHL{maxregret}{0} \\
			Contract B
			& \(90-80=10\)
			& \(70-50=20\)
			& \(60-40=20\)
			& \(65-55=10\)
			& \GTenToneHL{maxregret}{20} \\
			Contract C
			& \(90-70=20\)
			& \(70-45=25\)
			& \(60-35=25\)
			& \(65-50=15\)
			& \GTenToneHL{maxregret}{25} \\
			\bottomrule
		\end{tabular}
	\end{center}
	
	The \GTenToneHL{maxregret}{Minimax} \GTenToneHL{maxregret}{Regret} choice is Contract A because it has the smallest maximum
	\GTenToneHL{regret}{regret}, \(0\).
	
	\subsubsection*{Decisions summary}
	
	\begin{center}
		\begin{tabular}{l c}
			\toprule
			Criterion & Final decision \\
			\midrule
			\GTenToneHL{best}{Maximax}        & Contract A \\
			\GTenToneHL{worst}{Maximin}        & Contract A \\
			\GTenToneHL{maxregret}{Minimax} \GTenToneHL{maxregret}{Regret} & Contract A \\
			\bottomrule
		\end{tabular}
	\end{center}
	
	\secInicio[sec:supply-contract-4x3][sec:problems-4x3]
	
	
	% ================================================================
	\section{Problems \(3\times 4\)}
	\label{sec:problems-3x4}
	
	\SectionProblemsTableOfContents{
		\SectionProblemLink{sec:production-system-3x4}{Production System (AAB)} \\
		\SectionProblemLink{sec:service-network-3x4}{Service Network (ABA)} \\
		\SectionProblemLink{sec:inventory-strategy-3x4}{Inventory Strategy (ABB)} \\
		\SectionProblemLink{sec:event-venue-3x4}{Event Venue (ABC)} \\
		\SectionProblemLink{sec:supply-contract-3x4}{Supply Contract (AAA)}
	}
	
	% ----------------------------------------------------------------
	\subsection{Production System}
	\label{sec:production-system-3x4}
	
	A manufacturing company must choose among Production System A, System B,
	System C, and System D. Three \GTenToneHL{state}{states} \GTenToneHL{state}{of} \GTenToneHL{state}{nature} represent different operating
	conditions. The final \GTenToneHL{payoff}{payoff} table, in thousands of dollars, is given below.
	Use \GTenToneHL{best}{Maximax}, \GTenToneHL{worst}{Maximin}, and \GTenToneHL{maxregret}{Minimax} \GTenToneHL{maxregret}{Regret} to select a production system.
	
	\begin{center}
		\textit{Payoff table} \\
		\begin{tabular}{l c c c c}
			\toprule
			State of nature & System A & System B & System C & System D \\
			\midrule
			State \(\GTenToneHL{state}{S_1}\) & \GTenToneHL{payoff}{100} & \GTenToneHL{payoff}{70} & \GTenToneHL{payoff}{60} & \GTenToneHL{payoff}{55} \\
			State \(\GTenToneHL{state}{S_2}\) & \GTenToneHL{payoff}{50}  & \GTenToneHL{payoff}{49} & \GTenToneHL{payoff}{30} & \GTenToneHL{payoff}{35} \\
			State \(\GTenToneHL{state}{S_3}\) & \GTenToneHL{payoff}{50}  & \GTenToneHL{payoff}{90} & \GTenToneHL{payoff}{40} & \GTenToneHL{payoff}{45} \\
			\bottomrule
		\end{tabular}
	\end{center}
	
	\textbf{Questions/Tasks.}
	\begin{enumerate}
		\CriterionTask
		{subsec:production-system-3x4-c4}{4}
		{Applies the Maximin and Maximax criteria}

		\CriterionTask
		{subsec:production-system-3x4-c5}{5}
		{Constructs opportunity loss tables and applies the Minimax regret criterion}
	\end{enumerate}
	
	\subsubsection*{C4 - Applies the Maximin and Maximax criteria}
	\phantomsection
	\label{subsec:production-system-3x4-c4}
	
	\[
	\begin{aligned}
		\max(\text{System A}) &= \max\{100,50,50\}=100,\\
		\max(\text{System B}) &= \max\{70,49,90\}=90,\\
		\max(\text{System C}) &= \max\{60,30,40\}=60,\\
		\max(\text{System D}) &= \max\{55,35,45\}=55.
	\end{aligned}
	\]
	\[
	\begin{aligned}
		&\max\{\max(\text{System A}),\max(\text{System B}),
		\max(\text{System C}),\max(\text{System D})\}\\
		&\qquad=\max\{100,90,60,55\}
		=100.
	\end{aligned}
	\]
	The \GTenToneHL{best}{Maximax} choice is System A because it has the highest possible \GTenToneHL{payoff}{payoff}.
	
	\[
	\begin{aligned}
		\min(\text{System A}) &= \min\{100,50,50\}=50,\\
		\min(\text{System B}) &= \min\{70,49,90\}=49,\\
		\min(\text{System C}) &= \min\{60,30,40\}=30,\\
		\min(\text{System D}) &= \min\{55,35,45\}=35.
	\end{aligned}
	\]
	\[
	\begin{aligned}
		&\max\{\min(\text{System A}),\min(\text{System B}),
		\min(\text{System C}),\min(\text{System D})\}\\
		&\qquad=\max\{50,49,30,35\}
		=50.
	\end{aligned}
	\]
	The \GTenToneHL{worst}{Maximin} choice is System A because it has the largest worst-case \GTenToneHL{payoff}{payoff}.
	
	\subsubsection*{C5 - Constructs opportunity loss tables and applies the Minimax regret criterion}
	\phantomsection
	\label{subsec:production-system-3x4-c5}
	
	Best \GTenToneHL{payoff}{payoff} in each state:
	\[
	\begin{aligned}
		\text{State }\GTenToneHL{state}{S_1}:\;&\max\{100,70,60,55\}=100,\\
		\text{State }\GTenToneHL{state}{S_2}:\;&\max\{50,49,30,35\}=50,\\
		\text{State }\GTenToneHL{state}{S_3}:\;&\max\{50,90,40,45\}=90.
	\end{aligned}
	\]
	
	Regret table (best \GTenToneHL{payoff}{payoff} \( - \) \GTenToneHL{payoff}{payoff}):
	\begin{center}
		\begin{tabular}{l c c c c}
			\toprule
			Alternative & \(\GTenToneHL{state}{S_1}\) & \(\GTenToneHL{state}{S_2}\) & \(\GTenToneHL{state}{S_3}\) & \GTenToneHL{maxregret}{Maximum} \GTenToneHL{maxregret}{regret} \\
			\midrule
			System A
			& \(100-100=0\)
			& \(50-50=0\)
			& \(90-50=40\)
			& \GTenToneHL{maxregret}{40} \\
			System B
			& \(100-70=30\)
			& \(50-49=1\)
			& \(90-90=0\)
			& \GTenToneHL{maxregret}{30} \\
			System C
			& \(100-60=40\)
			& \(50-30=20\)
			& \(90-40=50\)
			& \GTenToneHL{maxregret}{50} \\
			System D
			& \(100-55=45\)
			& \(50-35=15\)
			& \(90-45=45\)
			& \GTenToneHL{maxregret}{45} \\
			\bottomrule
		\end{tabular}
	\end{center}
	
	The \GTenToneHL{maxregret}{Minimax} \GTenToneHL{maxregret}{Regret} choice is System B because it has the smallest maximum
	\GTenToneHL{regret}{regret}, \(30\).
	
	\subsubsection*{Decisions summary}
	
	\begin{center}
		\begin{tabular}{l c}
			\toprule
			Criterion & Final decision \\
			\midrule
			\GTenToneHL{best}{Maximax}        & System A \\
			\GTenToneHL{worst}{Maximin}        & System A \\
			\GTenToneHL{maxregret}{Minimax} \GTenToneHL{maxregret}{Regret} & System B \\
			\bottomrule
		\end{tabular}
	\end{center}
	
	\secInicio[sec:production-system-3x4][sec:problems-3x4]
	
	
	% ----------------------------------------------------------------
	\subsection{Service Network}
	\label{sec:service-network-3x4}
	
	A service company must choose among Service Network A, Network B, Network C,
	and Network D. Three \GTenToneHL{state}{states} \GTenToneHL{state}{of} \GTenToneHL{state}{nature} represent different demand conditions.
	The final \GTenToneHL{payoff}{payoff} table, in thousands of dollars, is given below. Use \GTenToneHL{best}{Maximax},
	\GTenToneHL{worst}{Maximin}, and \GTenToneHL{maxregret}{Minimax} \GTenToneHL{maxregret}{Regret} to select a service network.
	
	\begin{center}
		\textit{Payoff table} \\
		\begin{tabular}{l c c c c}
			\toprule
			State of nature & Network A & Network B & Network C & Network D \\
			\midrule
			State \(\GTenToneHL{state}{S_1}\) & \GTenToneHL{payoff}{100} & \GTenToneHL{payoff}{80} & \GTenToneHL{payoff}{60} & \GTenToneHL{payoff}{70} \\
			State \(\GTenToneHL{state}{S_2}\) & \GTenToneHL{payoff}{45}  & \GTenToneHL{payoff}{50} & \GTenToneHL{payoff}{30} & \GTenToneHL{payoff}{40} \\
			State \(\GTenToneHL{state}{S_3}\) & \GTenToneHL{payoff}{40}  & \GTenToneHL{payoff}{50} & \GTenToneHL{payoff}{35} & \GTenToneHL{payoff}{45} \\
			\bottomrule
		\end{tabular}
	\end{center}
	
	\textbf{Questions/Tasks.}
	\begin{enumerate}
		\CriterionTask
		{subsec:service-network-3x4-c4}{4}
		{Applies the Maximin and Maximax criteria}

		\CriterionTask
		{subsec:service-network-3x4-c5}{5}
		{Constructs opportunity loss tables and applies the Minimax regret criterion}
	\end{enumerate}
	
	\subsubsection*{C4 - Applies the Maximin and Maximax criteria}
	\phantomsection
	\label{subsec:service-network-3x4-c4}
	
	\[
	\begin{aligned}
		\max(\text{Network A}) &= \max\{100,45,40\}=100,\\
		\max(\text{Network B}) &= \max\{80,50,50\}=80,\\
		\max(\text{Network C}) &= \max\{60,30,35\}=60,\\
		\max(\text{Network D}) &= \max\{70,40,45\}=70.
	\end{aligned}
	\]
	\[
	\begin{aligned}
		&\max\{\max(\text{Network A}),\max(\text{Network B}),
		\max(\text{Network C}),\max(\text{Network D})\}\\
		&\qquad=\max\{100,80,60,70\}
		=100.
	\end{aligned}
	\]
	The \GTenToneHL{best}{Maximax} choice is Network A because it has the highest possible \GTenToneHL{payoff}{payoff}.
	
	\[
	\begin{aligned}
		\min(\text{Network A}) &= \min\{100,45,40\}=40,\\
		\min(\text{Network B}) &= \min\{80,50,50\}=50,\\
		\min(\text{Network C}) &= \min\{60,30,35\}=30,\\
		\min(\text{Network D}) &= \min\{70,40,45\}=40.
	\end{aligned}
	\]
	\[
	\begin{aligned}
		&\max\{\min(\text{Network A}),\min(\text{Network B}),
		\min(\text{Network C}),\min(\text{Network D})\}\\
		&\qquad=\max\{40,50,30,40\}
		=50.
	\end{aligned}
	\]
	The \GTenToneHL{worst}{Maximin} choice is Network B because it has the largest worst-case \GTenToneHL{payoff}{payoff}.
	
	\subsubsection*{C5 - Constructs opportunity loss tables and applies the Minimax regret criterion}
	\phantomsection
	\label{subsec:service-network-3x4-c5}
	
	Best \GTenToneHL{payoff}{payoff} in each state:
	\[
	\begin{aligned}
		\text{State }\GTenToneHL{state}{S_1}:\;&\max\{100,80,60,70\}=100,\\
		\text{State }\GTenToneHL{state}{S_2}:\;&\max\{45,50,30,40\}=50,\\
		\text{State }\GTenToneHL{state}{S_3}:\;&\max\{40,50,35,45\}=50.
	\end{aligned}
	\]
	
	Regret table (best \GTenToneHL{payoff}{payoff} \( - \) \GTenToneHL{payoff}{payoff}):
	\begin{center}
		\begin{tabular}{l c c c c}
			\toprule
			Alternative & \(\GTenToneHL{state}{S_1}\) & \(\GTenToneHL{state}{S_2}\) & \(\GTenToneHL{state}{S_3}\) & \GTenToneHL{maxregret}{Maximum} \GTenToneHL{maxregret}{regret} \\
			\midrule
			Network A
			& \(100-100=0\)
			& \(50-45=5\)
			& \(50-40=10\)
			& \GTenToneHL{maxregret}{10} \\
			Network B
			& \(100-80=20\)
			& \(50-50=0\)
			& \(50-50=0\)
			& \GTenToneHL{maxregret}{20} \\
			Network C
			& \(100-60=40\)
			& \(50-30=20\)
			& \(50-35=15\)
			& \GTenToneHL{maxregret}{40} \\
			Network D
			& \(100-70=30\)
			& \(50-40=10\)
			& \(50-45=5\)
			& \GTenToneHL{maxregret}{30} \\
			\bottomrule
		\end{tabular}
	\end{center}
	
	The \GTenToneHL{maxregret}{Minimax} \GTenToneHL{maxregret}{Regret} choice is Network A because it has the smallest maximum
	\GTenToneHL{regret}{regret}, \(10\).
	
	\subsubsection*{Decisions summary}
	
	\begin{center}
		\begin{tabular}{l c}
			\toprule
			Criterion & Final decision \\
			\midrule
			\GTenToneHL{best}{Maximax}        & Network A \\
			\GTenToneHL{worst}{Maximin}        & Network B \\
			\GTenToneHL{maxregret}{Minimax} \GTenToneHL{maxregret}{Regret} & Network A \\
			\bottomrule
		\end{tabular}
	\end{center}
	
	\secInicio[sec:service-network-3x4][sec:problems-3x4]
	
	
	% ----------------------------------------------------------------
	\subsection{Inventory Strategy}
	\label{sec:inventory-strategy-3x4}
	
	A retailer must choose among Inventory Strategy A, Strategy B, Strategy C,
	and Strategy D. Three \GTenToneHL{state}{states} \GTenToneHL{state}{of} \GTenToneHL{state}{nature} represent different demand conditions.
	The final \GTenToneHL{payoff}{payoff} table, in thousands of dollars, is given below. Use \GTenToneHL{best}{Maximax},
	\GTenToneHL{worst}{Maximin}, and \GTenToneHL{maxregret}{Minimax} \GTenToneHL{maxregret}{Regret} to select an inventory strategy.
	
	\begin{center}
		\textit{Payoff table} \\
		\begin{tabular}{l c c c c}
			\toprule
			State of nature & Strategy A & Strategy B & Strategy C & Strategy D \\
			\midrule
			State \(\GTenToneHL{state}{S_1}\) & \GTenToneHL{rate}{100} & \GTenToneHL{rate}{70} & \GTenToneHL{rate}{50} & \GTenToneHL{rate}{60} \\
			State \(\GTenToneHL{state}{S_2}\) & \GTenToneHL{rate}{20}  & \GTenToneHL{rate}{60} & \GTenToneHL{rate}{40} & \GTenToneHL{rate}{45} \\
			State \(\GTenToneHL{state}{S_3}\) & \GTenToneHL{rate}{0}   & \GTenToneHL{rate}{50} & \GTenToneHL{rate}{30} & \GTenToneHL{rate}{35} \\
			\bottomrule
		\end{tabular}
	\end{center}
	
	\textbf{Questions/Tasks.}
	\begin{enumerate}
		\CriterionTask
		{subsec:inventory-strategy-3x4-c4}{4}
		{Applies the Maximin and Maximax criteria}

		\CriterionTask
		{subsec:inventory-strategy-3x4-c5}{5}
		{Constructs opportunity loss tables and applies the Minimax regret criterion}
	\end{enumerate}
	
	\subsubsection*{C4 - Applies the Maximin and Maximax criteria}
	\phantomsection
	\label{subsec:inventory-strategy-3x4-c4}
	
	\[
	\begin{aligned}
		\max(\text{Strategy A}) &= \max\{100,20,0\}=100,\\
		\max(\text{Strategy B}) &= \max\{70,60,50\}=70,\\
		\max(\text{Strategy C}) &= \max\{50,40,30\}=50,\\
		\max(\text{Strategy D}) &= \max\{60,45,35\}=60.
	\end{aligned}
	\]
	\[
	\begin{aligned}
		&\max\{\max(\text{Strategy A}),\max(\text{Strategy B}),
		\max(\text{Strategy C}),\max(\text{Strategy D})\}\\
		&\qquad=\max\{100,70,50,60\}
		=100.
	\end{aligned}
	\]
	The \GTenToneHL{best}{Maximax} choice is Strategy A because it has the highest possible \GTenToneHL{payoff}{payoff}.
	
	\[
	\begin{aligned}
		\min(\text{Strategy A}) &= \min\{100,20,0\}=0,\\
		\min(\text{Strategy B}) &= \min\{70,60,50\}=50,\\
		\min(\text{Strategy C}) &= \min\{50,40,30\}=30,\\
		\min(\text{Strategy D}) &= \min\{60,45,35\}=35.
	\end{aligned}
	\]
	\[
	\begin{aligned}
		&\max\{\min(\text{Strategy A}),\min(\text{Strategy B}),
		\min(\text{Strategy C}),\min(\text{Strategy D})\}\\
		&\qquad=\max\{0,50,30,35\}
		=50.
	\end{aligned}
	\]
	The \GTenToneHL{worst}{Maximin} choice is Strategy B because it has the largest worst-case \GTenToneHL{payoff}{payoff}.
	
	\subsubsection*{C5 - Constructs opportunity loss tables and applies the Minimax regret criterion}
	\phantomsection
	\label{subsec:inventory-strategy-3x4-c5}
	
	Best \GTenToneHL{payoff}{payoff} in each state:
	\[
	\begin{aligned}
		\text{State }\GTenToneHL{state}{S_1}:\;&\max\{100,70,50,60\}=100,\\
		\text{State }\GTenToneHL{state}{S_2}:\;&\max\{20,60,40,45\}=60,\\
		\text{State }\GTenToneHL{state}{S_3}:\;&\max\{0,50,30,35\}=50.
	\end{aligned}
	\]
	
	Regret table (best \GTenToneHL{payoff}{payoff} \( - \) \GTenToneHL{payoff}{payoff}):
	\begin{center}
		\begin{tabular}{l c c c c}
			\toprule
			Alternative & \(\GTenToneHL{state}{S_1}\) & \(\GTenToneHL{state}{S_2}\) & \(\GTenToneHL{state}{S_3}\) & \GTenToneHL{maxregret}{Maximum} \GTenToneHL{maxregret}{regret} \\
			\midrule
			Strategy A
			& \(100-100=0\)
			& \(60-20=40\)
			& \(50-0=50\)
			& \GTenToneHL{maxregret}{50} \\
			Strategy B
			& \(100-70=30\)
			& \(60-60=0\)
			& \(50-50=0\)
			& \GTenToneHL{maxregret}{30} \\
			Strategy C
			& \(100-50=50\)
			& \(60-40=20\)
			& \(50-30=20\)
			& \GTenToneHL{maxregret}{50} \\
			Strategy D
			& \(100-60=40\)
			& \(60-45=15\)
			& \(50-35=15\)
			& \GTenToneHL{maxregret}{40} \\
			\bottomrule
		\end{tabular}
	\end{center}
	
	The \GTenToneHL{maxregret}{Minimax} \GTenToneHL{maxregret}{Regret} choice is Strategy B because it has the smallest maximum
	\GTenToneHL{regret}{regret}, \(30\).
	
	\subsubsection*{Decisions summary}
	
	\begin{center}
		\begin{tabular}{l c}
			\toprule
			Criterion & Final decision \\
			\midrule
			\GTenToneHL{best}{Maximax}        & Strategy A \\
			\GTenToneHL{worst}{Maximin}        & Strategy B \\
			\GTenToneHL{maxregret}{Minimax} \GTenToneHL{maxregret}{Regret} & Strategy B \\
			\bottomrule
		\end{tabular}
	\end{center}
	
	\secInicio[sec:inventory-strategy-3x4][sec:problems-3x4]
	
	
	% ----------------------------------------------------------------
	\subsection{Event Venue}
	\label{sec:event-venue-3x4}
	
	An event company must choose among Venue A, Venue B, Venue C, and Venue D.
	Three \GTenToneHL{state}{states} \GTenToneHL{state}{of} \GTenToneHL{state}{nature} represent different attendance levels. The final \GTenToneHL{payoff}{payoff}
	table, in thousands of dollars, is given below. Use \GTenToneHL{best}{Maximax}, \GTenToneHL{worst}{Maximin}, and
	\GTenToneHL{maxregret}{Minimax} \GTenToneHL{maxregret}{Regret} to select an event venue.
	
	\begin{center}
		\textit{Payoff table} \\
		\begin{tabular}{l c c c c}
			\toprule
			State of nature & Venue A & Venue B & Venue C & Venue D \\
			\midrule
			State \(\GTenToneHL{state}{S_1}\) & \GTenToneHL{payoff}{100} & \GTenToneHL{payoff}{70} & \GTenToneHL{payoff}{80} & \GTenToneHL{payoff}{60} \\
			State \(\GTenToneHL{state}{S_2}\) & \GTenToneHL{payoff}{0}   & \GTenToneHL{payoff}{60} & \GTenToneHL{payoff}{50} & \GTenToneHL{payoff}{45} \\
			State \(\GTenToneHL{state}{S_3}\) & \GTenToneHL{payoff}{20}  & \GTenToneHL{payoff}{55} & \GTenToneHL{payoff}{50} & \GTenToneHL{payoff}{40} \\
			\bottomrule
		\end{tabular}
	\end{center}
	
	\textbf{Questions/Tasks.}
	\begin{enumerate}
		\CriterionTask
		{subsec:event-venue-3x4-c4}{4}
		{Applies the Maximin and Maximax criteria}

		\CriterionTask
		{subsec:event-venue-3x4-c5}{5}
		{Constructs opportunity loss tables and applies the Minimax regret criterion}
	\end{enumerate}
	
	\subsubsection*{C4 - Applies the Maximin and Maximax criteria}
	\phantomsection
	\label{subsec:event-venue-3x4-c4}
	
	\[
	\begin{aligned}
		\max(\text{Venue A}) &= \max\{100,0,20\}=100,\\
		\max(\text{Venue B}) &= \max\{70,60,55\}=70,\\
		\max(\text{Venue C}) &= \max\{80,50,50\}=80,\\
		\max(\text{Venue D}) &= \max\{60,45,40\}=60.
	\end{aligned}
	\]
	\[
	\begin{aligned}
		&\max\{\max(\text{Venue A}),\max(\text{Venue B}),
		\max(\text{Venue C}),\max(\text{Venue D})\}\\
		&\qquad=\max\{100,70,80,60\}
		=100.
	\end{aligned}
	\]
	The \GTenToneHL{best}{Maximax} choice is Venue A because it has the highest possible \GTenToneHL{payoff}{payoff}.
	
	\[
	\begin{aligned}
		\min(\text{Venue A}) &= \min\{100,0,20\}=0,\\
		\min(\text{Venue B}) &= \min\{70,60,55\}=55,\\
		\min(\text{Venue C}) &= \min\{80,50,50\}=50,\\
		\min(\text{Venue D}) &= \min\{60,45,40\}=40.
	\end{aligned}
	\]
	\[
	\begin{aligned}
		&\max\{\min(\text{Venue A}),\min(\text{Venue B}),
		\min(\text{Venue C}),\min(\text{Venue D})\}\\
		&\qquad=\max\{0,55,50,40\}
		=55.
	\end{aligned}
	\]
	The \GTenToneHL{worst}{Maximin} choice is Venue B because it has the largest worst-case \GTenToneHL{payoff}{payoff}.
	
	\subsubsection*{C5 - Constructs opportunity loss tables and applies the Minimax regret criterion}
	\phantomsection
	\label{subsec:event-venue-3x4-c5}
	
	Best \GTenToneHL{payoff}{payoff} in each state:
	\[
	\begin{aligned}
		\text{State }\GTenToneHL{state}{S_1}:\;&\max\{100,70,80,60\}=100,\\
		\text{State }\GTenToneHL{state}{S_2}:\;&\max\{0,60,50,45\}=60,\\
		\text{State }\GTenToneHL{state}{S_3}:\;&\max\{20,55,50,40\}=55.
	\end{aligned}
	\]
	
	Regret table (best \GTenToneHL{payoff}{payoff} \( - \) \GTenToneHL{payoff}{payoff}):
	\begin{center}
		\begin{tabular}{l c c c c}
			\toprule
			Alternative & \(\GTenToneHL{state}{S_1}\) & \(\GTenToneHL{state}{S_2}\) & \(\GTenToneHL{state}{S_3}\) & \GTenToneHL{maxregret}{Maximum} \GTenToneHL{maxregret}{regret} \\
			\midrule
			Venue A
			& \(100-100=0\)
			& \(60-0=60\)
			& \(55-20=35\)
			& \GTenToneHL{maxregret}{60} \\
			Venue B
			& \(100-70=30\)
			& \(60-60=0\)
			& \(55-55=0\)
			& \GTenToneHL{maxregret}{30} \\
			Venue C
			& \(100-80=20\)
			& \(60-50=10\)
			& \(55-50=5\)
			& \GTenToneHL{maxregret}{20} \\
			Venue D
			& \(100-60=40\)
			& \(60-45=15\)
			& \(55-40=15\)
			& \GTenToneHL{maxregret}{40} \\
			\bottomrule
		\end{tabular}
	\end{center}
	
	The \GTenToneHL{maxregret}{Minimax} \GTenToneHL{maxregret}{Regret} choice is Venue C because it has the smallest maximum
	\GTenToneHL{regret}{regret}, \(20\).
	
	\subsubsection*{Decisions summary}
	
	\begin{center}
		\begin{tabular}{l c}
			\toprule
			Criterion & Final decision \\
			\midrule
			\GTenToneHL{best}{Maximax}        & Venue A \\
			\GTenToneHL{worst}{Maximin}        & Venue B \\
			\GTenToneHL{maxregret}{Minimax} \GTenToneHL{maxregret}{Regret} & Venue C \\
			\bottomrule
		\end{tabular}
	\end{center}
	
	\secInicio[sec:event-venue-3x4][sec:problems-3x4]
	
	
	% ----------------------------------------------------------------
	\subsection{Supply Contract}
	\label{sec:supply-contract-3x4}
	
	A company must choose among Supply Contract A, Contract B, Contract C, and
	Contract D. Three \GTenToneHL{state}{states} \GTenToneHL{state}{of} \GTenToneHL{state}{nature} represent different market and supply
	conditions. The final \GTenToneHL{payoff}{payoff} table, in thousands of dollars, is given below.
	Use \GTenToneHL{best}{Maximax}, \GTenToneHL{worst}{Maximin}, and \GTenToneHL{maxregret}{Minimax} \GTenToneHL{maxregret}{Regret} to select a supply contract.
	
	\begin{center}
		\textit{Payoff table} \\
		\begin{tabular}{l c c c c}
			\toprule
			State of nature & Contract A & Contract B & Contract C & Contract D \\
			\midrule
			State \(\GTenToneHL{state}{S_1}\) & \GTenToneHL{payoff}{90} & \GTenToneHL{payoff}{80} & \GTenToneHL{payoff}{70} & \GTenToneHL{payoff}{75} \\
			State \(\GTenToneHL{state}{S_2}\) & \GTenToneHL{payoff}{70} & \GTenToneHL{payoff}{50} & \GTenToneHL{payoff}{45} & \GTenToneHL{payoff}{55} \\
			State \(\GTenToneHL{state}{S_3}\) & \GTenToneHL{payoff}{60} & \GTenToneHL{payoff}{40} & \GTenToneHL{payoff}{35} & \GTenToneHL{payoff}{50} \\
			\bottomrule
		\end{tabular}
	\end{center}
	
	\textbf{Questions/Tasks.}
	\begin{enumerate}
		\CriterionTask
		{subsec:supply-contract-3x4-c4}{4}
		{Applies the Maximin and Maximax criteria}

		\CriterionTask
		{subsec:supply-contract-3x4-c5}{5}
		{Constructs opportunity loss tables and applies the Minimax regret criterion}
	\end{enumerate}
	
	\subsubsection*{C4 - Applies the Maximin and Maximax criteria}
	\phantomsection
	\label{subsec:supply-contract-3x4-c4}
	
	\[
	\begin{aligned}
		\max(\text{Contract A}) &= \max\{90,70,60\}=90,\\
		\max(\text{Contract B}) &= \max\{80,50,40\}=80,\\
		\max(\text{Contract C}) &= \max\{70,45,35\}=70,\\
		\max(\text{Contract D}) &= \max\{75,55,50\}=75.
	\end{aligned}
	\]
	\[
	\begin{aligned}
		&\max\{\max(\text{Contract A}),\max(\text{Contract B}),
		\max(\text{Contract C}),\max(\text{Contract D})\}\\
		&\qquad=\max\{90,80,70,75\}
		=90.
	\end{aligned}
	\]
	The \GTenToneHL{best}{Maximax} choice is Contract A because it has the highest possible \GTenToneHL{payoff}{payoff}.
	
	\[
	\begin{aligned}
		\min(\text{Contract A}) &= \min\{90,70,60\}=60,\\
		\min(\text{Contract B}) &= \min\{80,50,40\}=40,\\
		\min(\text{Contract C}) &= \min\{70,45,35\}=35,\\
		\min(\text{Contract D}) &= \min\{75,55,50\}=50.
	\end{aligned}
	\]
	\[
	\begin{aligned}
		&\max\{\min(\text{Contract A}),\min(\text{Contract B}),
		\min(\text{Contract C}),\min(\text{Contract D})\}\\
		&\qquad=\max\{60,40,35,50\}
		=60.
	\end{aligned}
	\]
	The \GTenToneHL{worst}{Maximin} choice is Contract A because it has the largest worst-case \GTenToneHL{payoff}{payoff}.
	
	\subsubsection*{C5 - Constructs opportunity loss tables and applies the Minimax regret criterion}
	\phantomsection
	\label{subsec:supply-contract-3x4-c5}
	
	Best \GTenToneHL{payoff}{payoff} in each state:
	\[
	\begin{aligned}
		\text{State }\GTenToneHL{state}{S_1}:\;&\max\{90,80,70,75\}=90,\\
		\text{State }\GTenToneHL{state}{S_2}:\;&\max\{70,50,45,55\}=70,\\
		\text{State }\GTenToneHL{state}{S_3}:\;&\max\{60,40,35,50\}=60.
	\end{aligned}
	\]
	
	Regret table (best \GTenToneHL{payoff}{payoff} \( - \) \GTenToneHL{payoff}{payoff}):
	\begin{center}
		\begin{tabular}{l c c c c}
			\toprule
			Alternative & \(\GTenToneHL{state}{S_1}\) & \(\GTenToneHL{state}{S_2}\) & \(\GTenToneHL{state}{S_3}\) & \GTenToneHL{maxregret}{Maximum} \GTenToneHL{maxregret}{regret} \\
			\midrule
			Contract A
			& \(90-90=0\)
			& \(70-70=0\)
			& \(60-60=0\)
			& \GTenToneHL{maxregret}{0} \\
			Contract B
			& \(90-80=10\)
			& \(70-50=20\)
			& \(60-40=20\)
			& \GTenToneHL{maxregret}{20} \\
			Contract C
			& \(90-70=20\)
			& \(70-45=25\)
			& \(60-35=25\)
			& \GTenToneHL{maxregret}{25} \\
			Contract D
			& \(90-75=15\)
			& \(70-55=15\)
			& \(60-50=10\)
			& \GTenToneHL{maxregret}{15} \\
			\bottomrule
		\end{tabular}
	\end{center}
	
	The \GTenToneHL{maxregret}{Minimax} \GTenToneHL{maxregret}{Regret} choice is Contract A because it has the smallest maximum
	\GTenToneHL{regret}{regret}, \(0\).
	
	\subsubsection*{Decisions summary}
	
	\begin{center}
		\begin{tabular}{l c}
			\toprule
			Criterion & Final decision \\
			\midrule
			\GTenToneHL{best}{Maximax}        & Contract A \\
			\GTenToneHL{worst}{Maximin}        & Contract A \\
			\GTenToneHL{maxregret}{Minimax} \GTenToneHL{maxregret}{Regret} & Contract A \\
			\bottomrule
		\end{tabular}
	\end{center}
	
	\secInicio[sec:supply-contract-3x4][sec:problems-3x4]
	
	
	% ================================================================
\end{document}